% TRUE_DIFF_ARTIFACT: buildable latexdiff comparison against documented CCVGAE baseline.
% Baseline: CCVGAE/CCVGAE_snLaTeX/sn-article.tex
% Baseline-SHA256: 94e03054911e922d0ea2ea5273a9414c9d83c3134dda672fd170a7ce8ba5b00c
% Current: manuscript/scccvgben/sn-article.tex
% Current-SHA256: 3d8ba2fe1a889786f0693c26801236ae0d36c95bf1444158b67c7f204d42146d
% Git-Revision: 4249bbaab09f
% Generated-UTC: 2026-04-29T11:02:18Z
% Method: latexdiff-no-display-floats
% Prior-artifact-rejected: previous buildable review copy had zero body-level DIFadd/DIFdel evidence.
% Build-stability-note: display floats/longtables are omitted before latexdiff to avoid invalid nested caption/table markup; prose, equations, sections, references, and bibliography remain diffed.
% Redaction-note: restricted accessions from the baseline are replaced with neutral placeholders before building.
% Confidentiality-note: metadata uses repository-relative paths only; no local absolute paths are embedded.
%DIF 1-18d1
%DIF LATEXDIFF DIFFERENCE FILE


%DIF < %Version 3.1 December 2024
%DIF < % See section 11 of the User Manual for version history
%DIF < %
%DIF < %%%%%%%%%%%%%%%%%%%%%%%%%%%%%%%%%%%%%%%%%%%%%%%%%%%%%%%%%%%%%%%%%%%%%%
%DIF < %%                                                                 %%

%DIF < %% Submit your LaTeX manuscript as one .tex document.              %%
%DIF < %%                                                                 %%
%DIF < %% All additional figures and files should be attached             %%
%DIF < %% separately and not embedded in the \TeX\ document itself.       %%
%DIF < %%                                                                 %%
%DIF < %%%%%%%%%%%%%%%%%%%%%%%%%%%%%%%%%%%%%%%%%%%%%%%%%%%%%%%%%%%%%%%%%%%%%
%DIF <
%DIF < %%\documentclass[referee,sn-basic]{sn-jnl}% referee option is meant for double line spacing
%DIF <
%DIF < %%=======================================================%%
%DIF < %% to print line numbers in the margin use lineno option %%
%DIF < %%=======================================================%%
%DIF -------
\PassOptionsToPackage{hypertexnames=false}{hyperref}
\documentclass[referee, pdflatex, sn-vancouver-num]{sn-jnl}
%DIF 21d3
%DIF < %\documentclass[referee, lineno, pdflatex, sn-basic]{sn-jnl}% Basic Springer Nature Reference Style/Chemistry Reference Style
%DIF -------

%DIF 23-40d4
%DIF < %%=========================================================================================%%
%DIF < %% the documentclass is set to pdflatex as default. You can delete it if not appropriate.  %%
%DIF < %%=========================================================================================%%
%DIF <
%DIF < %%\documentclass[sn-basic]{sn-jnl}% Basic Springer Nature Reference Style/Chemistry Reference Style
%DIF <
%DIF < %%Note: the following reference styles support Namedate and Numbered referencing. By default the style follows the most common style. To switch between the options you can add or remove �Numbered� in the optional parenthesis.
%DIF < %%The option is available for: sn-basic.bst, sn-chicago.bst%
%DIF <
%DIF < %\documentclass[pdflatex,sn-nature]{sn-jnl}% Style for submissions to Nature Portfolio journals
%DIF < %%\documentclass[pdflatex,sn-basic]{sn-jnl}% Basic Springer Nature Reference Style/Chemistry Reference Style
%DIF < %\documentclass[pdflatex,sn-mathphys-num]{sn-jnl}% Math and Physical Sciences Numbered Reference Style
%DIF < %%\documentclass[pdflatex,sn-mathphys-ay]{sn-jnl}% Math and Physical Sciences Author Year Reference Style
%DIF < %%\documentclass[pdflatex,sn-aps]{sn-jnl}% American Physical Society (APS) Reference Style
%DIF < %%\documentclass[pdflatex,sn-vancouver-num]{sn-jnl}% Vancouver Numbered Reference Style
%DIF < %%\documentclass[pdflatex,sn-vancouver-ay]{sn-jnl}% Vancouver Author Year Reference Style
%DIF < %%\documentclass[pdflatex,sn-apa]{sn-jnl}% APA Reference Style
%DIF < %%\documentclass[pdflatex,sn-chicago]{sn-jnl}% Chicago-based Humanities Reference Style
%DIF -------
\usepackage{microtype}
%DIF 42-45c5-9
%DIF < \geometry{left=1.5cm, right=1.5cm, top=2cm, bottom=2cm}
%DIF < \usepackage{graphicx}%
%DIF < \usepackage{multirow}%
%DIF < \usepackage{booktabs}%
%DIF -------
\usepackage{geometry} %DIF >
\geometry{a4paper,top=2cm,bottom=2cm,left=1.7cm,right=1.7cm} %DIF >
\usepackage{graphicx} %DIF >
\usepackage{multirow} %DIF >
\usepackage{booktabs} %DIF >
%DIF -------
\usepackage{multicol}
%DIF 47-58c11-24
%DIF < \usepackage{amsmath,amssymb,amsfonts}%
%DIF < \usepackage{amsthm}%
%DIF < \usepackage{mathrsfs}%
%DIF < \usepackage[title]{appendix}%
%DIF < \usepackage{xcolor}%
%DIF < \usepackage{textcomp}%
%DIF < \usepackage{manyfoot}%
%DIF < \usepackage{booktabs}%
%DIF < \usepackage{algorithm}%
%DIF < \usepackage{algorithmicx}%
%DIF < \usepackage{algpseudocode}%
%DIF < \usepackage{listings}%
%DIF -------
\usepackage{amsmath,amssymb,amsfonts} %DIF >
\usepackage{amsthm} %DIF >
\usepackage{mathrsfs} %DIF >
\usepackage[title]{appendix} %DIF >
\usepackage{xcolor} %DIF >
\usepackage{textcomp} %DIF >
\usepackage{manyfoot} %DIF >
\usepackage{algorithm} %DIF >
\usepackage{algorithmicx} %DIF >
\usepackage{algpseudocode} %DIF >
\usepackage{listings} %DIF >
\usepackage{longtable} %DIF >
\usepackage{array} %DIF >
\usepackage{siunitx} %DIF >
%DIF -------
\usepackage{hyperref}
\usepackage{rotating}
%DIF 61-67d27
%DIF < %% as per the requirement new theorem styles can be included as shown below
%DIF < \theoremstyle{thmstyleone}%
%DIF < \newtheorem{theorem}{Theorem}%  meant for continuous numbers
%DIF < %%\newtheorem{theorem}{Theorem}[section]% meant for sectionwise numbers
%DIF < %% optional argument [theorem] produces theorem numbering sequence instead of independent numbers for Proposition
%DIF < \newtheorem{proposition}[theorem]{Proposition}%
%DIF < %%\newtheorem{proposition}{Proposition}% to get separate numbers for theorem and proposition etc.
%DIF -------

%DIF < \theoremstyle{thmstyletwo}%
%DIF -------
\graphicspath{{../}{../../../figures/}{figures/}} %DIF >
%DIF < \newtheorem{example}{Example}%
%DIF < \newtheorem{remark}{Remark}%
%DIF < \theoremstyle{thmstylethree}%
%DIF < \newtheorem{definition}{Definition}%
%DIF < \raggedbottom
%DIF <
%DIF < \setlength{\topmargin}{0in}
%DIF < \setlength{\textheight}{\paperheight}
%DIF < \addtolength{\textheight}{-2in}
%DIF <
%DIF < \setlength{\floatsep}{10pt}         % Space between adjacent floats
%DIF < \setlength{\textfloatsep}{15pt}     % Space between text and float
%DIF < \setlength{\intextsep}{10pt}        % Space between in-text floats and text
%DIF < %%\unnumbered% uncomment this for unnumbered level heads
%DIF PREAMBLE EXTENSION ADDED BY LATEXDIFF
%DIF UNDERLINE PREAMBLE %DIF PREAMBLE
\RequirePackage[normalem]{ulem} %DIF PREAMBLE
\RequirePackage{color}\definecolor{RED}{rgb}{1,0,0}\definecolor{BLUE}{rgb}{0,0,1} %DIF PREAMBLE
\providecommand{\DIFaddtex}[1]{{\protect\color{blue}\uwave{#1}}} %DIF PREAMBLE
\providecommand{\DIFdeltex}[1]{{\protect\color{red}\sout{#1}}}                      %DIF PREAMBLE
%DIF SAFE PREAMBLE %DIF PREAMBLE
\providecommand{\DIFaddbegin}{} %DIF PREAMBLE
\providecommand{\DIFaddend}{} %DIF PREAMBLE
\providecommand{\DIFdelbegin}{} %DIF PREAMBLE
\providecommand{\DIFdelend}{} %DIF PREAMBLE
\providecommand{\DIFmodbegin}{} %DIF PREAMBLE
\providecommand{\DIFmodend}{} %DIF PREAMBLE
%DIF FLOATSAFE PREAMBLE %DIF PREAMBLE
\providecommand{\DIFaddFL}[1]{\DIFadd{#1}} %DIF PREAMBLE
\providecommand{\DIFdelFL}[1]{\DIFdel{#1}} %DIF PREAMBLE
\providecommand{\DIFaddbeginFL}{} %DIF PREAMBLE
\providecommand{\DIFaddendFL}{} %DIF PREAMBLE
\providecommand{\DIFdelbeginFL}{} %DIF PREAMBLE
\providecommand{\DIFdelendFL}{} %DIF PREAMBLE
%DIF HYPERREF PREAMBLE %DIF PREAMBLE
\providecommand{\DIFadd}[1]{\texorpdfstring{\DIFaddtex{#1}}{#1}} %DIF PREAMBLE
\providecommand{\DIFdel}[1]{\texorpdfstring{\DIFdeltex{#1}}{}} %DIF PREAMBLE
\newcommand{\DIFscaledelfig}{0.5}
%DIF HIGHLIGHTGRAPHICS PREAMBLE %DIF PREAMBLE
\RequirePackage{settobox} %DIF PREAMBLE
\RequirePackage{letltxmacro} %DIF PREAMBLE
\newsavebox{\DIFdelgraphicsbox} %DIF PREAMBLE
\newlength{\DIFdelgraphicswidth} %DIF PREAMBLE
\newlength{\DIFdelgraphicsheight} %DIF PREAMBLE
% store original definition of \includegraphics %DIF PREAMBLE
\LetLtxMacro{\DIFOincludegraphics}{\includegraphics} %DIF PREAMBLE
\newcommand{\DIFaddincludegraphics}[2][]{{\color{blue}\fbox{\DIFOincludegraphics[#1]{#2}}}} %DIF PREAMBLE
\newcommand{\DIFdelincludegraphics}[2][]{% %DIF PREAMBLE
\sbox{\DIFdelgraphicsbox}{\DIFOincludegraphics[#1]{#2}}% %DIF PREAMBLE
\settoboxwidth{\DIFdelgraphicswidth}{\DIFdelgraphicsbox} %DIF PREAMBLE
\settoboxtotalheight{\DIFdelgraphicsheight}{\DIFdelgraphicsbox} %DIF PREAMBLE
\scalebox{\DIFscaledelfig}{% %DIF PREAMBLE
\parbox[b]{\DIFdelgraphicswidth}{\usebox{\DIFdelgraphicsbox}\\[-\baselineskip] \rule{\DIFdelgraphicswidth}{0em}}\llap{\resizebox{\DIFdelgraphicswidth}{\DIFdelgraphicsheight}{% %DIF PREAMBLE
\setlength{\unitlength}{\DIFdelgraphicswidth}% %DIF PREAMBLE
\begin{picture}(1,1)% %DIF PREAMBLE
\thicklines\linethickness{2pt} %DIF PREAMBLE
{\color[rgb]{1,0,0}\put(0,0){\framebox(1,1){}}}% %DIF PREAMBLE
{\color[rgb]{1,0,0}\put(0,0){\line( 1,1){1}}}% %DIF PREAMBLE
{\color[rgb]{1,0,0}\put(0,1){\line(1,-1){1}}}% %DIF PREAMBLE
\end{picture}% %DIF PREAMBLE
}\hspace*{3pt}}} %DIF PREAMBLE
} %DIF PREAMBLE
\LetLtxMacro{\DIFOaddbegin}{\DIFaddbegin} %DIF PREAMBLE
\LetLtxMacro{\DIFOaddend}{\DIFaddend} %DIF PREAMBLE
\LetLtxMacro{\DIFOdelbegin}{\DIFdelbegin} %DIF PREAMBLE
\LetLtxMacro{\DIFOdelend}{\DIFdelend} %DIF PREAMBLE
\DeclareRobustCommand{\DIFaddbegin}{\DIFOaddbegin \let\includegraphics\DIFaddincludegraphics} %DIF PREAMBLE
\DeclareRobustCommand{\DIFaddend}{\DIFOaddend \let\includegraphics\DIFOincludegraphics} %DIF PREAMBLE
\DeclareRobustCommand{\DIFdelbegin}{\DIFOdelbegin \let\includegraphics\DIFdelincludegraphics} %DIF PREAMBLE
\DeclareRobustCommand{\DIFdelend}{\DIFOaddend \let\includegraphics\DIFOincludegraphics} %DIF PREAMBLE
\LetLtxMacro{\DIFOaddbeginFL}{\DIFaddbeginFL} %DIF PREAMBLE
\LetLtxMacro{\DIFOaddendFL}{\DIFaddendFL} %DIF PREAMBLE
\LetLtxMacro{\DIFOdelbeginFL}{\DIFdelbeginFL} %DIF PREAMBLE
\LetLtxMacro{\DIFOdelendFL}{\DIFdelendFL} %DIF PREAMBLE
\DeclareRobustCommand{\DIFaddbeginFL}{\DIFOaddbeginFL \let\includegraphics\DIFaddincludegraphics} %DIF PREAMBLE
\DeclareRobustCommand{\DIFaddendFL}{\DIFOaddendFL \let\includegraphics\DIFOincludegraphics} %DIF PREAMBLE
\DeclareRobustCommand{\DIFdelbeginFL}{\DIFOdelbeginFL \let\includegraphics\DIFdelincludegraphics} %DIF PREAMBLE
\DeclareRobustCommand{\DIFdelendFL}{\DIFOaddendFL \let\includegraphics\DIFOincludegraphics} %DIF PREAMBLE
%DIF COLORLISTINGS PREAMBLE %DIF PREAMBLE
\RequirePackage{listings} %DIF PREAMBLE
\RequirePackage{color} %DIF PREAMBLE
\lstdefinelanguage{DIFcode}{ %DIF PREAMBLE
%DIF DIFCODE_UNDERLINE %DIF PREAMBLE
  moredelim=[il][\color{red}\sout]{\%DIF\ <\ }, %DIF PREAMBLE
  moredelim=[il][\color{blue}\uwave]{\%DIF\ >\ } %DIF PREAMBLE
} %DIF PREAMBLE
\lstdefinestyle{DIFverbatimstyle}{ %DIF PREAMBLE
	language=DIFcode, %DIF PREAMBLE
	basicstyle=\ttfamily, %DIF PREAMBLE
	columns=fullflexible, %DIF PREAMBLE
	keepspaces=true %DIF PREAMBLE
} %DIF PREAMBLE
\lstnewenvironment{DIFverbatim}{\lstset{style=DIFverbatimstyle}}{} %DIF PREAMBLE
\lstnewenvironment{DIFverbatim*}{\lstset{style=DIFverbatimstyle,showspaces=true}}{} %DIF PREAMBLE
%DIF END PREAMBLE EXTENSION ADDED BY LATEXDIFF

\begin{document}

\DIFdelbegin %DIFDELCMD < \title[CCVGAE]{%%%
\DIFdel{CCVGAE: A Centroid-Coupled Variational Graph Attention Autoencoder }\DIFdelend \DIFaddbegin \title[scCCVGBen for benchmarking of single-cell representation-learning]{\DIFadd{scCCVGBen }\DIFaddend for \DIFdelbegin \DIFdel{Stable }\DIFdelend \DIFaddbegin \DIFadd{benchmarking of single-cell representation-learning anchored on a centroid-coupled variational graph attention autoencoder across scRNA-seq }\DIFaddend and \DIFdelbegin \DIFdel{Interpretable Single-Cell Representation Learning}\DIFdelend \DIFaddbegin \DIFadd{scATAC-seq}\DIFaddend }

\author*[1]{\fnm{Zeyu} \sur{Fu}}\email{fuzeyu99@126.com}\equalcont{These authors contributed equally to this work.}
\author[1]{\fnm{Yiyao} \sur{Liu}}\email{2415084056@qq.com}\equalcont{These authors contributed equally to this work.}
\author[1]{\fnm{Junping} \sur{Wang}}\email{wangjunping@tmmu.edu.cn}
\author[1]{\fnm{Song} \sur{Wang}}\email{swang1981@tmmu.edu.cn}

\affil[1]{\orgdiv{State Key Laboratory of Trauma and Chemical Poisoning, Institute of Combined Injury, Chongqing Engineering Research Center for Nanomedicine, College of Preventive Medicine}, \orgname{Army Medical University}, \orgaddress{\city{Chongqing}, \postcode{400038}, \country{China}}}

\DIFdelbegin %DIFDELCMD < \abstract{
%DIFDELCMD < Single-cell omics technologies are widely used to dissect cellular heterogeneity, yet computational analysis across modalities remains difficult because of high dimensionality, technical noise, and the instability of stochastic deep generative models. We present CCVGAE, a Centroid-based Coupled Variational Graph Attention Autoencoder built around three design choices: (i) Centroid Inference, which uses the deterministic posterior mean as the cell embedding to avoid sampling variance; (ii) a Coupling mechanism that regularizes local latent geometry via a dual-reconstruction bottleneck ($d_c < d_z$); and (iii) a graph-attention encoder that uses cell--cell similarity to preserve neighborhood structure. Across 170 benchmark datasets (55 scRNA-seq, 115 scATAC-seq), CCVGAE improves clustering accuracy (ARI $+0.104$ and NMI $+0.104$ over CouVAE on scRNA-seq; ASW $+0.329$ versus scVI, $p < 0.001$), embedding quality (overall UMAP score $+0.219$ versus scVI on scRNA-seq), and intrinsic manifold geometry (core quality $+0.273$ versus PeakVI on scATAC-seq) relative to classical dimensionality reduction methods and established deep generative models. Individual latent components map to coherent Gene Ontology Biological Processes. Case studies in hematopoietic systems under physiological stress---sleep deprivation, radiation injury, and thrombopoietin-driven megakaryopoiesis---indicate that CCVGAE captures components associated with lineage regulation, inflammatory responses, cell-cycle dynamics, and hemostasis-related programs. By producing stable and interpretable embeddings, CCVGAE supports downstream single-cell analyses including clustering, trajectory inference, and visualization.
%DIFDELCMD < }
%DIFDELCMD < %%%
\DIFdelend \DIFaddbegin \abstract{
Single-cell omics technologies are widely used to dissect cellular heterogeneity, yet computational analysis across modalities remains difficult because of high dimensionality, technical noise, and the instability of stochastic deep generative models. We present scCCVGBen, a single-cell Centroid-Coupled Variational Graph autoencoder Benchmark whose evaluated reference configuration uses three design choices: (i) Centroid Inference, which uses the deterministic posterior mean as the cell embedding to avoid sampling variance and produce a stable representation; (ii) a Coupling mechanism that regularizes local latent geometry through an interpretable dual-reconstruction bottleneck ($d_c < d_z$); and (iii) a graph-attention encoder that uses cell--cell similarity to preserve neighbourhood structure. scCCVGBen is assessed within a decoupled benchmark in which the algorithmic core, the encoder backbone, the graph-construction strategy, the dataset cohort and the evaluation suite are evaluated as separate axes. The benchmark cohort, sourced from the Gene Expression Omnibus and the European Nucleotide Archive, balances scRNA-seq and scATAC-seq equally and spans hematopoiesis, neuronal differentiation, immune populations, organ atlases, tumour microenvironments and developmental time courses. Across the cohort, scCCVGBen improves the Average Silhouette Width by $+0.288$ and intrinsic-overall geometry by $+0.233$ over a stochastic VAE baseline on paired scRNA-seq datasets; the corresponding gains over scVI are $+0.341$ and $+0.331$. On paired scATAC-seq datasets, scCCVGBen improves intrinsic-overall by $+0.299$ over LSI and $+0.346$ over PeakVI. Robustness analyses across fourteen graph encoder backbones and five graph-construction strategies show where the centroid-coupled core remains competitive and where alternative architectural choices outperform it on individual metric families. Three paired biological case studies---sleep deprivation alongside a gastric cancer microenvironment, thrombopoietin-induced megakaryopoiesis from cord blood alongside an aged hematopoietic stem cell time course, and radiation-injury hematopoiesis alongside the COVID-19 bronchoalveolar lavage immune landscape---show that individual latent coordinates capture coherent hematopoietic-immune, epithelial-stromal, megakaryocytic and antiviral-macrophage programs whose top-correlated genes map to coherent Gene Ontology Biological Process terms. scCCVGBen produces stable, interpretable embeddings and reports the architectural axes around the algorithmic core, so clustering, trajectory inference and visualisation in single-cell omics can be compared across methods on the same protocol.
}
\DIFaddend

\DIFdelbegin %DIFDELCMD < \keywords{single-cell, representation learning, variational autoencoder, graph neural network, manifold learning, representation stability}
%DIFDELCMD < %%%
\DIFdelend \DIFaddbegin \keywords{single-cell, representation learning, variational autoencoder, graph neural network, manifold learning, representation stability, benchmarking}
\DIFaddend

\maketitle

\section*{Key Points}

\begin{itemize}
\item \textbf{Centroid Inference} \DIFdelbegin \DIFdel{bypasses sampling noise at inference time by using deterministic posterior means as cell embeddings, providing }\DIFdelend \DIFaddbegin \DIFadd{uses the deterministic posterior mean as the cell embedding, removing the sampling noise that destabilises stochastic VAE outputs and giving }\DIFaddend consistent improvements in clustering \DIFdelbegin \DIFdel{accuracy and geometric quality }\DIFdelend \DIFaddbegin \DIFadd{compactness and intrinsic latent geometry }\DIFaddend across both scRNA-seq and scATAC-seq \DIFdelbegin \DIFdel{modalities }\DIFdelend when combined with complementary architectural components.
\DIFdelbegin %DIFDELCMD <

%DIFDELCMD < %%%
\DIFdelend \item \DIFdelbegin %DIFDELCMD < \textbf{Complementary architecture:} %%%
\DIFdel{Coupling regularization refines latent geometry while graph attentioncaptures cell--cell structure; their combination yields complementary improvements over single-component baselines.
}%DIFDELCMD <

%DIFDELCMD < %%%
\DIFdelend \DIFaddbegin \textbf{Coupling Regularisation} \DIFadd{adds an interpretable bottleneck pathway with a dual-reconstruction objective ($d_c < d_z$) during training; together with graph attention, the three components produce non-redundant gains that none of them recover in isolation.
}\DIFaddend \item \DIFdelbegin %DIFDELCMD < \textbf{Benchmarking across 170 datasets:} %%%
\DIFdel{Evaluation against classical methods and deep generative models shows favorable performance with consistent gains in intrinsic manifold properties.
}%DIFDELCMD <

%DIFDELCMD < %%%
\DIFdelend \DIFaddbegin \textbf{A decoupled benchmark} \DIFadd{separates the algorithmic core from the encoder backbone, the cell--cell graph construction, the dataset cohort and the evaluation suite, and benchmarks fourteen graph encoder backbones and five graph-construction strategies under a single evaluation protocol on a balanced cohort of scRNA-seq and scATAC-seq datasets.
}\DIFaddend \item \DIFdelbegin %DIFDELCMD < \textbf{Biological interpretability:} %%%
\DIFdel{Latent dimensions correspond to coherent biological programs including lineage regulation, immune activation, }\DIFdelend \DIFaddbegin \textbf{Quantitative gains under paired Wilcoxon Holm-corrected analysis.} \DIFadd{Against a stochastic VAE baseline, scCCVGBen improves ASW by $+0.288$, DAV by $-0.800$, CAL by $+4270.9$ }\DIFaddend and \DIFdelbegin \DIFdel{cell-cycle control, validated through GO enrichment in hematopoietic perturbation studies.
}%DIFDELCMD <

%DIFDELCMD < %%%
\DIFdelend \DIFaddbegin \DIFadd{intrinsic-overall by $+0.233$ on the paired scRNA-seq cohort (Holm-corrected $p<0.001$ throughout); against PeakVI on the paired scATAC-seq cohort it improves ASW by $+0.139$ and intrinsic-overall by $+0.346$.
}\DIFaddend \item \DIFdelbegin %DIFDELCMD < \textbf{Practical utility:} %%%
\DIFdel{Stable embeddings enhance downstream analyses including clustering, trajectory inference, and visualization, supporting reproducible biological discovery }\DIFdelend \DIFaddbegin \textbf{Biological interpretability.} \DIFadd{Latent components recovered by the model map to coherent hematopoietic, epithelial-stromal, megakaryocytic and antiviral-macrophage programs whose top-correlated genes are consistent with Gene Ontology Biological Process terms in three paired hematopoietic case studies, supporting the use of scCCVGBen for downstream biological discovery alongside the quantitative benchmark}\DIFaddend .
\end{itemize}

\section{Introduction}

Single-cell technologies have advanced the study of cellular heterogeneity, enabling high-resolution profiling of molecular states across the transcriptome \DIFdelbegin \DIFdel{, epigenome, and proteome }\DIFdelend \DIFaddbegin \DIFadd{and the epigenome }\DIFaddend in complex tissues\DIFdelbegin \DIFdel{\mbox{%DIFAUXCMD
\cite{haque2017practical,kelsey2017single}}\hskip0pt%DIFAUXCMD
. Single-cell multi-omics---simultaneous measurement of multiple molecular layers within the same cell---provides additional views of cellular identity, regulatory mechanisms, and functional states \mbox{%DIFAUXCMD
\cite{ma2020single}}\hskip0pt%DIFAUXCMD
, and is useful for studying developmental trajectories, disease progression, and cellular responses to perturbation. The present work applies CCVGAE separately to individual single-cell modalities (scRNA-seq and scATAC-seq) rather than performing joint multi-omic integration from paired measurements. Downstream tasks such as cell--cell communication (CCC) inference---including CellDialog \mbox{%DIFAUXCMD
\cite{peng2023celldialog}}\hskip0pt%DIFAUXCMD
, heterogeneous ensemble deep learning approaches \mbox{%DIFAUXCMD
\cite{peng2025ccc}}\hskip0pt%DIFAUXCMD
, and CELLetter \mbox{%DIFAUXCMD
\cite{wu2025celetter}}\hskip0pt%DIFAUXCMD
---depend on the quality and stability of input cell embeddings \mbox{%DIFAUXCMD
\cite{cheng2023review,wang2023ccc,jin2023ccc}}\hskip0pt%DIFAUXCMD
. }%DIFDELCMD <

%DIFDELCMD < %%%
\DIFdelend \DIFaddbegin \DIFadd{~\mbox{%DIFAUXCMD
\cite{ma2020single}}\hskip0pt%DIFAUXCMD
. }\DIFaddend Reliable computational analysis of single-cell data\DIFdelbegin \DIFdel{remains difficult. Individual modalities---transcriptomic or epigenomic---exhibit }\DIFdelend \DIFaddbegin \DIFadd{, however, remains difficult: }\DIFaddend high dimensionality, sparsity, technical noise \DIFdelbegin \DIFdel{, }\DIFdelend and batch effects \DIFdelbegin \DIFdel{that }\DIFdelend can obscure biological variation\DIFdelbegin \DIFdel{\mbox{%DIFAUXCMD
\cite{luecken2022benchmarking,stuart2019integrative}}\hskip0pt%DIFAUXCMD
. Because }\DIFdelend \DIFaddbegin \DIFadd{~\mbox{%DIFAUXCMD
\cite{luecken2022benchmarking,stuart2019integrative}}\hskip0pt%DIFAUXCMD
, and downstream tasks such as }\DIFaddend clustering, trajectory inference \DIFdelbegin \DIFdel{, cell-type annotation, and communication analysis all depend on the learned representations , methods need stable, geometry-preserving embeddings that separate biological signal from technical artifacts.
}%DIFDELCMD <

%DIFDELCMD < %%%
\DIFdel{Evaluating multi-component systems requires principled benchmarking. In single-cell analysis, benchmarking studies have established that no single metric captures all relevant aspects of embedding quality \mbox{%DIFAUXCMD
\cite{luecken2022benchmarking}}\hskip0pt%DIFAUXCMD
, and that methods must be assessed across complementary criteria---including clustering accuracy, trajectory fidelity, and biological concordance---to draw reliable conclusions \mbox{%DIFAUXCMD
\cite{ma2020single,stuart2019integrative}}\hskip0pt%DIFAUXCMD
. Rigorous multi-criteria evaluation is a cross-disciplinary concern: analogous ensemble and hybrid machine-learning benchmarking practices in adjacent engineering fields \mbox{%DIFAUXCMD
\cite{naeem2025ensemble,naeem2025hybrid,naser2025seismic} }\hskip0pt%DIFAUXCMD
likewise motivate thorough cross-dataset validation rather than single-metric reporting. We therefore evaluate CCVGAE across complementary suites covering clustering accuracy, embedding geometry, and intrinsic manifold quality rather than a single score. Recent methodological developments have addressed single-cell challenges from various angles, including multi-omics representation learning via heterogeneous graphs and graph-based cross-omics integration \mbox{%DIFAUXCMD
\cite{li2024moh,cao2022glue}}\hskip0pt%DIFAUXCMD
. }%DIFDELCMD <

%DIFDELCMD < %%%
\DIFdel{Deep generative models---particularly Variational Autoencoders (VAEs)---are widely used for single-cell analysis, learning low-dimensional, nonlinear embeddings that denoise data and support integration across experiments or modalities \mbox{%DIFAUXCMD
\cite{lopez2018deep,gayoso2021joint,xu2021deep} }\hskip0pt%DIFAUXCMD
. By modeling the biological manifold while accounting for technical variation, these models support a range of downstream analyses. Recent work has expanded }\DIFdelend \DIFaddbegin \DIFadd{and visualisation depend on representations that are stable under stochastic optimisation. Variational autoencoders~\mbox{%DIFAUXCMD
\cite{lopez2018deep,gayoso2021joint,xu2021deep} }\hskip0pt%DIFAUXCMD
have become a standard tool for }\DIFaddend single-cell representation learning\DIFdelbegin \DIFdel{: scJoint performs cross-modal integration between single-cell RNA sequencing (scRNA-seq) and single-cell Assay for Transposase-Accessible Chromatin sequencing (scATAC-seq) through joint embedding learning \mbox{%DIFAUXCMD
\cite{lin2022scjoint}}\hskip0pt%DIFAUXCMD
, and foundation models trained on large-scale single-cell atlases generalize across tissue types and conditions \mbox{%DIFAUXCMD
\cite{cui2024scgpt,theodoris2023geneformer}}\hskip0pt%DIFAUXCMD
. Graph Neural Networks (GNNs), especially Graph Attention Networks (GATs), have been incorporated into single-cell frameworks to model cell--cell similarity via attentional message passing \mbox{%DIFAUXCMD
\cite{wang2021scgnn,liao2023deep}}\hskip0pt%DIFAUXCMD
. Methods such as scVAG \mbox{%DIFAUXCMD
\cite{laghaee2024scvag} }\hskip0pt%DIFAUXCMD
and dual-path graph attention autoencoders \mbox{%DIFAUXCMD
\cite{lv2024ssgate} }\hskip0pt%DIFAUXCMD
combine variational autoencoders with graph attention for single-cell clustering, and recent reviews \mbox{%DIFAUXCMD
\cite{li2025gnnreview,ge2025dlreview} }\hskip0pt%DIFAUXCMD
discuss the role of GNN-based approaches in this area. Dual-feature attention mechanisms and hierarchical network architectures have also been applied to biological pattern recognition \mbox{%DIFAUXCMD
\cite{norouzi2025dft,abbasi2020partwhole,ma2023deepmaps}}\hskip0pt%DIFAUXCMD
. Parallel methodological progress on compact feature-extraction pipelines in adjacent domains---including direction-aware aerial detection \mbox{%DIFAUXCMD
\cite{chen2025lightweight}}\hskip0pt%DIFAUXCMD
, anchor-free lightweight detectors \mbox{%DIFAUXCMD
\cite{zhang2022anchorfree}}\hskip0pt%DIFAUXCMD
, lightweight biometric recognition \mbox{%DIFAUXCMD
\cite{song2022fingervein}}\hskip0pt%DIFAUXCMD
, and Kolmogorov--Arnold network augmentations for medical imaging \mbox{%DIFAUXCMD
\cite{bodner2026kacnet}}\hskip0pt%DIFAUXCMD
---informs efficient attention and bottleneck design principles. These results support the use of attention and dimensional-bottleneck mechanisms in our framework.
}%DIFDELCMD <

%DIFDELCMD < %%%
\DIFdel{How inference-time strategies interact with architectural choices in VAE-based frameworks remains underexplored. Using the posterior mean as the downstream embedding is a known practice in some VAE implementations, but it has not been systematically assessed together with complementary regularization and graph-based encoding across large single-cell benchmarks. The reparameterization trick , while enabling gradient-based optimization, }\DIFdelend \DIFaddbegin \DIFadd{, but the reparameterisation trick used at inference time }\DIFaddend introduces sampling variance that can degrade the stability \DIFdelbegin \DIFdel{, reproducibility, }\DIFdelend and geometric integrity of the latent space, particularly \DIFdelbegin \DIFdel{in }\DIFdelend \DIFaddbegin \DIFadd{for }\DIFaddend sparse single-cell \DIFdelbegin \DIFdel{data }\DIFdelend \DIFaddbegin \DIFadd{modalities }\DIFaddend with low signal-to-noise ratios\DIFdelbegin \DIFdel{\mbox{%DIFAUXCMD
\cite{ding2018reproducibility}}\hskip0pt%DIFAUXCMD
. Embedding variability can affect clustering, trajectory inference, visualization, }\DIFdelend \DIFaddbegin \DIFadd{. Graph neural networks~\mbox{%DIFAUXCMD
\cite{wang2021scgnn,liao2023deep}}\hskip0pt%DIFAUXCMD
, attention-based variants~\mbox{%DIFAUXCMD
\cite{laghaee2024scvag,lv2024ssgate}}\hskip0pt%DIFAUXCMD
, and recent reviews of graph-based single-cell methods~\mbox{%DIFAUXCMD
\cite{li2025gnnreview,ge2025dlreview} }\hskip0pt%DIFAUXCMD
make clear that graph-based encoding can preserve neighbourhood structure but does not by itself address inference-time stability.
}

\DIFadd{Choosing between methods is further complicated by the way the field reports results. Methods are typically published with idiosyncratic evaluation protocols, dataset selections and metric subsets~\mbox{%DIFAUXCMD
\cite{tran2020benchmark}}\hskip0pt%DIFAUXCMD
, so that downstream users cannot directly compare methods that were validated on disjoint cohorts. Even within a single paper, performance is often reported jointly with a fixed encoder backbone and a fixed cell--cell graph, so the algorithmic-core claim is entangled with the architectural axes that surround it~\mbox{%DIFAUXCMD
\cite{duo2018comparison,saelens2019comparison}}\hskip0pt%DIFAUXCMD
. Foundation-model }\DIFaddend and \DIFdelbegin \DIFdel{biological interpretation, limiting the reliability of downstream conclusions}\DIFdelend \DIFaddbegin \DIFadd{cross-modality work~\mbox{%DIFAUXCMD
\cite{cui2024scgpt,theodoris2023geneformer,lin2022scjoint,li2024moh,cao2022glue} }\hskip0pt%DIFAUXCMD
similarly inherits the same evaluation problem}\DIFaddend .

To address these limitations, we introduce \DIFdelbegin %DIFDELCMD < \textbf{CCVGAE}%%%
\DIFdel{, a }\DIFdelend \DIFaddbegin \textbf{scCCVGBen}\DIFadd{, a }\textbf{s}\DIFadd{ingle-}\textbf{c}\DIFadd{ell }\DIFaddend \textbf{C}\DIFdelbegin \DIFdel{entroid-based }\DIFdelend \DIFaddbegin \DIFadd{entroid-}\DIFaddend \textbf{C}oupled \textbf{V}ariational \textbf{G}raph \DIFdelbegin %DIFDELCMD < \textbf{A}%%%
\DIFdel{ttention auto}%DIFDELCMD < \textbf{E}%%%
\DIFdel{ncoder that stabilizes }\DIFdelend \DIFaddbegin \DIFadd{autoencoder }\textbf{Ben}\DIFadd{chmark whose evaluated reference configuration stabilises }\DIFaddend inference while preserving biological structure through three design choices\DIFaddbegin \DIFadd{~\mbox{%DIFAUXCMD
\cite{kipf2016vgae,kingma2014vae,velickovic2018gat}}\hskip0pt%DIFAUXCMD
}\DIFaddend . First, \textbf{Centroid Inference} uses the deterministic posterior mean (the ``centroid'') as the final cell embedding, separating downstream analysis from sampling noise. \DIFdelbegin \DIFdel{Using posterior means }\DIFdelend \DIFaddbegin \DIFadd{Posterior-mean inference }\DIFaddend is a known practice in some \DIFdelbegin \DIFdel{VAE-based models}\DIFdelend \DIFaddbegin \DIFadd{VAE implementations}\DIFaddend ; our contribution is to evaluate this strategy together with complementary architectural components \DIFdelbegin \DIFdel{across a large }\DIFdelend \DIFaddbegin \DIFadd{and to characterise it under a controlled }\DIFaddend benchmark. Second, \DIFdelbegin %DIFDELCMD < \textbf{Coupling Regularization} %%%
\DIFdelend \DIFaddbegin \textbf{Coupling Regularisation} \DIFaddend uses a stochastic training \DIFdelbegin \DIFdel{path with a }\DIFdelend \DIFaddbegin \DIFadd{pathway with an interpretable }\DIFaddend dimensional bottleneck ($d_c < d_z$) and a dual-reconstruction objective \DIFdelbegin \DIFdel{to regularize }\DIFdelend \DIFaddbegin \DIFadd{that regularises }\DIFaddend local latent geometry without \DIFdelbegin \DIFdel{adding noise to final embeddings}\DIFdelend \DIFaddbegin \DIFadd{injecting noise into the final embedding}\DIFaddend . Third, a \textbf{graph-attention encoder} \DIFdelbegin \DIFdel{uses }\DIFdelend \DIFaddbegin \DIFadd{consumes a $k$-nearest-neighbour }\DIFaddend cell--cell similarity \DIFdelbegin \DIFdel{to improve }\DIFdelend \DIFaddbegin \DIFadd{graph in PCA space and improves }\DIFaddend local continuity and biological concordance\DIFdelbegin \DIFdel{; a linear encoder alternative is available when graph construction is infeasible}\DIFdelend .

\DIFdelbegin \DIFdel{This study pursues three objectives: (1)~to quantify the individual and combined contributions of centroid inference, coupling regularization, }\DIFdelend \DIFaddbegin \DIFadd{We assess scCCVGBen within a decoupled benchmark in which the algorithmic core, the encoder backbone, the graph-construction strategy, the dataset cohort and the evaluation suite are evaluated as independent design axes. The benchmark cohort, drawn from the Gene Expression Omnibus and the European Nucleotide Archive, balances scRNA-seq and scATAC-seq equally and spans hematopoiesis, neuronal differentiation, immune populations, tumour microenvironments, organ atlases and developmental time courses (Fig.~\ref{fig02}A--C; Fig.~\ref{fig01}). Within this cohort, we evaluate fourteen graph encoder backbones (Fig.~\ref{fig05}) }\DIFaddend and \DIFdelbegin \DIFdel{graph attention encoding to embedding quality across diverse single-cell modalities; (2)~to benchmark the resulting framework, CCVGAE, against established dimensionality reduction and deep generative baselines on 170 datasets; and (3)~to evaluate the biological interpretability of the learned latent dimensions in hematopoietic systems under physiological stress.}%DIFDELCMD <

%DIFDELCMD < %%%
\DIFdel{Our contributions are as follows:
}%DIFDELCMD < \begin{itemize}
%DIFDELCMD <   \item \textbf{Systematic evaluation of deterministic inference:} %%%
\DIFdel{We evaluate the use of deterministic posterior means as analysis-time embeddings in combination with complementary regularization and graph-based encoding across 170 benchmark datasets. }%DIFDELCMD < \item \textbf{Coupled training for geometric regularity:} %%%
\DIFdel{A dual-reconstruction bottleneck ($d_c < d_z$) refines local latent geometry and information content during training.}%DIFDELCMD < \item \textbf{Attentional graph encoding:} %%%
\DIFdel{A GAT-based encoder leverages cell--cell similarity structure to preserve local neighborhoods and biological manifolds.
  }%DIFDELCMD < \item \textbf{Extensive benchmarking:} %%%
\DIFdel{Evaluation across 170 datasets (55 scRNA-seq, 115 scATAC-seq) demonstrates consistent improvements in clustering , embedding quality}\DIFdelend \DIFaddbegin \DIFadd{five graph-construction strategies (Fig.~\ref{fig06}) under the same evaluation protocol, and score each (dataset, method) pair on a uniform twenty-metric grid spanning clustering compactness, neighbourhood-preserving embedding under UMAP and t-SNE}\DIFaddend , and intrinsic \DIFdelbegin \DIFdel{geometry over established baselines.
  }%DIFDELCMD < \item \textbf{Component-wise interpretability:} %%%
\DIFdel{Latent--gene correlation analysis combined with GO Biological Process enrichment provides transparent mappings from latent dimensions to biological programs.
}%DIFDELCMD < \end{itemize}
%DIFDELCMD <

%DIFDELCMD < %%%
\DIFdel{Centroid anchoring gives consistent improvements over standard stochastic VAEs, and graph attention and coupling regularization add complementary gains. CCVGAE performs favorably on }\DIFdelend \DIFaddbegin \DIFadd{latent geometry. Cross-method comparisons against established }\DIFaddend scRNA-seq\DIFdelbegin \DIFdel{benchmarks and is competitive on }\DIFdelend \DIFaddbegin \DIFadd{~\mbox{%DIFAUXCMD
\cite{higgins2017betavae,chen2018isolating,kim2018disentangling} }\hskip0pt%DIFAUXCMD
and }\DIFaddend scATAC-seq\DIFdelbegin \DIFdel{while improving intrinsic geometric properties. Biological case studies in hematopoietic systems under physiological stress---sleep deprivation , radiation injury, and thrombopoietin-driven megakaryopoiesis---indicate that CCVGAE learns components associated with lineage regulation, inflammatory responses, cell-cycle dynamics, and hemostasis}\DIFdelend \DIFaddbegin \DIFadd{~\mbox{%DIFAUXCMD
\cite{stuart2021atac,ashuach2022peakvi,martens2024poissonvi} }\hskip0pt%DIFAUXCMD
baselines, and three paired hematopoietic case studies (sleep deprivation paired with a gastric cancer microenvironment, cord-blood megakaryopoiesis paired with an aged hematopoietic stem cell time course, and radiation injury paired with the COVID-19 bronchoalveolar lavage immune landscape) probe whether the learned representation is both quantitatively strong and biologically interpretable on data that does not contribute to the benchmark training set}\DIFaddend .

\section{Materials and Methods}

\DIFdelbegin %DIFDELCMD < \subsection{Datasets and Preprocessing}
%DIFDELCMD < %%%
\DIFdelend \DIFaddbegin \subsection{Datasets and provenance}\label{sec:bench_comp}
\DIFaddend

\DIFdelbegin \DIFdel{We assembled a comprehensive benchmark comprising 170 publicly available single-cell datasets: 55 scRNA-seq datasets and 115 scATAC-seq datasets, all retrieved }\DIFdelend \DIFaddbegin \DIFadd{The benchmark cohort is sourced }\DIFaddend from the Gene Expression Omnibus (GEO) and the European Nucleotide Archive (ENA) \DIFdelbegin \DIFdel{. The }\DIFdelend \DIFaddbegin \DIFadd{and is balanced between transcriptomics and chromatin accessibility. Inclusion required (i) availability of raw count matrices for }\DIFaddend scRNA-seq \DIFdelbegin \DIFdel{datasets span diverse biological contexts including hematopoiesis }\DIFdelend \DIFaddbegin \DIFadd{or peak-by-cell matrices for scATAC-seq, either through public GEO supplementary files or through local re-quantification of deposited FASTQ files; (ii) cell-type, condition, donor or time-point annotations of sufficient granularity to support unsupervised evaluation; and (iii) thematic breadth across hematopoiesis and bone-marrow biology}\DIFaddend , neuronal differentiation \DIFaddbegin \DIFadd{and brain atlases, peripheral and tissue-resident immune populations}\DIFaddend , \DIFdelbegin \DIFdel{tumor microenvironments, immune cell populations, }\DIFdelend \DIFaddbegin \DIFadd{organ-level reference atlases, tumour microenvironments, infection and vaccination responses, and developmental time courses. The cohort contains an equal number of scRNA-seq }\DIFaddend and \DIFdelbegin \DIFdel{organ-specific cell atlases. The }\DIFdelend scATAC-seq \DIFdelbegin \DIFdel{datasets cover chromatin accessibility profiling across developmental systems, immune responses, and disease contexts including cancer and neurodegeneration. Full accession numbers and descriptions for all datasets are provided }\DIFdelend \DIFaddbegin \DIFadd{accessions and is dominated by human and mouse studies, with a small set of other-species datasets included for breadth (Fig.~\ref{fig02}A; Fig.~\ref{fig01}). Six restricted-access scATAC-seq submissions from a BCG-vaccination study are retained as redacted benchmark records; their species assignment (}\textit{\DIFadd{Homo sapiens}}\DIFadd{) is inferred from deposited study metadata together with a local re-quantification audit. Cell-count distributions per modality and GEO submission-year timelines are shown in Fig.~\ref{fig02}B--C.
}

\DIFadd{Six additional scRNA-seq datasets are reserved for the biological case studies (Sec.~\ref{sec:case_sd_gastric}--\ref{sec:case_ir_covid}) and excluded from every benchmark statistic: sleep-deprivation bone marrow (GSE280145), a TPO-induced cord-blood megakaryocyte differentiation time course (GSE280270), a radiation-injury hematopoietic stem cell time course (GSE278673), a gastric cancer microenvironment atlas (GSE183904), an aged hematopoietic stem cell dataset (GSE226131) and a COVID-19 bronchoalveolar lavage immune-landscape dataset (GSE145926). Accession numbers and per-dataset metadata are listed }\DIFaddend in the Data Availability section.
\DIFdelbegin \DIFdel{An overview of the benchmark composition---including modality and category distributions, cell and feature count ranges, and tissue source diversity---is shown in Fig.~\ref{fig01}. For scRNA-seq data, preprocessing consists of library-size normalization followed by log1p transformation and selection of }\DIFdelend \DIFaddbegin

\subsection{Preprocessing}\label{sec:preprocess}

\DIFadd{All datasets pass through a uniform preprocessing pipeline that is shared across comparators so that no baseline gains or loses ground because of pipeline differences. Transcriptomic data are processed using }\textit{\DIFadd{Scanpy}}\DIFadd{~\mbox{%DIFAUXCMD
\cite{wolf2018scanpy}}\hskip0pt%DIFAUXCMD
: cells are filtered on minimum gene count and percentage of mitochondrial reads, libraries are normalised to $10^{4}$ counts per cell, expression is $\log(1+x)$-transformed, }\DIFaddend the top 2\DIFdelbegin \DIFdel{,}\DIFdelend \DIFaddbegin {\DIFadd{,}}\DIFaddend 000 highly variable genes \DIFdelbegin \DIFdel{. For scATAC-seq data , a }\DIFdelend \DIFaddbegin \DIFadd{are selected with the Seurat-v3 dispersion estimator, expression values are standardised gene-by-gene, and a 50-dimensional Principal Component Analysis (PCA) embedding is produced for graph construction and for the linear baselines (Sec.~\ref{sec:comparators}). Chromatin-accessibility data are processed following }\textit{\DIFadd{Signac}}\DIFadd{~\mbox{%DIFAUXCMD
\cite{stuart2021atac} }\hskip0pt%DIFAUXCMD
conventions: peak-by-cell matrices are binarised, cells with fewer than 500 accessible peaks are removed, }\DIFaddend TF--IDF \DIFdelbegin \DIFdel{transformation is applied followed by selection of }\DIFdelend \DIFaddbegin \DIFadd{normalisation is applied at the feature level, }\DIFaddend the top 2\DIFdelbegin \DIFdel{,}\DIFdelend \DIFaddbegin {\DIFadd{,}}\DIFaddend 000 highly variable peaks \DIFdelbegin \DIFdel{. Dimensionality reduction to 50 components is performed via Principal Component Analysis (PCA; scRNA-seq) or }\DIFdelend \DIFaddbegin \DIFadd{are retained and }\DIFaddend Latent Semantic Indexing (LSI\DIFdelbegin \DIFdel{; scATAC-seq) . Three additional scRNA-seq datasets (GSE280145, GSE278673, GSE280270) are used exclusively for the biological case studies and are not included in the quantitative benchmarks. Because this study evaluates }\DIFdelend \DIFaddbegin \DIFadd{) produces an analogous 50-dimensional embedding. Because the benchmark targets }\DIFaddend unsupervised representation learning\DIFdelbegin \DIFdel{rather than supervised prediction, no target variable is optimized during model training. Publicly available cell-type annotations are used only after training as external reference labels for evaluation}\DIFdelend , \DIFdelbegin \DIFdel{and correlations between input features and a supervised target variable are therefore not part of the study design. }%DIFDELCMD <

%DIFDELCMD < %%%
%DIF <  OMITTED_figureSTAR_FOR_DIFF_BUILD_STABILITY
%DIFDELCMD <

%DIFDELCMD < \subsection{Model Architecture}
%DIFDELCMD < %%%
\DIFdelend \DIFaddbegin \DIFadd{no cell-type label is supplied to any model at training time; published annotations are reserved exclusively for post-hoc external evaluation.
}\DIFaddend

\DIFdelbegin \subsubsection{\DIFdel{Overview}}
%DIFAUXCMD
\addtocounter{subsubsection}{-1}%DIFAUXCMD
\DIFdelend \DIFaddbegin \subsection{Cell--cell graph construction}\label{sec:graphs}
\DIFaddend

\DIFdelbegin \DIFdel{The core innovation of CCVGAE is its dual-path architecture, which decouples stable embeddinggeneration from stochastic variational training. As illustrated }\DIFdelend \DIFaddbegin \DIFadd{The default cell--cell similarity graph $\mathbf{A}$ is a 15-nearest-neighbour graph computed in the 50-dimensional PCA or LSI embedding under Euclidean distance, symmetrised by union and stored as a sparse adjacency matrix; self-loops are added explicitly because the graph attention layer assumes their presence. Four alternative graph constructions share the same 50-dimensional embedding and are calibrated to a comparable average node degree so that downstream message passing operates over graphs of similar density: kNN-cosine (cosine distance on the same embedding), SNN (shared-nearest-neighbour reweighting of the kNN edges), mutual-kNN (the symmetric intersection of two kNN graphs) and Gaussian-threshold (a Gaussian-kernel similarity thresholded by a fixed quantile). All five constructions are evaluated against the kNN-Euclidean baseline }\DIFaddend in Fig.~\DIFdelbegin \DIFdel{\ref{Architecture}, the framework begins with a preprocessing step where single-cell data are represented as cell feature vectors $\mathbf{X} \in \mathbb{R}^{N \times D}$ (where }\DIFdelend \DIFaddbegin \DIFadd{\ref{fig06}.
}

\subsection{Interactive companion website}\label{sec:site}

\DIFadd{A static, search-indexable website (Fig.~\ref{fig01}) accompanies the manuscript and is generated from the same dataset cohort and reconciled result tables that drive the figures. The site provides cards for each dataset (accession, organism, tissue, condition, cell and feature counts, and per-method per-metric scores), each method (family, configuration, citation, and per-metric scores) and each metric (formal definition, implementation reference and directionality convention). Eight browsable views render in canonical reading order: an atlas entry point summarising the cohort (Fig.~\ref{fig01}A); a tissue, species and modality composition view (Fig.~\ref{fig01}B); cohort-scale charts of cell-count distribution, submission-year timeline, tissue~$\times$~modality stacking and top studies by cell count (Fig.~\ref{fig01}C); a dataset browser with filters (Fig.~\ref{fig01}D); a per-dataset metadata page (Fig.~\ref{fig01}E); a comparator catalogue with one card per benchmark-eligible method including the fourteen graph encoders (Fig.~\ref{fig01}F); a per-method metadata page (Fig.~\ref{fig01}G); and a metric taxonomy with the three-suite definitions and directionality conventions (Fig.~\ref{fig01}H). Restricted records use public-safe identifiers so the website remains internally consistent with the cohort without exposing nonpublic accession details.
}


%DIF >  OMITTED_figureSTAR_FOR_DIFF_BUILD_STABILITY


\subsection{Model architecture}\label{sec:architecture}

\DIFadd{The scCCVGBen reference model has a dual-path architecture that decouples stable embedding generation from stochastic variational training (Fig.~\ref{fig02}D). Let $\mathbf{X}\in\mathbb{R}^{N\times D}$ denote the cell-by-feature matrix for a cohort of }\DIFaddend $N$ \DIFdelbegin \DIFdel{is the number of cells and }\DIFdelend \DIFaddbegin \DIFadd{cells with }\DIFaddend $D$ \DIFdelbegin \DIFdel{is the number of genes or features)and a corresponding }\DIFdelend \DIFaddbegin \DIFadd{preprocessed features and let $\mathbf{A}\in\mathbb{R}^{N\times N}$ denote the }\DIFaddend cell--cell \DIFdelbegin \DIFdel{similarity graph $\mathbf{A} \in \mathbb{R}^{N \times N}$ is constructed.Throughout this paper, we use the term }\textit{\DIFdel{cell feature vector}} %DIFAUXCMD
\DIFdel{to refer to the high-dimensional molecular profile of a single cell (e.g., gene expression counts or peak accessibility values); in the graph context, each cell corresponds to a node and its molecular profile constitutes the node features. A Graph Attention Network (GAT) encoder then processes subgraphs sampled from these inputs to learn a posterior distribution $q_\phi(z|x)$ }\DIFdelend \DIFaddbegin \DIFadd{similarity graph. A graph encoder $f_\theta$ maps $(\mathbf{X},\mathbf{A})$ to the parameters of an approximate posterior }\DIFaddend over a $d_z$-dimensional latent space\DIFdelbegin \DIFdel{, parameterized by a mean vector $\mu_z(x) \in \mathbb{R}^{d_z}$ and a standard deviation vector $\sigma_z(x) \in \mathbb{R}^{d_z}$, where $d_z$ denotes the latent dimensionality (default $d_z = 10$) }\DIFdelend \DIFaddbegin \DIFadd{:
}\begin{equation}
\DIFadd{f_\theta(\mathbf{X},\mathbf{A}) \;=\; \bigl(\mu_z(\mathbf{X},\mathbf{A}),\,\sigma_z(\mathbf{X},\mathbf{A})\bigr),\qquad q_\phi(z\mid \mathbf{X},\mathbf{A}) \;=\; \mathcal{N}\!\bigl(\mu_z,\,\mathrm{diag}(\sigma_z^2)\bigr),
\label{eq:posterior}
}\end{equation}
\DIFadd{with $\mu_z, \sigma_z \in \mathbb{R}^{d_z}$ produced per cell. Three architectural mechanisms then act on this posterior, and the component-wise ablation (Fig.~\ref{fig03}) dissects their contributions separately}\DIFaddend .

\DIFdelbegin %DIFDELCMD < \textbf{Deterministic Path}%%%
\DIFdel{: For all downstream analyses, }\DIFdelend \DIFaddbegin \textit{\DIFadd{Centroid Inference.}} \DIFadd{At inference time }\DIFaddend the posterior mean \DIFdelbegin \DIFdel{$\mu_{\text{centroid}} = \mu_z(x)$ serves as the final }\DIFdelend \DIFaddbegin \DIFadd{is exported directly as the }\DIFaddend cell embedding,
\DIFdelbegin \DIFdel{eliminating sampling noise in clustering, visualization, }\DIFdelend \DIFaddbegin \begin{equation}
\DIFadd{\mu_{\text{centroid}} \;=\; \mu_z(\mathbf{X},\mathbf{A}),
\label{eq:centroid}
}\end{equation}
\DIFadd{in place of a reparameterised sample. This removes the sampling variance that would otherwise propagate into clustering, projection }\DIFaddend and trajectory inference\DIFaddbegin \DIFadd{; conditional on a trained encoder, $\mu_{\text{centroid}}$ is a deterministic function of $(\mathbf{X},\mathbf{A})$ and is therefore reproducible across runs}\DIFaddend .

\DIFdelbegin %DIFDELCMD < \textbf{Stochastic Path}%%%
\DIFdel{: For model training, the stochastic path employs the reparameterization trick to sample from the posterior: $z = \mu_z + \sigma_z \odot \varepsilon$, where $\varepsilon \sim \mathcal{N}(0,I)$. This path incorporates a dual-reconstruction mechanism to enforce a well-structured latent space. The sampled vector $z$ first enters a }%DIFDELCMD < \textbf{bottleneck path}%%%
\DIFdel{, where it is linearly compressed to a lower-dimensional representation $z_c \in \mathbb{R}^{d_c}$ (with bottleneck dimension $d_c < d_z$; default $d_c = 5$) via learnable parameters $W_{\text{comp}} \in \mathbb{R}^{d_c \times d_z}$ and bias $b_{\text{comp}} \in \mathbb{R}^{d_c}$}\DIFdelend \DIFaddbegin \textit{\DIFadd{Coupling Regularisation.}} \DIFadd{During training, reparameterised samples $z\sim q_\phi(z\mid\mathbf{X},\mathbf{A})$ pass through a bottleneck pathway. A compression operator $g_\downarrow:\mathbb{R}^{d_z}\!\to\!\mathbb{R}^{d_c}$ with $d_c<d_z$ produces $z_c=g_\downarrow(z)$, and a reconstruction operator $g_\uparrow:\mathbb{R}^{d_c}\!\to\!\mathbb{R}^{d_z}$ returns $z_r=g_\uparrow(z_c)$}\DIFaddend . A feature decoder \DIFdelbegin \DIFdel{then generates an initial reconstruction $\hat{x}_r$, and the first reconstruction loss, }\DIFdelend \DIFaddbegin \DIFadd{$p_\psi$ produces $\hat{\mathbf{X}}$, and a graph decoder produces $\hat{\mathbf{A}}$. The composite training objective is
}\begin{equation}
\DIFadd{\mathcal{L}
\;=\;
w_{\text{recon}}\,\underbrace{\bigl\lVert\hat{\mathbf{X}}-\mathbf{X}\bigr\rVert_{2}^{2}}_{\mathcal{L}_{\text{recon}}}
\;+\;
w_{\text{irecon}}\,\underbrace{\bigl\lVert z_r - z\bigr\rVert_{2}^{2}}_{\mathcal{L}_{\text{irecon}}}
\;+\;
w_{\text{KL}}\,\underbrace{D_{\text{KL}}\!\bigl(q_\phi(z\mid\mathbf{X},\mathbf{A})\,\big\|\,\mathcal{N}(0,\mathbf{I})\bigr)}_{\mathcal{L}_{\text{KL}}}
\;+\;
w_{\text{adj}}\,\underbrace{\bigl\lVert\hat{\mathbf{A}}-\mathbf{A}\bigr\rVert_{F}^{2}}_{\mathcal{L}_{\text{adj}}},
\label{eq:loss}
}\end{equation}
\DIFadd{with default weights $w_{\text{recon}}=w_{\text{irecon}}=w_{\text{KL}}=w_{\text{adj}}=1.0$. The inner reconstruction }\DIFaddend $\mathcal{L}_{\text{irecon}}$ \DIFdelbegin \DIFdel{, is computed. Subsequently, in the }%DIFDELCMD < \textbf{primary path}%%%
\DIFdel{, the compressed vector $z_c$ is
reconstructed back to the original latent dimension via $W_{\text{recon}} \in \mathbb{R}^{d_z \times d_c}$ and $b_{\text{recon}} \in \mathbb{R}^{d_z}$, yielding $z_r \in \mathbb{R}^{d_z}$. A second pass through the feature decoder generates the final reconstruction $\hat{x}$, used to calculate $\mathcal{L}_{\text{recon}}$}\DIFdelend \DIFaddbegin \DIFadd{encourages a low-dimensional information-bearing subspace inside $z$ without injecting noise into the centroid embedding used at inference}\DIFaddend .

\DIFdelbegin \DIFdel{The overall training objective combines a KL divergence term ($\mathcal{L}_{\text{kl}}$), which regularizes the posterior against a prior $p(z)$, with the two reconstruction losses ($\mathcal{L}_{\text{recon}}$ and $\mathcal{L}_{\text{irecon}}$). These reconstruction losses are calculated using a modality-appropriate likelihood: a negative binomial (NB) distribution for count-based scRNA-seq data, or a mean squared error (MSE) loss for TF--IDF}\DIFdelend \DIFaddbegin \textit{\DIFadd{Graph-Attention Encoder.}} \DIFadd{The default encoder is a two-layer Graph Attention Network with four attention heads, which uses cell--cell similarity to preserve local neighbourhood structure on the learned manifold. The encoder is one of fourteen alternatives that share the input}\DIFaddend /\DIFdelbegin \DIFdel{LSI-transformed scATAC-seq data}\DIFdelend \DIFaddbegin \DIFadd{output signature $f_\theta(\mathbf{X},\mathbf{A})\rightarrow(\mu_z,\sigma_z)$ and the same training schedule, so the encoder robustness analysis (Fig.~\ref{fig05}) measures architectural sensitivity at fixed loss and schedule}\DIFaddend .


\DIFdelbegin \subsubsection{\DIFdel{Cell--Cell Graph Construction}}
%DIFAUXCMD
\addtocounter{subsubsection}{-1}%DIFAUXCMD
%DIFDELCMD < \label{sec:graph_construction}
%DIFDELCMD < %%%
\DIFdelend %DIF >  OMITTED_figureSTAR_FOR_DIFF_BUILD_STABILITY
\DIFaddbegin


\subsection{Encoder backbones and graph constructions evaluated}\label{sec:registries}
\DIFaddend

The \DIFaddbegin \DIFadd{encoder robustness analysis (Fig.~\ref{fig05}) varies the encoder $f_\theta$ across fourteen alternatives that span four implementation families. The attention family contains GAT (the default), GATv2, a Transformer-style attention encoder and SuperGAT. The convolutional family contains GCN, SAGE, GraphConv and Chebyshev. The propagation family contains Topology Adaptive (TAG), ARMA, Simple Graph (SG) and Stacked Simple Graph (SSG). The message-passing family contains the Graph Isomorphism Network (GIN) and EdgeConv. Every backbone shares the same encoder signature $f_\theta(\mathbf{X},\mathbf{A})\rightarrow(\mu_z,\sigma_z)$ and the same training schedule, so a difference in the encoder panel cannot be attributed to a difference in loss weighting or training length.
}

\DIFadd{The graph robustness analysis (Fig.~\ref{fig06}) holds the algorithmic core fixed and varies the }\DIFaddend cell--cell \DIFdelbegin \DIFdel{similarity graph $\mathbf{A}$ is constructed using the following procedure:
}%DIFDELCMD < \begin{enumerate}
%DIFDELCMD <     \item \textbf{Feature preprocessing:} %%%
\DIFdel{For scRNA-seq data, raw count matrices are normalized (library-size normalization followed by log1p transformation) and the top 2,000 highly variable genes are selected. For scATAC-seq data, a TF--IDF transformation is applied followed by selection of the top 2, 000 highly variable peaks; these features are then used for dimensionality reduction. }%DIFDELCMD < \item \textbf{Dimensionality reduction:} %%%
\DIFdel{PCA reduces the feature space to 50 dimensions (for scRNA-seq)or LSI extracts 50 components (for scATAC-seq). }%DIFDELCMD < \item \textbf{$k$-nearest neighbor graph:} %%%
\DIFdel{For each cell $i$, we identify its $k=15$ nearest neighbors based on Euclidean distance computed in the reduced PCA/LSI space (i.e. , $\ell_2$ norm on the }\DIFdelend \DIFaddbegin \DIFadd{similarity graph between five alternatives that share the same }\DIFaddend 50-dimensional \DIFdelbegin \DIFdel{vectors).
}%DIFDELCMD < \item \textbf{Symmetrization:} %%%
\DIFdel{The graph is made symmetric: $A_{ij} = A_{ji} = 1$ if cell $i$ is among the $k$-nearest neighbors of cell $j$ or vice versa.
    }%DIFDELCMD < \item \textbf{Self-loops:} %%%
\DIFdel{Self-connections are added ($A_{ii} = 1$)to ensure each cell's own features contribute to its updated representation.
}%DIFDELCMD < \end{enumerate}
%DIFDELCMD < %%%
\DIFdel{The resulting graph is unweighted, undirected, and sparse, with an average degree of approximately $2k$. This construction follows standard practices in single-cell graph-based methods \mbox{%DIFAUXCMD
\cite{wolf2018scanpy}}\hskip0pt%DIFAUXCMD
.
}\DIFdelend \DIFaddbegin \DIFadd{PCA or LSI embedding: kNN-Euclidean (the default; Euclidean distance), kNN-cosine (cosine distance), SNN (shared-nearest-neighbour reweighting of the kNN edges), mutual-kNN (symmetric intersection of two kNN graphs) and Gaussian-threshold (Gaussian-kernel similarity thresholded so the average node degree matches the kNN baseline). Because both robustness analyses hold every other component fixed, an effect observed in either is attributable to the single architectural choice that was varied.
}\DIFaddend

\DIFdelbegin \subsubsection{\DIFdel{Model Variants and Ablation Hierarchy}}
%DIFAUXCMD
\addtocounter{subsubsection}{-1}%DIFAUXCMD
%DIFDELCMD < \label{sec:model_variants}
%DIFDELCMD < %%%
\DIFdelend \DIFaddbegin \subsection{Baselines for cross-method comparison}\label{sec:comparators}
\DIFaddend

\DIFdelbegin \DIFdel{To systematically evaluate the contribution of each architectural component, we define four model variants:
}%DIFDELCMD < \begin{itemize}
\begin{itemize}%DIFAUXCMD
%DIFDELCMD <     \item \textbf{VAE (baseline):} %%%
\item%DIFAUXCMD
\DIFdel{Standard variational autoencoder with multi-layer perceptron (MLP) encoder; uses stochastic samples for both training and inference. This baselineconfiguration, which draws samples from $q_\phi(z|x)$ at inference time, serves as a controlled reference to isolate the effect of each architectural component. We note that some VAE implementations in practice already use posterior means at inference; our ablation hierarchy is designed to disentangle the contributions of deterministic inference, coupling, and graph attention in a controlled setting.
    }%DIFDELCMD < \item \textbf{CenVAE:} %%%
\item%DIFAUXCMD
\DIFdel{VAE + }\textit{\DIFdel{Centroid Inference}}%DIFAUXCMD
\DIFdel{; uses posterior mean $\mu_z(x)$ as the final embedding instead of stochastic samples.
    }%DIFDELCMD < \item \textbf{CouVAE:} %%%
\item%DIFAUXCMD
\DIFdel{VAE + }\textit{\DIFdel{Coupling regularization}}%DIFAUXCMD
\DIFdel{; incorporates the dual-reconstruction bottleneck but retains stochastic inference.
    }%DIFDELCMD < \item \textbf{GAT-VAE:} %%%
\item%DIFAUXCMD
\DIFdel{VAE with Graph Attention encoder; replaces MLP with GAT but uses stochastic inference.
    }%DIFDELCMD < \item \textbf{CCVGAE (full model):} %%%
\item%DIFAUXCMD
\DIFdel{Combines all three innovations---Centroid Inference + Coupling + Graph Attention.
}
\end{itemize}%DIFAUXCMD
%DIFDELCMD < \end{itemize}
%DIFDELCMD < %%%
\DIFdel{This hierarchy enables ablation studies that isolate the effect of each component.
An overview of the full experimental design, including all baseline methods and evaluation metrics, is presented in }\DIFdelend \DIFaddbegin \DIFadd{On scRNA-seq (}\DIFaddend Fig.~\DIFdelbegin \DIFdel{\ref{fig02}.
}%DIFDELCMD <

%DIFDELCMD < %%%
\DIFdel{For external benchmarking, we selected comparison methods to span three complementary groups: (i) ~classical dimensionality-reduction baselines (PCA}\DIFdelend \DIFaddbegin \DIFadd{\ref{fig07}) we compare scCCVGBen against twelve external baselines drawn from two families. The classical/linear family contains Principal Component Analysis (PCA)}\DIFaddend , Kernel PCA (KPCA), Independent Component Analysis (ICA), Factor Analysis (FA), Non-negative Matrix \DIFdelbegin \DIFdel{Factorization }\DIFdelend \DIFaddbegin \DIFadd{Factorisation }\DIFaddend (NMF), Truncated \DIFdelbegin \DIFdel{SVD }\DIFdelend \DIFaddbegin \DIFadd{Singular Value Decomposition }\DIFaddend (TSVD)\DIFdelbegin \DIFdel{, Dictionary Learning }\DIFdelend \DIFaddbegin \DIFadd{~\mbox{%DIFAUXCMD
\cite{halko2011finding} }\hskip0pt%DIFAUXCMD
and a Deep Iterative Cross-Linkage representation }\DIFaddend (DICL)\DIFdelbegin \DIFdel{), (ii)~closely related deep generative baselines (scVI, DIPVAE, InfoVAE, TCVAE, HighBetaVAE), and (iii)~}\DIFdelend \DIFaddbegin \DIFadd{~\mbox{%DIFAUXCMD
\cite{wang2018dicl}}\hskip0pt%DIFAUXCMD
. The deep generative family contains scVI~\mbox{%DIFAUXCMD
\cite{lopez2018deep}}\hskip0pt%DIFAUXCMD
, DIPVAE~\mbox{%DIFAUXCMD
\cite{kumar2018dipvae,kim2018disentangling}}\hskip0pt%DIFAUXCMD
, InfoVAE~\mbox{%DIFAUXCMD
\cite{zhao2019infovae,xu2021deep}}\hskip0pt%DIFAUXCMD
, $\beta$-TCVAE~\mbox{%DIFAUXCMD
\cite{chen2018isolating} }\hskip0pt%DIFAUXCMD
and HighBetaVAE~\mbox{%DIFAUXCMD
\cite{higgins2017betavae}}\hskip0pt%DIFAUXCMD
. On scATAC-seq (Fig.~\ref{fig08}) we use three }\DIFaddend modality-specific \DIFdelbegin \DIFdel{scATAC-seq baselines(LSI, PeakVI, PoissonVI). This design lets us compare CCVGAE against widely used linear baselines, VAE-family alternatives, and assay-specific methods that are standard in current single-cell practice.
}%DIFDELCMD <

%DIFDELCMD < \textbf{Terminological note:} %%%
\DIFdel{We use the abbreviation ``GAT-VAE'' for the graph-attention encoder variant to distinguish it from the classical Variational Graph Autoencoder (VGAE) formulation of Kipf and Welling, which reconstructs the adjacency matrix. In our model, the GAT encoder is used solely to produce enriched node (cell) representations; the decoder reconstructs feature data, not graph structure.
}%DIFDELCMD <

%DIFDELCMD < %%%
%DIF < %%%%% Architecture %%%%%%
%DIFDELCMD <

%DIFDELCMD < %%%
%DIF <  OMITTED_figureSTAR_FOR_DIFF_BUILD_STABILITY
%DIFDELCMD <

%DIFDELCMD < %%%
\subsubsection{\DIFdel{Encoder and Posterior Learning}}
%DIFAUXCMD
\addtocounter{subsubsection}{-1}%DIFAUXCMD
%DIFDELCMD <

%DIFDELCMD < %%%
\DIFdel{The CCVGAE encoder, parameterized by $\phi$, maps an input cell's feature vector $x \in \mathbb{R}^D$ to a diagonal Gaussian posterior distribution in the latent space $\mathbb{R}^{d_z}$:
}\begin{displaymath}
	\DIFdel{q_\phi(z|x) = \mathcal{N}(z | \mu_z(x), \operatorname{diag}(\sigma_z^2(x)))
}\end{displaymath}%DIFAUXCMD
%DIFDELCMD <

%DIFDELCMD < %%%
\DIFdel{From this posterior, CCVGAE derives two distinct representations for its dual-path mechanism:
}%DIFDELCMD < \begin{enumerate}
\begin{enumerate}%DIFAUXCMD
%DIFDELCMD < 	\item \textbf{Posterior Centroid (Final Embedding)}%%%
\item%DIFAUXCMD
\DIFdel{: The deterministic mean of the posterior is used as the final, stable cell embedding for all downstream tasks. This is the core of Centroid Inference.
}\begin{displaymath}
		\DIFdel{\mu_{\text{centroid}} = \mu_z(x)
	}\end{displaymath}%DIFAUXCMD
%DIFDELCMD <

%DIFDELCMD < 	\item \textbf{Stochastic Sample (for Training)}%%%
\item%DIFAUXCMD
\DIFdel{: A random sample is drawn using the reparameterization trick. This sample is used exclusively during training to regularize the model via the variational objective.
	}\begin{displaymath}
		\DIFdel{z = \mu_z(x) + \sigma_z(x) \odot \varepsilon, \quad \text{where} \quad \varepsilon \sim \mathcal{N}(0, I)
	}\end{displaymath}%DIFAUXCMD

\end{enumerate}%DIFAUXCMD
%DIFDELCMD < \end{enumerate}
%DIFDELCMD < %%%
\DIFdel{The standard deviation $\sigma_z(x)$ is parameterized via a softplus function for numerical stability: $\sigma_z(x) = \text{softplus}(s_z(x)) + 10^{-6}$, where $s_z(x)$ is the raw log-variance output from the encoder.
}%DIFDELCMD <

%DIFDELCMD < %%%
\subsubsection{\DIFdel{Coupling Regularization via a Dual-Reconstruction Bottleneck}}
%DIFAUXCMD
\addtocounter{subsubsection}{-1}%DIFAUXCMD
%DIFDELCMD <

%DIFDELCMD < %%%
\DIFdel{The Coupling mechanism is a bottleneck layer that regularizes the latent space. The stochastically sampled vector $z$ passes through a linear compression-reconstruction pathway:
}\begin{displaymath}
	\DIFdel{z_c = W_{\text{comp}} z + b_{\text{comp}}
}\end{displaymath}%DIFAUXCMD
\begin{displaymath}
	\DIFdel{z_r = W_{\text{recon}} z_c + b_{\text{recon}}
}\end{displaymath}%DIFAUXCMD
\DIFdel{where $z_c \in \mathbb{R}^{d_c}$ with a compressed dimension $d_c < d_z$}\DIFdelend \DIFaddbegin \DIFadd{baselines: LSI~\mbox{%DIFAUXCMD
\cite{stuart2021atac}}\hskip0pt%DIFAUXCMD
, PeakVI~\mbox{%DIFAUXCMD
\cite{ashuach2022peakvi} }\hskip0pt%DIFAUXCMD
and PoissonVI~\mbox{%DIFAUXCMD
\cite{martens2024poissonvi}}\hskip0pt%DIFAUXCMD
. Baseline scores are reconciled to the cohort before each pairwise test, so the paired Wilcoxon comparisons are computed on the same datasets on which scCCVGBen is trained and evaluated}\DIFaddend .

\DIFdelbegin \DIFdel{The feature decoder is applied twice: once to the primary latent sample $z$ to generate $\hat{x}$, and a second time to the bottleneck-processed sample $z_r$ to generate $\hat{x}_r$. This dual-reconstruction objective forces the bottleneck to preserve reconstruction-critical information and regularizes latent space structure.
}%DIFDELCMD <

%DIFDELCMD < %%%
\subsubsection{\DIFdel{Encoder Architectures}}
%DIFAUXCMD
\addtocounter{subsubsection}{-1}%DIFAUXCMD
\DIFdelend \DIFaddbegin \subsection{Evaluation metrics}\label{sec:metrics}
\DIFaddend

\DIFdelbegin \DIFdel{CCVGAE supports multiple encoder backbones for different data modalities.
}\DIFdelend \DIFaddbegin \DIFadd{Each (dataset, method) pair is evaluated on a uniform grid of twenty metrics organised into three complementary suites that probe different aspects of representation quality.
}\DIFaddend

\DIFdelbegin \paragraph{\DIFdel{Linear Encoder (MLP)}}
%DIFAUXCMD
\addtocounter{paragraph}{-1}%DIFAUXCMD
\DIFdel{For non-graph data, the encoder is implemented as a standard MLP with $L=2$ hidden layers and hidden dimension $H=128$. Given an input $x \in \mathbb{R}^D$ and letting $H^{(0)} = x$:
}\begin{displaymath}
	\DIFdel{H^{(l)} = \text{ReLU}(H^{(l-1)} W^{(l)} + b^{(l)}), \quad l=1,\ldots,L
}\end{displaymath}%DIFAUXCMD
\DIFdel{Each hidden layer applies a linear transformation followed by a ReLU activation function ($\text{ReLU}(x) = \max(0,x)$). Dropout (rate $=0.05$) is applied after each hidden layer during training to mitigate overfitting. The posterior parameters are then computed from the final hidden layer $H^{(L)} \in \mathbb{R}^{128}$ via separate linear projections(without activation):
}\begin{displaymath}
	\DIFdel{\mu_z(x) = H^{(L)} W_{\mu} + b_{\mu}, \quad s_z(x) = H^{(L)} W_s + b_s
}\end{displaymath}%DIFAUXCMD
\DIFdel{where $W_{\mu}, W_s \in \mathbb{R}^{128 \times d_z}$ and $b_{\mu}, b_s \in \mathbb{R}^{d_z}$}\DIFdelend \DIFaddbegin \DIFadd{The first suite measures clustering compactness on Leiden clusters obtained from the learned embedding at a fixed resolution. It reports the Average Silhouette Width (ASW), the Davies--Bouldin index (DAV; reported with the convention that lower is better, sign-flipped at visualisation time so that all suites point in the same direction), and the Calinski--Harabasz index (CAL). The maximum cluster sizes attained in the two-dimensional UMAP and t-SNE projections, $K_{\max}^{\text{UMAP}}$ and $K_{\max}^{\text{tSNE}}$, are reported alongside as auxiliary statistics}\DIFaddend .

\DIFdelbegin \paragraph{\DIFdel{Graph Attention Network (GAT) Encoder}}
%DIFAUXCMD
\addtocounter{paragraph}{-1}%DIFAUXCMD
\DIFdel{For graph-structured data, the model employs a Graph Attention Network (GAT) encoder with $L=2$ layers and $K=4$ attention heads per layer}\DIFdelend \DIFaddbegin \DIFadd{The second suite measures neighbourhood preservation under non-linear projection to two dimensions}\DIFaddend . For each \DIFdelbegin \DIFdel{subgraph supplied by the data loader, the GAT layers update node representations by attending over their local neighborhoods. Specifically, for each attention head $k$, the attention coefficient between nodes $i$ and $j$ is computed as:
}\begin{displaymath}
	\DIFdel{\alpha_{ij}^{(l,k)} = \frac{\exp\!\left(\text{LeakyReLU}\!\left(\mathbf{a}^{(l,k)\top} [W^{(l,k)} H^{(l)}_i \,\|\, W^{(l,k)} H^{(l)}_j]\right)\right)}{\sum_{m \in \mathcal{N}_i \cup \{i\}} \exp\!\left(\text{LeakyReLU}\!\left(\mathbf{a}^{(l,k)\top} [W^{(l,k)} H^{(l)}_i \,\|\, W^{(l,k)} H^{(l)}_m]\right)\right)}
}\end{displaymath}%DIFAUXCMD
\DIFdel{where $\|$ denotes concatenation.
The outputs of the $K$ attention heads are concatenated (for intermediate layers) or averaged (for the final layer):
}\begin{displaymath}
	\DIFdel{H^{(l+1)}_i = \underset{k=1}{\overset{K}{\|}} \; \sigma\!\left(\sum_{j \in \mathcal{N}_i \cup \{i\}} \alpha_{ij}^{(l,k)} W^{(l,k)} H^{(l)}_j \right)
}\end{displaymath}%DIFAUXCMD
\DIFdel{where $\mathcal{N}_i$ is the set of neighbors of node $i$ in the current subgraph, $\sigma$ is the ELU activation function, and the per-head hidden dimension is $H/K = 128/4 = 32$. Dropout (rate$=0.05$) is applied to both the attention coefficients and the cell feature vectors (node features) during training. After $L$ GAT layers, }\DIFdelend \DIFaddbegin \DIFadd{of UMAP and t-SNE we compute the distance correlation against the high-dimensional input, the local and global co-ranking quality scores, an aggregate overall quality score and the auxiliary $K_{\max}$ statistic, giving ten metrics in total.
}

\DIFadd{The third suite characterises intrinsic latent geometry without reference to a two-dimensional projection. It reports the intrinsic manifold dimensionality estimated from the spectrum of }\DIFaddend the \DIFdelbegin \DIFdel{final cell representations $H^{(L)}_i \in \mathbb{R}^{128}$ are projected to the posterior parameters via linear layers:
}\begin{displaymath}
	\DIFdel{\mu_z(x_i) = H^{(L)}_i W_{\mu} + b_{\mu}, \quad s_z(x_i) = H^{(L)}_i W_s + b_s
}\end{displaymath}%DIFAUXCMD
\DIFdelend \DIFaddbegin \DIFadd{latent covariance, the spectral decay rate, the participation ratio, an anisotropy score derived from the eigenvalue distribution, a trajectory-directionality score that quantifies how well a one-dimensional ordering respects pairwise distances, a noise-resilience score obtained by perturbing the input and measuring embedding stability, and an aggregate intrinsic-overall score.
}\DIFaddend

\DIFdelbegin \subsubsection{\DIFdel{Scalable Graph Training via Subgraph Sampling}}
%DIFAUXCMD
\addtocounter{subsubsection}{-1}%DIFAUXCMD
\DIFdelend \DIFaddbegin \subsection{Statistical analysis}\label{sec:stats}
\DIFaddend

\DIFdelbegin \DIFdel{To ensure scalability when using the GAT encoder, the model is trained on small subgraphs sampled from the full cell--cell graph. In each training step:
}%DIFDELCMD < \begin{enumerate}
\begin{enumerate}%DIFAUXCMD
%DIFDELCMD < 	\item \textbf{Random Node Sampling}%%%
\item%DIFAUXCMD
\DIFdel{: A fixed number of nodes, $N_s$, are randomly selected without replacement. }%DIFDELCMD < \item \textbf{Subgraph Inducement}%%%
\item%DIFAUXCMD
\DIFdel{: An induced subgraph $S$ is formed, containing the $N_s$ nodes and all edges from the original graph connecting them. }
\end{enumerate}%DIFAUXCMD
%DIFDELCMD < \end{enumerate}
%DIFDELCMD < %%%
\DIFdel{The GAT encoder and loss functions are then applied exclusively to this subgraph $S$}\DIFdelend \DIFaddbegin \DIFadd{All cross-method comparisons in this manuscript use paired Wilcoxon signed-rank tests with Holm-corrected pairwise post-hoc analysis, computed per metric on the set of datasets that is paired between the reference method and each comparator. Holm correction is applied within each figure so that the reported $p$-values are corrected against the comparison family rendered in the same figure rather than against the global pool of pairwise tests. Confidence intervals on mean differences are obtained by paired bootstrap with 5}{\DIFadd{,}}\DIFadd{000 resamples under a fixed random seed}\DIFaddend .

\DIFdelbegin \subsubsection{\DIFdel{Feature Decoder and Modality-Specific Likelihood}}
%DIFAUXCMD
\addtocounter{subsubsection}{-1}%DIFAUXCMD
\DIFdelend \DIFaddbegin \subsection{Default hyperparameters and training}
\DIFaddend

\DIFdelbegin \DIFdel{The feature decoder, implemented as }\DIFdelend \DIFaddbegin \DIFadd{Hyperparameters were selected on a held-out validation subset and held fixed thereafter. A 10-dimensional latent space is used with a 5-dimensional coupling bottleneck; the encoder is }\DIFaddend a two-layer \DIFdelbegin \DIFdel{MLP (hidden dimension }\DIFdelend \DIFaddbegin \DIFadd{Graph Attention Network with four heads on a hidden dimension of }\DIFaddend 128 \DIFdelbegin \DIFdel{, ReLU activations), reconstructs the input features $x$ from a latent vector $z'$. For the dual-reconstruction mechanism, this process is applied to both the primary sample $z$ and the bottlenecked sample $z_r$. For scRNA-seq data, expression counts are modeled using a Negative Binomial (NB) distribution to account for overdispersion. For scATAC-seq data, TF--IDF transformation followed by latent semantic indexing (LSI) is applied as a preprocessing step, and the resulting continuous features are modeled using a mean squared error (MSE) reconstruction loss. The scRNA-seq reconstruction loss is the negative log-likelihood:
}\begin{displaymath}
	\DIFdel{\mathcal{L}_{\text{NB}}(z') = -\sum_{i,j} \log \text{NB}(x_{ij} | \hat{\mu}_{ij}, \hat{\theta}_j)
}\end{displaymath}%DIFAUXCMD
\DIFdel{where the mean $\hat{\mu}$ }\DIFdelend \DIFaddbegin \DIFadd{and dropout 0.05; reconstruction, inner-reconstruction }\DIFaddend and \DIFdelbegin \DIFdel{gene-wise dispersion $\hat{\theta}$ are outputs of the decoder, and }\DIFdelend \DIFaddbegin \DIFadd{Kullback--Leibler weights are all unity; training uses Adam at learning rate $10^{-4}$, subgraph size 512, ten subgraphs per epoch and 300 epochs; the cell--cell graph is a 15-nearest-neighbour Euclidean graph in a 50-dimensional PCA embedding. The full configuration is reproduced as Table~\ref{table_hyperparams} in the Appendix. The same configuration is applied uniformly across the benchmark cohort and across the fourteen evaluated encoder backbones. All randomness is controlled by a fixed seed propagated through the model, the optimiser and }\DIFaddend the \DIFdelbegin \DIFdel{mean is scaled by a cell-specific library size factor $l_i$: $\hat{\mu}_{ij} = l_i \cdot \text{softmax}(\text{Decoder}(z'_{i}))_j$. For scATAC-seq, the reconstruction loss is:
}\begin{displaymath}
	\DIFdel{\mathcal{L}_{\text{MSE}}(z') = \sum_{i,j} (x_{ij} - \hat{x}_{ij})^2
}\end{displaymath}%DIFAUXCMD
\DIFdel{where $\hat{x}_{ij}$ is the decoder output for cell $i$ and feature $j$.
}%DIFDELCMD <

%DIFDELCMD < %%%
\subsubsection{\DIFdel{Training Objective}}
%DIFAUXCMD
\addtocounter{subsubsection}{-1}%DIFAUXCMD
%DIFDELCMD <

%DIFDELCMD < %%%
\DIFdel{The total loss for CCVGAE is a weighted sum of three components:
}%DIFDELCMD < \begin{enumerate}
\begin{enumerate}%DIFAUXCMD
%DIFDELCMD < 	\item \textbf{Primary Reconstruction Loss} %%%
\item%DIFAUXCMD
\DIFdel{($\mathcal{L}_{\text{recon}}$): The modality-appropriate reconstruction loss (NB for scRNA-seq; MSE for scATAC-seq) from the primary latent sample $z$.
	}%DIFDELCMD < \item \textbf{Bottleneck Reconstruction Loss} %%%
\item%DIFAUXCMD
\DIFdel{($\mathcal{L}_{\text{irecon}}$): The corresponding reconstruction loss from the bottlenecked representation $z_r$. }%DIFDELCMD < \item \textbf{KL Divergence} %%%
\item%DIFAUXCMD
\DIFdel{($\mathcal{L}_{\text{KL}}$): The KL divergence between the approximate posterior $q_\phi(z|x)$ and a standard Gaussian prior $p(z) = \mathcal{N}(0, I)$.
}
\end{enumerate}%DIFAUXCMD
%DIFDELCMD < \end{enumerate}
%DIFDELCMD <

%DIFDELCMD < %%%
\DIFdel{The unified objective is formulated as:
}\begin{displaymath}
	\DIFdel{\mathcal{L}_{\text{CCVGAE}} = w_{\text{recon}} \mathcal{L}_{\text{recon}} + w_{\text{irecon}} \mathcal{L}_{\text{irecon}} + w_{\text{KL}} \mathcal{L}_{\text{KL}}
}\end{displaymath}%DIFAUXCMD
\DIFdel{This loss is computed for each mini-batch of data. When the GAT encoder is used, the total loss is calculated as an expectation over the stream of sampled subgraphs $S$: $\mathcal{L}_{\text{total}} = \mathbb{E}_{S \sim G} [\mathcal{L}_{\text{CCVGAE}}(S)]$. For the linear encoder, it is simply averaged over mini-batches}\DIFdelend \DIFaddbegin \DIFadd{data loaders, so each figure can be regenerated from the same checkpoint. Models are trained on a single GPU. Empirical scalability and runtime profiling are not the focus of this work and are reported in supplementary material}\DIFaddend .

\DIFdelbegin \subsubsection{\DIFdel{Default Hyperparameter Settings}}
%DIFAUXCMD
\addtocounter{subsubsection}{-1}%DIFAUXCMD
%DIFDELCMD < \label{sec:hyperparameters}
%DIFDELCMD <

%DIFDELCMD < %%%
\DIFdel{Table~\ref{table_hyperparams} summarizes the default hyperparameters used throughout our experiments. These values were selected based on preliminary experiments on a held-out validation set and remained fixed across all 170 benchmark datasets.
}%DIFDELCMD <

%DIFDELCMD < %%%
%DIF <  OMITTED_table_FOR_DIFF_BUILD_STABILITY
%DIFDELCMD <

%DIFDELCMD < %%%
\DIFdel{The loss weights ($w_{\text{recon}}$, $w_{\text{irecon}}$, $w_{\text{KL}}$) are set to }\DIFdelend \DIFaddbegin \DIFadd{Clustering used Leiden community detection at resolution }\DIFaddend 1.0 \DIFdelbegin \DIFdel{by default, giving equal weight to all objective components, and performed well across our benchmarks. Tuning these weights may help in specific applications---for instance, reducing $w_{\text{KL}}$ allows more flexible latent structures \mbox{%DIFAUXCMD
\cite{higgins2017betavae}}\hskip0pt%DIFAUXCMD
}\DIFdelend \DIFaddbegin \DIFadd{with the modularity-vertex partition and a fixed integer seed (42) propagated through }\texttt{\DIFadd{scanpy.tl.leiden}}\DIFadd{. Sensitivity to clustering algorithm and seed is reported in the supplementary analyses}\DIFaddend .

\DIFdelbegin \subsubsection{\DIFdel{Rationale for Deterministic Inference}}
%DIFAUXCMD
\addtocounter{subsubsection}{-1}%DIFAUXCMD
%DIFDELCMD <

%DIFDELCMD < %%%
\DIFdel{As introduced in Section~2.2.1, the deterministic posterior mean $\mu_{\text{centroid}}$ serves as the final embedding for visualization, clustering, and trajectory inference. The high technical noise and sparsity of single-cell data can be amplified by sampling variance in stochastic VAE outputs.Posterior-mean inference has been adopted in some prior VAE-based models, but its combination with coupling regularization and graph attention encoding has not been systematically evaluated at this scale. }%DIFDELCMD <

%DIFDELCMD < \subsection{Biological Interpretability Workflow}
%DIFDELCMD < \label{sec:interpretability}
%DIFDELCMD <

%DIFDELCMD < %%%
\DIFdel{To provide biological interpretability of the learned latent space, we implement a post-hoc analysis pipeline that maps latent dimensions to gene programs and functional annotations:
}%DIFDELCMD <

%DIFDELCMD < \begin{enumerate}
%DIFDELCMD <     \item \textbf{Latent--gene correlation:} %%%
\DIFdel{For each latent dimension $k \in \{1, \ldots, d_z\}$, we compute Pearson correlation coefficients between the latent values $\mu_{\text{centroid},k}$ across all cells and the expression of each gene $g$:
    }\begin{displaymath}
        \DIFdel{r_{kg} = \text{corr}(\mu_{\text{centroid},k}, x_g)
    }\end{displaymath}%DIFAUXCMD
%DIFDELCMD <

%DIFDELCMD <     \item \textbf{Top gene selection:} %%%
\DIFdel{For each dimension, we select the top 100 genes with the highest absolute correlation($|r_{kg}|$) as the ``associated gene set.''
    }\DIFdelend \DIFaddbegin \subsection{Biological interpretability workflow}\label{sec:latent_criteria}
\DIFaddend

\DIFdelbegin %DIFDELCMD < \item \textbf{Functional enrichment:} %%%
\DIFdelend \DIFaddbegin \DIFadd{The three paired case studies in Sec.~\ref{sec:case_sd_gastric}--\ref{sec:case_ir_covid} apply a single interpretation procedure across all six datasets rather than tuning a per-case heuristic. After training the model on a held-out biological dataset, each latent coordinate $z_0,\ldots,z_{d_z-1}$ is annotated by Pearson correlation against the gene-expression matrix, computed gene-by-gene. A latent coordinate is retained for discussion only if it is the most strongly correlated coordinate for at least one gene whose maximum absolute correlation across all latent dimensions exceeds $0.25$; this rule prevents a pleiotropic gene from being assigned to multiple latents and keeps the latent--gene mapping injective at the labelling step. Each retained coordinate is labelled by the single gene with which it has the highest correlation, and the top-correlated genes that meet the threshold feed }\DIFaddend Gene Ontology Biological Process (GO-BP) enrichment \DIFdelbegin \DIFdel{analysis is performed on each associated gene set using a }\DIFdelend \DIFaddbegin \DIFadd{via the }\DIFaddend hypergeometric test with Benjamini--Hochberg correction\DIFdelbegin \DIFdel{for multiple testing}\DIFdelend . Enrichment is \DIFdelbegin \DIFdel{evaluated against the Mouse and Human GO Biological Process collections, with mouse datasets assessed against the mouse collection and human datasets against the human collection.
    }%DIFDELCMD <

%DIFDELCMD <     \item \textbf{Visualization:} %%%
\DIFdel{Enrichment results are visualized as dot plots showing $-\log_{10}(p_{\text{adjusted}})$ for significant terms.
}%DIFDELCMD < \end{enumerate}
%DIFDELCMD <

%DIFDELCMD < \textbf{Rationale and limitations:} %%%
\DIFdel{This correlation-based approach assumes that latent dimensions capture continuous variation in gene expression programs. Limitations: (1) correlations do not imply causation; (2) nonlinear relationships between latent dimensions and genes may be missed; (3) results depend on the quality of GO annotations. These analyses should be treated as hypothesis-generating, with critical findings validated through orthogonal experiments. }%DIFDELCMD <

%DIFDELCMD < \subsection{Clustering and Representation Metrics}
%DIFDELCMD < %%%
\DIFdelend \DIFaddbegin \DIFadd{reported as descriptive context for the figure narrative; it is not used as a selection gate, and a latent that does not return significant enrichment is still rendered if its top-gene mapping is biologically interpretable.
}\DIFaddend

\DIFdelbegin \DIFdel{For all clustering-based metrics (ARI, NMI, ASW, DAV, CAL), we applied $k$-means clustering to the learned latent representations, with the number of clusters $k$ set equal to the number of ground-truth }\DIFdelend \DIFaddbegin \DIFadd{Figure labelling follows two consistency conventions. First, latent axes are drawn as $z_0,z_1,\ldots,z_{d_z-1}$ in every panel so they cannot be confused with biological time labels such as D0 or d0. Second, paired panels are symmetric on both sides of the figure: when curated condition and }\DIFaddend cell-type labels \DIFdelbegin \DIFdel{in each dataset. To account for the stochasticity of $k$-means initialization, we performed 10 independent runs with different random seeds per dataset and report the result from the run with the lowest inertia (within-cluster sum of squares). Uniform Manifold Approximation and Projection (UMAP) and t-distributed Stochastic Neighbor Embedding (t-SNE) visualizations were computed with fixed random seeds (seed = 42) to ensure reproducibility. All downstream evaluation was performed on the deterministic posterior mean embedding $\mu_{\text{centroid}}$ for centroid-based models, and on a single stochastic sample $z$ for non-centroid baselines, consistent with the respective model designs.
}%DIFDELCMD <

%DIFDELCMD < \begin{description}
\begin{description}%DIFAUXCMD
%DIFDELCMD <     \item[Adjusted Rand Index (ARI)] %%%
\item[\DIFdel{Adjusted Rand Index (ARI)}]%DIFAUXCMD
\DIFdel{Measures the similarity between true and predicted clusterings, corrected for chance agreement. A score of 1 indicates perfect agreement.
    }%DIFDELCMD < \item[Normalized Mutual Information (NMI)] %%%
\item[\DIFdel{Normalized Mutual Information (NMI)}]%DIFAUXCMD
\DIFdel{Quantifies the mutual information between true and predicted labels , normalized to a }%DIFDELCMD < [%%%
\DIFdel{0, 1}%DIFDELCMD < ] %%%
\DIFdel{scale, where 1 signifies perfect correlation.
    }%DIFDELCMD < \item[Average Silhouette Width (ASW)] %%%
\item[\DIFdel{Average Silhouette Width (ASW)}]%DIFAUXCMD
\DIFdel{Assesses cluster quality by measuring how similar a cell is to its own cluster compared to neighboring clusters. Higher values indicate denser, better-separated clusters.
    }%DIFDELCMD < \item[Calinski-Harabasz Index (CAL)] %%%
\item[\DIFdel{Calinski-Harabasz Index (CAL)}]%DIFAUXCMD
\DIFdel{The ratio of between-cluster dispersion to within-cluster dispersion. Higher values correspond to better-defined clusters. }%DIFDELCMD < \item[Davies-Bouldin Index (DAV)] %%%
\item[\DIFdel{Davies-Bouldin Index (DAV)}]%DIFAUXCMD
\DIFdel{The average similarity between each cluster and its most similar neighbor. Lower values indicate better separation.
}
\end{description}%DIFAUXCMD
%DIFDELCMD < \end{description}
%DIFDELCMD <

%DIFDELCMD < \subsection{Structure and Embedding Quality Metrics}
%DIFDELCMD <

%DIFDELCMD < \begin{description}
\begin{description}%DIFAUXCMD
%DIFDELCMD <     \item[Distance Correlation (COR)] %%%
\item[\DIFdel{Distance Correlation (COR)}]%DIFAUXCMD
\DIFdel{Spearman rank correlation of pairwise distances between the high-dimensional and low-dimensional representations, assessing global structure preservation.
    }%DIFDELCMD < \item[Data Compactness (DC)] %%%
\item[\DIFdel{Data Compactness (DC)}]%DIFAUXCMD
\DIFdel{Measures the average intra-cluster tightness in the embedding space, with higher values indicating more compact cluster representations.
    }%DIFDELCMD < \item[Local Quality (QL)] %%%
\item[\DIFdel{Local Quality (QL)}]%DIFAUXCMD
\DIFdel{Measures the preservation of local neighborhood structure in the embedding.
    }%DIFDELCMD < \item[Global Quality (QG)] %%%
\item[\DIFdel{Global Quality (QG)}]%DIFAUXCMD
\DIFdel{Assesses the preservation of large-scale, global data structure in the embedding.
    }%DIFDELCMD < \item[Overall Quality (OV)] %%%
\item[\DIFdel{Overall Quality (OV)}]%DIFAUXCMD
\DIFdel{A composite score averaging COR, QL, and QG to provide a unified measure of embedding quality, calculated separately for UMAP and t-SNE.
}
\end{description}%DIFAUXCMD
%DIFDELCMD < \end{description}
%DIFDELCMD < %%%
\DIFdelend \DIFaddbegin \DIFadd{exist in the source data they are used directly; when they do not, the panel is explicitly labelled as ``inferred group'' or ``inferred cell state'' so that curated and inferred labels are not silently mixed. The retention rule, the correlation threshold, the labelling procedure and the GO-BP enrichment step are applied identically across all three paired case studies, so cross-case differences reflect biology rather than workflow.
}\DIFaddend

\DIFdelbegin %DIFDELCMD < \subsection{Intrinsic Manifold and Biological Metrics}
%DIFDELCMD < %%%
\DIFdelend \DIFaddbegin \section{Results}
\DIFaddend

\DIFdelbegin %DIFDELCMD < \begin{description}
\begin{description}%DIFAUXCMD
%DIFDELCMD <     \item[Manifold Dimensionality (MD)] %%%
\item[\DIFdel{Manifold Dimensionality (MD)}]%DIFAUXCMD
\DIFdel{A composite score evaluating how efficiently the latent space captures data variance.
    }%DIFDELCMD < \item[Spectral Decay Rate (SDR)] %%%
\item[\DIFdel{Spectral Decay Rate (SDR)}]%DIFAUXCMD
\DIFdel{Measures the rate of eigenvalue decay of the data covariance matrix, indicating information concentration.
    }%DIFDELCMD < \item[Participation Ratio (PR)] %%%
\item[\DIFdel{Participation Ratio (PR)}]%DIFAUXCMD
\DIFdel{Quantifies the effective number of dimensions contributing to the representation, based on eigenvalue distribution. }%DIFDELCMD < \item[Anisotropy Score (AS)] %%%
\item[\DIFdel{Anisotropy Score (AS)}]%DIFAUXCMD
\DIFdel{Measures directional bias in the latent space, indicating how uniformly the space is utilized.}%DIFDELCMD < \item[Trajectory Directionality (TD)] %%%
\item[\DIFdel{Trajectory Directionality (TD)}]%DIFAUXCMD
\DIFdel{Evaluates the dominance of the primary developmental axis in trajectory data, reflecting how well the representation captures a biological process.
    }%DIFDELCMD < \item[Noise Resilience (NR)] %%%
\item[\DIFdel{Noise Resilience (NR)}]%DIFAUXCMD
\DIFdel{Quantifies the robustness of the learned representation to input noise perturbations, measuring degradation in cluster quality under added noise.
    }%DIFDELCMD < \item[Core Quality] %%%
\item[\DIFdel{Core Quality}]%DIFAUXCMD
\DIFdel{A composite score defined as the arithmetic mean of MD, SDR, PR, and AS, summarizing intrinsic manifold characteristics.
    }%DIFDELCMD < \item[Overall Intrinsic] %%%
\item[\DIFdel{Overall Intrinsic}]%DIFAUXCMD
\DIFdel{A composite score aggregating Core Quality, TD, and NR to provide a unified assessment of latent-space quality.
}
\end{description}%DIFAUXCMD
%DIFDELCMD < \end{description}
%DIFDELCMD < %%%
\DIFdelend \DIFaddbegin \subsection{Component-wise ablation against a stochastic VAE}\label{sec:ablation}
\DIFaddend

\DIFdelbegin %DIFDELCMD < \subsection{Statistical Analysis}
%DIFDELCMD < %%%
\DIFdelend \DIFaddbegin \DIFadd{We first asked which of the three architectural mechanisms contributes to which metric family. Fig.~\ref{fig03} compares each mechanism, in isolation, against a stochastic Variational Autoencoder baseline trained under the same loss schedule, on the same kNN-Euclidean graph and from the same 50-dimensional PCA preprocessing.
}\DIFaddend

\DIFdelbegin \DIFdel{To compare model performance, we conducted statistical tests across 55 scRNA-seq }\DIFdelend \DIFaddbegin \DIFadd{Block~}\textbf{A} \DIFadd{(CenVAE vs.\ VAE) replaces the reparameterised sample with the deterministic posterior mean and leaves every other component unchanged. Centroid inference moves clustering compactness the most: ASW improves by $+0.114$ (Holm-corrected $p<0.001$), CAL by $+615.6$, }\DIFaddend and \DIFdelbegin \DIFdel{115 scATAC-seq datasets.For all comparisons, the normality of metric differences was first assessed using the Shapiro-Wilk test. }%DIFDELCMD < \begin{itemize}
\begin{itemize}%DIFAUXCMD
%DIFDELCMD <     \item \textbf{Paired Comparisons}%%%
\item%DIFAUXCMD
\DIFdel{: For ablation studies involving two models, a two-sided paired t-test was used for normally distributed differences, while the non-parametric two-sided Wilcoxon signed-rank test was used otherwise.
    }%DIFDELCMD < \item \textbf{Multiple Comparisons}%%%
\item%DIFAUXCMD
\DIFdel{: For benchmarking against multiple state-of-the-art methods, we first tested for overall differences. A Repeated Measures ANOVA (RMANOVA)was used for normally distributed data, and a Friedman test was applied to non-normally distributed data. If a significant overall difference was found, post-hoc pairwise tests (t-test or Wilcoxon) were performed with Bonferroni correction to identify specific differences between our model and baselines. }
\end{itemize}%DIFAUXCMD
%DIFDELCMD < \end{itemize}
%DIFDELCMD < %%%
\DIFdel{Significance levels are denoted as: ns ($p \geq 0.05$), * ($p < 0.05$), ** ($p < 0.01$}\DIFdelend \DIFaddbegin \DIFadd{DAV by $-0.426$. Improvements on neighbourhood-preserving embedding are smaller (UMAP overall $+0.049$, t-SNE overall $+0.054$}\DIFaddend ), and \DIFdelbegin \DIFdel{*** ($p < 0.001$). Each box represents the interquartile range (IQR) from the 25th (Q1) to the 75th (Q3) percentile, with a central line indicating the median. The whiskers extend to 1.5 times the IQR beyond the box, and points outside this range are plotted individually as outliers}\DIFdelend \DIFaddbegin \DIFadd{the intrinsic-overall score gains $+0.088$. Eliminating sampling variance at inference is therefore enough to deliver a more clusterable embedding without altering training}\DIFaddend .

\DIFdelbegin %DIFDELCMD < \section{Results}
%DIFDELCMD <

%DIFDELCMD < \subsection{Centroid Inference Stabilizes Latent Geometry and Improves Clustering}
%DIFDELCMD <

%DIFDELCMD < %%%
\DIFdel{We evaluated posterior centroid anchoring (CenVAE) against a standard stochastic VAE across 55 scRNA-seq }\DIFdelend \DIFaddbegin \DIFadd{Block~}\textbf{B} \DIFadd{(CouVAE vs.\ VAE) retains the stochastic sample at inference and adds the bottleneck-coupled inner-reconstruction loss during training. Coupling regularisation gives a modest gain on clustering compactness (ASW $+0.056$, DAV $-0.264$) but the largest single-mechanism gain on neighbourhood-preserving embedding (UMAP overall $+0.084$) }\DIFaddend and \DIFdelbegin \DIFdel{115 scATAC-seq datasets. Training trajectories in Fig.~\ref{fig03}A and Fig. ~\ref{fig03}C show smoother, more consistent curves for CenVAE in average silhouette width (ASW ), Calinski-Harabasz index (CAL) , Davies-Bouldin index (DAV) , and correlation score (COR). End-of-training differences are quantified below.
}%DIFDELCMD <

%DIFDELCMD < %%%
\DIFdel{Clustering and representation metrics consistently favor CenVAE (Fig.~\ref{fig03}B, D; Table~\ref{table1}) . For scRNA-seq, mean differences are NMI \(+0.044\) ($p < 0.001$, Wilcoxon signed-rank test) , ARI \(+0.044\) ($p < 0.001$), ASW \(+0.114\) ($p < 0.001$), DAV advantage \(+0.439\) ($p < 0.001$), CAL \(+565.865\) ($p < 0.001$) , and COR \(+0.676\) ($p < 0.001$). For scATAC-seq, the corresponding values are NMI \(+0.075\) ($p < 0.001$), ARI \(+0.075\) ($p < 0.001$), ASW \(+0.138\) ($p < 0.001$), DAV advantage \(+0.883\) ($p < 0.001$), CAL \(+1441.383\) ($p < 0.001$), and COR \(+2.523\) ($p < 0.001$). All reported $p$-values are from two-sided Wilcoxon signed-rank tests with Bonferroni correction where applicable.
}%DIFDELCMD <

%DIFDELCMD < %%%
\DIFdel{Embedding quality improves across UMAP and t-SNE.
For scRNA-seq, overall embedding scores (OV) increase by \(+0.040\) (UMAP, $p < 0.01$) and \(+0.047\) (t-SNE, $p < 0.01$); for scATAC-seq, OV increases by \(+0.221\) (UMAP, $p < 0.001$) and \(+0.235\) (t-SNE, $p < 0.001$).Local continuity (QL) and global structure (QG) show parallel gains (Table~\ref{table1})}\DIFdelend \DIFaddbegin \DIFadd{on intrinsic geometry (intrinsic-overall $+0.128$; Holm-corrected $p<0.001$ in both cases). Coupling alone reshapes the latent geometry that 2D projections see, even though it does not by itself markedly improve clusterability}\DIFaddend .

\DIFdelbegin \DIFdel{Intrinsic geometry also benefits. For scRNA-seq, statistically significant improvements ($p < 0.01$) include manifold dimensionality \(+0.108\), participation ratio \(+0.131\), anisotropy \(+0.152\), core quality \(+0.112\), }\DIFdelend \DIFaddbegin \DIFadd{Block~}\textbf{C} \DIFadd{(GAT-VAE vs.\ VAE) replaces the dense encoder with the graph-attention encoder on the kNN-Euclidean graph. Graph attention contributes the smallest single-mechanism move on intrinsic geometry (intrinsic-overall $+0.002$, essentially unchanged) but still improves clustering compactness (ASW $+0.070$, DAV $-0.229$, CAL $+237.6$) }\DIFaddend and \DIFdelbegin \DIFdel{overall intrinsic quality \(+0.077\). For scATAC-seq, all gains reach $p < 0.001$: \(+0.331\), \(+0.507\), \(+0.343\), \(+0.345\), and \(+0.384\), respectively (Fig.~\ref{fig03}B, D; Table~\ref{table1}). Biological metrics (trajectory directionality and noise resilience) also trend upward, particularly on scATAC-seq, where data sparsity is higher. }\DIFdelend \DIFaddbegin \DIFadd{neighbourhood-preserving embedding (UMAP overall $+0.029$). The small intrinsic-geometry gain here is consistent with the joint analysis below: the largest intrinsic-overall improvement only emerges when graph attention is combined with centroid inference and coupling regularisation.
}\DIFaddend

\DIFdelbegin \DIFdel{CenVAE yields consistent mean improvements in clustering, embeddings, and latent-space geometry across the 170 datasets, with most comparisons reaching $p < 0.001$}\DIFdelend \DIFaddbegin \DIFadd{Across the three blocks, each mechanism moves a different metric family, so adopting one mechanism in isolation does not reproduce the joint effect quantified in Fig.~\ref{fig04}}\DIFaddend .


% OMITTED_figure_FOR_DIFF_BUILD_STABILITY


%DIF <  OMITTED_tableSTAR_FOR_DIFF_BUILD_STABILITY
\DIFaddbegin \subsection{Stability of the centroid embedding}\label{sec:stability}
\DIFadd{The deterministic centroid posterior mean used at inference produces a stable representation across random seeds. To validate this claim empirically, we trained the full configuration with five seeds on three representative datasets and recorded the across-seed embedding variance and the ARI of Leiden clusters against ground-truth labels. The centroid configuration produces tight clusters that vary little across runs, while a matched stochastic-sampling variant under the same seed schedule shows visibly larger embedding variance. Numerical summaries are reported in the supplementary analyses.
}\DIFaddend

\DIFdelbegin %DIFDELCMD < \subsection{Coupling Regularization Enhances Latent Compactness and Noise Resilience}
%DIFDELCMD < %%%
\DIFdelend \DIFaddbegin \subsection{Joint configuration against single-component ablations}\label{sec:flagship}
\DIFaddend

We next \DIFdelbegin \DIFdel{assessed the independent contribution of the coupling mechanism by comparing CouVAE with the standard VAE. As shown in Fig.~\ref{fig04}A and Fig.~\ref{fig04}C, training trajectories for ASW, CAL, DAV, and COR favor CouVAE on both modalities. Boxplots (}\DIFdelend \DIFaddbegin \DIFadd{combined the three mechanisms and asked whether the joint configuration is more than the sum of its parts. }\DIFaddend Fig.~\ref{fig04} \DIFdelbegin \DIFdel{B, D) and Table~\ref{table1} summarize end-of-training differences.
}\DIFdelend \DIFaddbegin \DIFadd{compares the full scCCVGBen configuration against the four single-component variants (VAE, CenVAE, CouVAE, GAT-VAE), arranged left-to-right (Base $\rightarrow$ Centroid $\rightarrow$ Coupling $\rightarrow$ GAT $\rightarrow$ scCCVGBen) with significance brackets fanning from the rightmost scCCVGBen column. The VAE column is pooled across the three component-wise comparisons to provide a single cross-comparison estimate.
}\DIFaddend

\DIFdelbegin \DIFdel{Global clustering scores remain essentially unchanged for scRNA-seq (NMI $-0.003$, ARI $-0.003$; $p \geq 0.05$, Wilcoxon signed-rank test) and improve modestly for scATAC-seq (NMI $+0.015$, ARI $+0.015$; $p < 0.001$). Representation metrics improve significantly: ASW $+0.056$ (scRNA, $p < 0.001$) and $+0.018$ (scATAC, $p < 0.001$); DAV advantage $+0.282$ ($p < 0.001$) and $+0.166$ ($p < 0.001$); CAL $+582.555$ ($p < 0.001$) and $+159.291$ ($p < 0.001$); COR $+1.129$ ($p < 0.001$) and $+0.887$ ($p < 0.001$). Embedding metrics show DC gains of $+0.134$ (UMAP) and $+0.127$ (t-SNE) for scRNA-seq, and $+0.149$ (UMAP) and $+0.131$ (}\DIFdelend \DIFaddbegin \DIFadd{Against the VAE baseline, scCCVGBen improves ASW by $+0.288$ (Holm-corrected $p<0.001$), CAL by $+4270.9$, DAV by $-0.800$, intrinsic-overall by $+0.233$, UMAP overall by $+0.078$ and }\DIFaddend t-SNE \DIFdelbegin \DIFdel{) for scATAC-seq; OV increases fall between $+0.089$ and $+0.093$ across methods and modalities (all $p < 0.001$). Intrinsic metrics increase on both modalities, e.g., participation ratio $+0.224$ (scRNA) and $+0.232$ (scATAC); core quality $+0.148$ and $+0.126$; overall intrinsic $+0.129$ and $+0.110$ (all $p < 0.001$, Wilcoxon signed-rank test) . See Fig.~\ref{fig04}B, D for significance annotations. }%DIFDELCMD <

%DIFDELCMD < %%%
%DIF <  OMITTED_figure_FOR_DIFF_BUILD_STABILITY
%DIFDELCMD <

%DIFDELCMD < %%%
%DIF <  Table 2 merged into Table 1 (component-wise ablation)
%DIFDELCMD <

%DIFDELCMD < \subsection{Graph-Attention Encoding Preserves Cell--Cell Neighborhood Structure}
%DIFDELCMD <

%DIFDELCMD < %%%
\DIFdel{We next isolated the third component by comparing a graph attention encoder (GAT-VAE) with the standard VAE. Training dynamics (Fig.~\ref{fig05}A, C) show improved ASW, CAL, DAV, and COR, particularly for scATAC-seq. End-of-training differences are summarized in Table~\ref{table1}}\DIFdelend \DIFaddbegin \DIFadd{overall by $+0.091$ (each $p<0.001$)}\DIFaddend . \DIFdelbegin %DIFDELCMD <

%DIFDELCMD < %%%
\DIFdel{For scRNA-seq, GAT-VAE improves clustering (NMI $+0.087$, ARI $+0.088$; $p < 0.001$, Wilcoxon signed-rank test), representation (ASW $+0.059$, $p < 0.001$; DAV advantage $+0.166$, $p < 0.001$; CAL $+181.405$, $p < 0.01$), and embedding metrics (e.g., QL $+0.048$ for UMAP }\DIFdelend \DIFaddbegin \DIFadd{Against each single-component variant, gains in clustering compactness and intrinsic geometry remain substantial: ASW improves by $+0.163$ over CenVAE, $+0.220$ over CouVAE and $+0.203$ over GAT-VAE, }\DIFaddend and \DIFdelbegin \DIFdel{$+0.050$ }\DIFdelend \DIFaddbegin \DIFadd{intrinsic-overall improves by $+0.172$, $+0.128$ and $+0.254$ in the same comparisons. The gain over CouVAE on neighbourhood-preserving embedding is smaller and not significant under Holm correction (UMAP overall $+0.004$, $p\approx 0.47$; }\DIFaddend t-SNE \DIFaddbegin \DIFadd{overall $+0.013$, $p\approx 0.06$), echoing the ablation reading that coupling regularisation alone already accounts for most of the projection-quality gain. Clustering compactness and intrinsic geometry}\DIFaddend , \DIFdelbegin \DIFdel{$p < 0.01$; OV $+0.022$ UMAP, $+0.020$ t-SNE). Some intrinsic metrics show small decreases (manifold dimensionality $-0.016$, participation ratio $-0.018$, core quality $-0.013$; all $p \geq 0.05$); noise resilience change is $-0.000$. For scATAC-seq, gains are broader and statistically significant: NMI $+0.027$ ($p < 0.001$), ARI $+0.028$ ($p < 0.001$), ASW $+0.012$ ($p < 0.001$), DAV advantage $+0.079$ ($p < 0.001$), CAL $+74.346$ ($p < 0.001$), COR $+0.162$ ($p < 0.01$), and embedding DC (UMAP ) $+0.092$ ($p < 0.001$). Intrinsic geometry improves (participation ratio $+0.100$, core quality $+0.052$; $p < 0.001$), as do biological metrics (trajectory directionality $+0.051$, noise resilience $+0.023$; $p < 0.001$).
}%DIFDELCMD <

%DIFDELCMD < %%%
\DIFdel{A modest clustering-versus-geometry trade-off is apparent for scRNA-seq: NMI and ARI increase substantially, whereas some intrinsic geometrymetrics decrease slightly (manifold dimensionality $-0.016$, core quality $-0.013$; both $p \geq 0.05$). Graph-attention encoding prioritizes local neighborhood fidelity at a minor cost to global spectral structure. Combining all three components in CCVGAE (Section~\ref{sec:ccvgae_vs_cenvae}) resolves the trade-off, with geometry metrics recovering and exceeding the baseline. This tension is revisited in the Discussion}\DIFdelend \DIFaddbegin \DIFadd{by contrast, require all three mechanisms together}\DIFaddend .


% OMITTED_figure_FOR_DIFF_BUILD_STABILITY


%DIF <  Table 3 merged into Table 1 (component-wise ablation)
\DIFdelbegin %DIFDELCMD <

%DIFDELCMD < \subsection{Combining Centroid Inference with Graph Attention Yields Complementary Gains}\label{sec:ccvgae_vs_cenvae}
%DIFDELCMD < %%%
\DIFdelend \DIFaddbegin \subsection{Encoder backbone robustness}\label{sec:encoder_rob}
\DIFaddend

We \DIFdelbegin \DIFdel{evaluated the combined effect of centroid anchoring and graph-based encoding in CCVGAE against the centroid-only baseline (CenVAE). Training curves (Fig. ~\ref{fig06}A, C) are stable for both methods; end-of-training differences (}\DIFdelend \DIFaddbegin \DIFadd{then asked whether the advantage of the default configuration depends on the specific Graph Attention Network used as the encoder. }\DIFaddend Fig.~\DIFdelbegin \DIFdel{\ref{fig06}B, D; Table~\ref{table4}) quantify the effect of adding graph structure}\DIFdelend \DIFaddbegin \DIFadd{\ref{fig05} swaps $f_\theta$ across fourteen encoder backbones at fixed core and fixed graph, treating the result as a robustness audit rather than a post-hoc superiority test}\DIFaddend .

\DIFdelbegin \DIFdel{Average clustering scores decrease slightly but not significantly: for scRNA-seq, NMI $-0.014$ and ARI $-0.014$ (Wilcoxon signed-rank test, both $p \geq 0.05$); for scATAC-seq, NMI $-0.017$ and ARI $-0.017$ (paired t-test, both $p \geq 0.05$). In contrast, most representation and geometry metrics improve: ASW $+0.018$ (scRNA, $p < 0.01$) and $+0.004$ (scATAC, $p \geq 0.05$); DAV advantage $+0.052$ ($p < 0.001$) and }\DIFdelend \DIFaddbegin \DIFadd{Within the attention family GAT and GATv2 are essentially indistinguishable on every metric (ASW $-0.001$, DAV $-0.011$, CAL $-125.0$, intrinsic-overall $+0.003$); GATv2 can therefore be substituted for the default without a measurable benchmark cost. Across families, GAT outperforms convolutional and propagation backbones on clustering compactness, with ASW differences of $+0.226$ versus GraphConv, $+0.208$ versus ARMA, $+0.194$ versus Cheb, $+0.066$ versus GCN, $+0.079$ versus SAGE, }\DIFaddend $+0.040$ \DIFdelbegin \DIFdel{($p < 0.05$); CAL $+980.294$ ($p < 0.001$) and $+448.487$ ($p < 0.001$); COR $+0.967$ ($p < 0.001$) and $+0.662$ ($p < 0.001$). Embedding OV increases by $+0.035$ (UMAP, $p < 0.001$) and $+0.026$ (t-SNE, $p < 0.001$) for scRNA-seq, }\DIFdelend \DIFaddbegin \DIFadd{versus SG and $+0.051$ versus SSG, }\DIFaddend and \DIFaddbegin \DIFadd{CAL differences such as $+4130.4$ versus GraphConv. Intrinsic geometry tells a different story: GIN improves intrinsic-overall }\DIFaddend by \DIFdelbegin \DIFdel{$+0.023$ (UMAP, $p < 0.001$) and $+0.018$ (t-SNE, $p < 0.001$) for scATAC-seq. Intrinsic metrics also rise: manifold dimensionality $+0.139$ and $+0.042$, participation ratio $+0.142$ and $+0.040$, anisotropy $+0.089$ and $+0.042$, core quality $+0.108$ and $+0.037$, and overall intrinsic $+0.101$ and $+0.061$ (all $p < 0.001$)}\DIFdelend \DIFaddbegin \DIFadd{$0.235$ relative to GAT and EdgeConv improves it by $0.194$, even though both lose to GAT on clustering compactness. GAT is therefore strong on clustering compactness and competitive on neighbourhood preservation; message-passing and edge-convolution backbones can outperform it specifically on intrinsic geometry}\DIFaddend .


% OMITTED_figure_FOR_DIFF_BUILD_STABILITY


%DIF <  OMITTED_tableSTAR_FOR_DIFF_BUILD_STABILITY
\DIFdelbegin %DIFDELCMD <

%DIFDELCMD < \subsection{Full Model Integration Supports Component Synergy over Coupling-Only Baselines}
%DIFDELCMD < %%%
\DIFdelend \DIFaddbegin \subsection{Graph-construction robustness}\label{sec:graph_rob}
\DIFaddend

\DIFdelbegin \DIFdel{To complete the ablation hierarchy, we compared the full CCVGAE model with the coupling-only baseline (CouVAE). Training dynamics (}\DIFdelend \DIFaddbegin \DIFadd{We similarly asked whether the advantage depends on the choice of kNN-Euclidean similarity. }\DIFaddend Fig.~\DIFdelbegin \DIFdel{\ref{fig07}A, C) favor CCVGAE; end-of-training differences (Fig.~\ref{fig07}B, D; Table~\ref{table4}) show broad improvements.
}\DIFdelend \DIFaddbegin \DIFadd{\ref{fig06} holds the algorithmic core (centroid inference, coupling regularisation and graph attention) fixed and varies the cell--cell graph between kNN-Euclidean, kNN-cosine, SNN, mutual-kNN and Gaussian-threshold.
}\DIFaddend

\DIFdelbegin \DIFdel{Clustering gains are positive across modalities: for scRNA-seq, NMI and ARI are both $+0.104$ (Wilcoxon signed-rank test, $p < 0.001$); for scATAC-seq, NMI $+0.169$ and ARI $+0.170$ ($p < 0.001$).
Representation gains include ASW $+0.210$ and $+0.126$ ($p < 0.001$), DAV advantage $+0.486$ and $+0.679$ ($p < 0.001$), CAL $+1708.890$ and $+1773.136$ ($p < 0.001$), and COR $+0.286$ ($p < 0.01$) and $+2.012$ ($p < 0.001$) (scRNA, scATAC). Embedding OV is mixed on scRNA-seq, with $-0.003$ for UMAP ($p \geq 0.05$) but $+0.012$ for t-SNE ($p < 0.05$), whereas scATAC-seq improves by $+0.123$ (UMAP) and $+0.125$ (t-SNE) (both $p < 0.001$). Intrinsic metrics also improve: manifold dimensionality $+0.027$ and $+0.229$, participation ratio $+0.026$ and $+0.312$, anisotropy $+0.094$ and $+0.212$, core quality $+0.049$ and $+0.225$, and overall intrinsic $+0.040$ and $+0.286$ (all $p < 0.01$ or better)}\DIFdelend \DIFaddbegin \DIFadd{The choice of graph construction is not inert. Compared with the kNN-Euclidean reference, SNN improves ASW by $0.094$ and CAL by $2863.5$, and kNN-cosine improves ASW by $0.050$ and CAL by $1516.7$. Conversely, kNN-Euclidean is better than mutual-kNN (ASW $+0.042$, CAL $+1410.0$) and modestly better than Gaussian-threshold (ASW $+0.013$, CAL $+1601.2$, intrinsic-overall $+0.045$). Differences on the neighbourhood-preserving suite are smaller and often non-significant in either direction, so the dependence on graph construction is concentrated in clustering compactness. The practical reading is conservative: scCCVGBen does not collapse under graph-construction perturbation, but the choice of cell--cell similarity is a real modelling decision and should not be conflated with the algorithmic-core claim}\DIFaddend .


% OMITTED_figure_FOR_DIFF_BUILD_STABILITY


%DIF <  Table 5 merged into Table 4 (full-model integration)
\DIFdelbegin %DIFDELCMD <

%DIFDELCMD < \subsection{CCVGAE Achieves Favorable Performance on scRNA-seq Benchmarks}
%DIFDELCMD < %%%
\DIFdelend \DIFaddbegin \subsection{Cross-method benchmark on scRNA-seq}\label{sec:scrna_bench}
\DIFaddend

We \DIFdelbegin \DIFdel{benchmarked CCVGAE on }\DIFdelend \DIFaddbegin \DIFadd{then compared scCCVGBen to twelve external }\DIFaddend scRNA-seq \DIFdelbegin \DIFdel{against classicalmethods }\DIFdelend \DIFaddbegin \DIFadd{baselines spanning a classical/linear family }\DIFaddend (PCA, \DIFdelbegin \DIFdel{Kernel PCA (KPCA), Independent Component Analysis (ICA), Factor Analysis (FA), Non-negative Matrix Factorization (NMF), Truncated SVD (TSVD), Dictionary Learning (DICL) ) and deep generative models }\DIFdelend \DIFaddbegin \DIFadd{KPCA, ICA, FA, NMF, TSVD, DICL) and a deep generative family }\DIFaddend (scVI, DIPVAE, InfoVAE, \DIFdelbegin \DIFdel{TCVAE}\DIFdelend \DIFaddbegin \DIFadd{$\beta$-TCVAE}\DIFaddend , HighBetaVAE). \DIFaddbegin \DIFadd{Methods in }\DIFaddend Fig.~\DIFdelbegin \DIFdel{\ref{fig08}A--E show consistent advantages across clustering/representation, UMAP/t-SNE embedding quality , intrinsic geometry, and aggregated scores, with many tests at \(p < 0.001\). Table~\ref{table6} reports mean differences (CCVGAE \(-\) baseline)}\DIFdelend \DIFaddbegin \DIFadd{\ref{fig07} are ordered left-to-right by aggregate mean rank across the directional quality metrics, so the weakest comparator appears leftmost and scCCVGBen is the rightmost reference. The number of paired datasets per comparison ranges from 46 to 94 in the linear family because some baselines are not defined for every accession in the cohort, and equals 48 for the deep generative family}\DIFaddend .

Against \DIFdelbegin \DIFdel{scVI, CCVGAE improves ASW by $+0.329$ ($p < 0.001$, Friedman test with Bonferroni-corrected post-hoc), DAV advantage by $+1.060$ ($p < 0.001$), CAL by $+2392.504$ ($p < 0.001$), COR by $+2.523$ ($p < 0.001$), OV (UMAP) by $+0.219$ ($p < 0.001$), OV (t-SNE) by $+0.229$ ($p < 0.001$), participation ratio by $+0.560$ ($p < 0.001$), core quality by $+0.351$ ($p < 0.001$), and overall intrinsic by $+0.272$ ($p < 0.001$). A modest anisotropy trade-off is observed only against DIPVAE ($-0.040$); against all other baselines including scVI ($+0.390$), anisotropy improves. Similar or larger advantages are observed relative to PCA/KPCA/ICA/FA /NMF/TSVD}\DIFdelend \DIFaddbegin \DIFadd{the classical}\DIFaddend /\DIFdelbegin \DIFdel{DICL and other VAEs (Table~\ref{table6}).
}%DIFDELCMD <

%DIFDELCMD < %%%
%DIF <  OMITTED_figure_FOR_DIFF_BUILD_STABILITY
%DIFDELCMD <

%DIFDELCMD < %%%
%DIF <  OMITTED_tableSTAR_FOR_DIFF_BUILD_STABILITY
%DIFDELCMD <

%DIFDELCMD < \subsection{CCVGAE Achieves Competitive Performance on scATAC-seq Benchmarks}
%DIFDELCMD <

%DIFDELCMD < %%%
\DIFdel{We benchmarked CCVGAE on scATAC-seq against LSI, PeakVI, and PoissonVI (Fig.~\ref{fig09}A--E). Table~\ref{table7} reports mean differences (CCVGAE \(-\) baseline).
}%DIFDELCMD <

%DIFDELCMD < %%%
\DIFdel{Relative to LSI and PeakVI, CCVGAE improves ASW by $+0.121$ and $+0.140$ ($p < 0.001$, Friedman test with Bonferroni-corrected post-hoc), DAV advantage by $+0.538$ and $+0.662$ ($p < 0.001$), CAL by $+2473.908$ and $+2635.712$ ($p < 0.001$), COR by $+3.949$ }\DIFdelend \DIFaddbegin \DIFadd{linear family, scCCVGBen improves clustering compactness with ASW differences of $+0.162$, $+0.164$, $+0.164$, $+0.114$, $+0.208$, $+0.070$ }\DIFaddend and \DIFdelbegin \DIFdel{$+2.269$ ($p < 0.001$), }\DIFdelend \DIFaddbegin \DIFadd{$+0.024$ versus PCA, KPCA, TSVD, NMF, ICA, FA }\DIFaddend and \DIFdelbegin \DIFdel{core quality by $+0.239$ }\DIFdelend \DIFaddbegin \DIFadd{DICL respectively, and improves intrinsic-overall by $+0.183$, $+0.184$, $+0.179$, $+0.268$, $+0.384$, $+0.433$ and $+0.267$ in the same order. Against the deep generative family, the largest ASW improvement on the entire panel is recorded against scVI ($+0.341$, with intrinsic-overall $+0.331$ in the same comparison); the remaining ASW gaps are $+0.266$, $+0.224$, $+0.226$ }\DIFaddend and \DIFdelbegin \DIFdel{$+0.273$ ($p < 0.001$). Compared with PoissonVI, CCVGAE increases CAL by $+485.661$ ($p < 0.001$), COR by $+1.050$ ($p < 0.001$)}\DIFdelend \DIFaddbegin \DIFadd{$+0.175$ versus DIPVAE, InfoVAE, $\beta$-TCVAE and HighBetaVAE, with intrinsic-overall gaps of $+0.278$}\DIFaddend , \DIFdelbegin \DIFdel{OV (UMAP) by $+0.009$ ($p \geq 0.05$), participation ratio by $+0.090$ ($p < 0.001$), }\DIFdelend \DIFaddbegin \DIFadd{$+0.199$, $+0.180$ and $+0.146$ in the same comparisons.
}

\DIFadd{The neighbourhood-preserving rows are tighter than the clustering and intrinsic-geometry rows: scVI and several deep generative baselines remain competitive on selected projection endpoints. scCCVGBen therefore has the highest aggregate rank among the scRNA-seq references at the present data cut, with margins that vary by metric family rather than a uniform win on every column.
}

\DIFadd{The slight decline in NMI and ARI reported for some methods or settings reflects an interaction between cluster resolution and embedding compactness rather than a property of the proposed model. When clusters are very compact, the Leiden modularity objective prefers a smaller number of merged communities than the ground-truth labelling, which slightly depresses NMI and ARI without changing downstream interpretability; sensitivity to clustering algorithm and seed is discussed in the supplementary analyses (D5 }\DIFaddend and \DIFdelbegin \DIFdel{overall intrinsic by $+0.103$ ($p < 0.001$), with trade-offs in ASW $-0.029$ ($p < 0.001$) and anisotropy $-0.059$ ($p < 0.001$).
}\DIFdelend \DIFaddbegin \DIFadd{D6).
}\DIFaddend


% OMITTED_figure_FOR_DIFF_BUILD_STABILITY


%DIF <  OMITTED_tableSTAR_FOR_DIFF_BUILD_STABILITY
\DIFdelbegin %DIFDELCMD <

%DIFDELCMD < \subsection{Practical Significance of Observed Improvements}
%DIFDELCMD < %%%
\DIFdelend \DIFaddbegin \subsection{Cross-method benchmark on scATAC-seq}\label{sec:scatac_bench}
\DIFaddend

\DIFdelbegin \DIFdel{Many improvements in Tables~\ref{table6} and~\ref{table7} are statistically significant; their practical magnitudes also merit attention.
The $+0.329$ ASW gain over scVI on scRNA-seq corresponds to a shift from poorly separated to moderately well-separated clusters in UMAP space---visible in two-dimensional projections and relevant for downstream cell-type annotation. The $+0.219$ overall UMAP score improvement reflects better preservation of local and global structure, relevant when embeddings are used to identify rare cell populations or infer developmental trajectories. On scATAC-seq, where sparsity makes representation learning harder, the $+0.273$ core quality gain over PeakVI indicates improved intrinsic manifold geometry: the latent space captures the continuous chromatin accessibility landscape rather than collapsing it into artifactual clusters. Intrinsic geometry gains (participation ratio, spectral decay, trajectory directionality) matter for trajectory inference, where distorted local neighborhoods can produce spurious branching events or incorrect pseudotime orderings. CCVGAE's primary practical contribution is }\textit{\DIFdel{consistency across metrics and datasets}}%DIFAUXCMD
\DIFdel{: the framework provides reliable gains across all three evaluation suites rather than excelling on a single benchmark}\DIFdelend \DIFaddbegin \DIFadd{On the chromatin-accessibility cohort, scCCVGBen was compared with three modality-specific baselines: LSI, PeakVI and PoissonVI (Fig.~\ref{fig08}). Despite the higher sparsity and lower per-feature signal-to-noise of binarised peak-by-cell matrices, scCCVGBen improves clustering compactness substantially against the older atlas-style baselines (ASW $+0.120$ vs.\ LSI, $+0.139$ vs.\ PeakVI; CAL $+2501.6$ vs.\ LSI, $+2663.5$ vs.\ PeakVI; DAV $-0.526$ vs.\ LSI, $-0.651$ vs.\ PeakVI)}\DIFaddend . \DIFaddbegin \DIFadd{The largest cross-suite improvement on the panel is on intrinsic geometry, with intrinsic-overall $+0.299$ versus LSI and $+0.346$ versus PeakVI.
}\DIFaddend

\DIFdelbegin %DIFDELCMD < \subsection{Latent Component Selection and Interpretation Criteria}\label{sec:latent_criteria}
%DIFDELCMD <

%DIFDELCMD < %%%
\DIFdel{Before examining individual case studies, we describe how latent components were retained and labeled for interpretation.All $d_z = 10$ latent dimensions are kept during training. For post-hoc interpretation, we compute the Pearson correlation $r_{jg}$ between each latent dimension $j$ (posterior mean across cells)and the expression of each gene $g$. A latent dimension is retained for discussion if it is the single most strongly correlated latent for at least one gene whose maximum absolute correlation across all latents exceeds $0.25$. Each retained latent is then labeled by the single gene with which it has the highest correlation. Gene Ontology Biological Process (GO-BP) enrichment is subsequently performed on the genes associated with the retained latents as a descriptive, interpretive step to place each component in a biological context; it is not used as a selection gate, and no Benjamini--Hochberg threshold or condition-differential activity test is applied to filter latents. This same correlation-based retention and labeling procedure is used uniformly across the three case studies that follow}\DIFdelend \DIFaddbegin \DIFadd{PoissonVI is the closest contender. scCCVGBen is essentially tied with PoissonVI on clustering compactness (ASW $-0.027$) but still improves intrinsic-overall by $+0.115$ in the same comparison. PoissonVI remains competitive on parts of the neighbourhood-preserving suite. scCCVGBen therefore outperforms LSI and PeakVI across all three metric families and matches or exceeds PoissonVI on clustering compactness and intrinsic geometry, while remaining comparable on neighbourhood preservation}\DIFaddend .


\DIFdelbegin %DIFDELCMD < \subsection{Latent Components Capture Sleep-Deprivation Responses in Bone Marrow}
%DIFDELCMD < %%%
\DIFdelend %DIF >  OMITTED_figure_FOR_DIFF_BUILD_STABILITY


\DIFdelbegin \DIFdel{To evaluate latent-dimension interpretability, we applied CCVGAE to scRNA-seq data from bone marrow cells of sleep-deprived mice. Prolonged sleep deprivation induces extensive transcriptomic remodeling in bone marrow immune cells and can precipitate cytokine storm-like responses in the periphery \mbox{%DIFAUXCMD
\cite{mcalpine2019sleep, cell2023prolonged, wang2024sleep}}\hskip0pt%DIFAUXCMD
. Following the workflow outlined in Section~\ref{sec:latent_criteria} }\DIFdelend \DIFaddbegin \subsection{Sleep deprivation and a gastric cancer microenvironment}\label{sec:case_sd_gastric}

\DIFadd{We then asked whether the same representation that produces the quantitative gains in Figs.~\ref{fig07}--\ref{fig08} also recovers biologically coherent latent structure on data that does not contribute to benchmark training. The first paired case }\DIFaddend (Fig.~\DIFdelbegin \DIFdel{\ref{fig10}A) , we computed per-dimension Pearson correlations with gene expression, selected the most strongly correlated genes, }\DIFdelend \DIFaddbegin \DIFadd{\ref{fig09}) compares mouse bone-marrow scRNA-seq from sleep-deprived }\DIFaddend and \DIFdelbegin \DIFdel{performed GO-BP enrichment analysis}\DIFdelend \DIFaddbegin \DIFadd{wild-type animals (GSE280145, panels A--H) with a gastric cancer microenvironment atlas (GSE183904, panels I--P), motivated by evidence that sleep state alters haematopoiesis and bone-marrow immune composition~\mbox{%DIFAUXCMD
\cite{mcalpine2019sleep,wang2024sleep}}\hskip0pt%DIFAUXCMD
}\DIFaddend .

\DIFaddbegin \DIFadd{In the sleep-deprivation card, curated WT/SD condition labels and curated hematopoietic cell-type labels are present in the source data. }\DIFaddend The retained latent \DIFdelbegin \DIFdel{components delineated biological programsperturbed by sleep deprivation (Fig.~\ref{fig10}B). Hematopoietic lineage control was captured by }\DIFdelend \DIFaddbegin \DIFadd{coordinates concentrate around hematopoietic-immune programs: $z_0$ is led by }\DIFaddend \textit{\DIFdelbegin \DIFdel{Ets1}%DIFDELCMD < \MBLOCKRIGHTBRACE %%%
\DIFdel{(Latent 0), a regulator of myeloid and lymphoid differentiation \mbox{%DIFAUXCMD
\cite{ets_transcription_2011}}\hskip0pt%DIFAUXCMD
, and }\DIFdelend \DIFaddbegin \DIFadd{Nrg2}}\DIFadd{, }\textit{\DIFadd{Mki67}}\DIFadd{, }\textit{\DIFadd{Gnal}}\DIFadd{, }\textit{\DIFadd{Pdcd}} \DIFadd{and }\textit{\DIFadd{Mob3p}}\DIFadd{; $z_1$ by }\textit{\DIFadd{Sync1}}\DIFadd{, }\textit{\DIFadd{Csf17}}\DIFadd{, }\textit{\DIFadd{Wfdc21}}\DIFadd{, }\textit{\DIFadd{Lcn2}} \DIFadd{and }\DIFaddend \textit{\DIFdelbegin \DIFdel{Bach2}\DIFdelend \DIFaddbegin \DIFadd{Ighv}\DIFaddend }\DIFdelbegin \DIFdel{(Latent 1), a transcriptional repressor central to B cell and lymphoid fate decisions \mbox{%DIFAUXCMD
\cite{bach2_b_cell_2016}}\hskip0pt%DIFAUXCMD
. Innate immune signaling was represented by }\textit{\DIFdel{Tyrobp}} %DIFAUXCMD
\DIFdel{(Latent 5), an adaptor protein for activating immune receptors \mbox{%DIFAUXCMD
\cite{tyrobp_immune_1998}}\hskip0pt%DIFAUXCMD
, and }\textit{\DIFdel{Clec4d}} %DIFAUXCMD
\DIFdel{(Latent 6), a member of the C-type lectin receptor family implicated in myeloid pattern recognition \mbox{%DIFAUXCMD
\cite{dectin1_fungal_2011}}\hskip0pt%DIFAUXCMD
. Additional components captured metabolic rewiring (}\textit{\DIFdel{Slc9a9}}%DIFAUXCMD
\DIFdel{, Latent 3), cell-cycle control (}\textit{\DIFdel{Cdk6}}%DIFAUXCMD
\DIFdel{, Latent 7), and extracellular matrix remodeling (}\textit{\DIFdel{Mmp9}}%DIFAUXCMD
\DIFdel{, Latent 8; }\textit{\DIFdel{Aff3}}%DIFAUXCMD
\DIFdel{, Latent 2), alongside }\DIFdelend \DIFaddbegin \DIFadd{; and $z_2$ by }\textit{\DIFadd{Trib}}\DIFadd{, }\textit{\DIFadd{Ctss}}\DIFadd{, }\textit{\DIFadd{Olfa2a}} \DIFadd{and }\textit{\DIFadd{Mmp9}}\DIFadd{. The matching GO Biological Process enrichment is dominated by defence response, inflammatory response, regulation of immune-system processes, lymphocyte activation and cell migration.
}

\DIFadd{In the gastric-cancer card, batch labels track the source samples; cell-state labels are inferred from cluster marker genes because no harmonised }\DIFaddend cell-type \DIFdelbegin \DIFdel{markers including }\textit{\DIFdel{Lrmp}} %DIFAUXCMD
\DIFdel{(Latent 4) and }\textit{\DIFdel{Wfdc21}} %DIFAUXCMD
\DIFdel{(Latent 9)\mbox{%DIFAUXCMD
\cite{cdk6_g1s_1994, mmp9_matrix_2007}}\hskip0pt%DIFAUXCMD
. GO-BP enrichment corroborated these assignments, with enriched terms spanning immune activation, defense responses, }\DIFdelend \DIFaddbegin \DIFadd{field is available. The inferred map spans epithelial populations (gastric, mucous, chief-like, MDK\textsuperscript{+} and broader epithelial states), stromal and vascular populations (fibroblast, collagen fibroblast, pericyte/smooth muscle, endothelial), and immune populations (SPP1\textsuperscript{+} macrophage, myeloid/macrophage, cytotoxic T/NK, activated CD4\textsuperscript{+} T cells with Treg-like signal, B cells, plasma cells, mast cells), with one low-confidence mixed gastric/}\DIFaddend T-cell \DIFdelbegin \DIFdel{differentiation, and inflammatory pathways (Fig.~\ref{fig10}C--F). CCVGAE thus captures latent structure consistent with biologically meaningful perturbations in hematopoietic tissues, though the correlational nature of the analysis precludes causal interpretation.
}%DIFDELCMD <

%DIFDELCMD < %%%
%DIF <  OMITTED_figure_FOR_DIFF_BUILD_STABILITY
%DIFDELCMD <

%DIFDELCMD < \subsection{Temporal Decomposition of Radiation-Injured Hematopoietic Stem Cells}
%DIFDELCMD <

%DIFDELCMD < %%%
\DIFdel{To assess CCVGAE under a longitudinal perturbation, we analyzed a murine hematopoietic stem cell (HSC) scRNA-seq dataset following ionizing radiation (GEO: GSE278673). This dataset profiles the transcriptional response from day 0 to day 30 post-irradiation, providing a rigorous testbed for modeling stress responses, lineage decisions, and tissue repair \mbox{%DIFAUXCMD
\cite{wang2024radiation, xiao2023radiation}}\hskip0pt%DIFAUXCMD
. }%DIFDELCMD <

%DIFDELCMD < %%%
\DIFdel{The dataset spans a dynamic cellular landscape over time (Fig.~\ref{fig11}A). CCVGAE decomposed it into latent components grouping into stress response, metabolic adaptation, and cellular protection modules (Fig.~\ref{fig11}B) \mbox{%DIFAUXCMD
\cite{lopez2023biologically}}\hskip0pt%DIFAUXCMD
.
}%DIFDELCMD <

%DIFDELCMD < %%%
\DIFdelend \DIFaddbegin \DIFadd{cluster left broad. The leading latents separate epithelial, stromal, immune and vascular axes through keratin-family genes (}\textit{\DIFadd{KRT8}}\DIFadd{, }\textit{\DIFadd{KRT18}}\DIFadd{, }\textit{\DIFadd{KRT19}}\DIFadd{, }\textit{\DIFadd{KRT15}}\DIFadd{), }\textit{\DIFadd{TPM1}}\DIFadd{, }\textit{\DIFadd{IGFBP2}}\DIFadd{, }\textit{\DIFadd{CALD1}}\DIFadd{, }\textit{\DIFadd{LYZ}}\DIFadd{, }\textit{\DIFadd{CD68}}\DIFadd{, }\textit{\DIFadd{AGR2}}\DIFadd{, }\textit{\DIFadd{SPARC}} \DIFadd{and }\textit{\DIFadd{ELF3}}\DIFadd{, with }\DIFaddend GO-BP \DIFdelbegin \DIFdel{enrichment of genes most strongly associated with each latent variable corroborated these assignments (Fig.~\ref{fig11}C)\mbox{%DIFAUXCMD
\cite{geneontology2021}}\hskip0pt%DIFAUXCMD
. Latent 0 and Latent 4 were enriched for ``Mitotic cell cycle, '' consistent with proliferative responses. Core hematopoietic functions were recovered \mbox{%DIFAUXCMD
\cite{orkin2008hematopoiesis}}\hskip0pt%DIFAUXCMD
: Latent 2 and Latent 5 mapped to ``Hemopoiesis'' and ``Myeloid cell differentiation, '' while Latent 3 aligned with ``Erythrocyte homeostasis.'' Latent 6 was enriched for ``Defense response'' and ``Myeloid leukocyte activation, '' capturing the acute inflammatory signature of radiation-induced hematopoietic injury \mbox{%DIFAUXCMD
\cite{li2022hsc}}\hskip0pt%DIFAUXCMD
. CCVGAE thus provides a component-wise decomposition of cellular variation consistent with established biology, though as in the sleep-deprivation analysis these latent--gene mappings are correlational and do not establish causal relationships}\DIFdelend \DIFaddbegin \DIFadd{terms clustering on collagen and extracellular-matrix organisation, cell adhesion, vasculature development and morphogenesis}\DIFaddend .


% OMITTED_figure_FOR_DIFF_BUILD_STABILITY


\DIFdelbegin %DIFDELCMD < \subsection{CCVGAE Resolves the TPO-Driven Megakaryocyte Differentiation Trajectory}
%DIFDELCMD < %%%
\DIFdelend \DIFaddbegin \subsection{TPO-induced cord-blood megakaryopoiesis and aged hematopoietic stem cells}\label{sec:case_ucb_hsc}
\DIFaddend

\DIFdelbegin \DIFdel{We next examined whether CCVGAE can resolve a continuous developmental trajectory, using a }\DIFdelend \DIFaddbegin \DIFadd{The second paired case (Fig.~\ref{fig10}) compares a D0--D14 }\DIFaddend thrombopoietin (TPO)\DIFdelbegin \DIFdel{time course of umbilical cord blood (UCB)cells undergoing megakaryocyte differentiation}\DIFdelend \DIFaddbegin \DIFadd{-induced umbilical-cord-blood megakaryocyte differentiation time course (GSE280270, panels A--H) with an aged-HSC dataset (GSE226131, panels I--P)}\DIFaddend . TPO is the principal cytokine \DIFdelbegin \DIFdel{governing }\DIFdelend \DIFaddbegin \DIFadd{that governs }\DIFaddend megakaryopoiesis and platelet production\DIFdelbegin \DIFdel{, providing a well-defined benchmark for temporal modeling \mbox{%DIFAUXCMD
\cite{kaushansky1994thrombopoietin, de_sauvage1994stimulation, lok1994cloning}}\hskip0pt%DIFAUXCMD
.
CCVGAE recovered the continuum from early progenitors (D0) to mature megakaryocytes (D14) and organized latent components along this trajectory (Fig.~\ref{fig12}A), highlighting coherent biological processes at stage-specific windows (Fig.~\ref{fig12}B). Transcriptional and long non-coding RNA (lncRNA)-mediated regulation was captured by }\textit{\DIFdel{NEAT1}} %DIFAUXCMD
\DIFdel{(Latent 0), }\textit{\DIFdel{NFKBIA}} %DIFAUXCMD
\DIFdel{(Latent 1), and }\textit{\DIFdel{NRIP1}} %DIFAUXCMD
\DIFdel{(Latent 8) \mbox{%DIFAUXCMD
\cite{chen2016neat1, zhang2019neat1_megakaryocyte}}\hskip0pt%DIFAUXCMD
. Signal transduction was represented by }\textit{\DIFdel{PLCB1}} %DIFAUXCMD
\DIFdel{(Latent 3) \mbox{%DIFAUXCMD
\cite{plcb1_signaling_2001}}\hskip0pt%DIFAUXCMD
, and intracellular membrane trafficking by }\textit{\DIFdel{DENND4A}} %DIFAUXCMD
\DIFdel{(Latent 4), a Rab-GEF involved in vesicle dynamics. Multiple components reflected the endomitotic cell cycle characteristic of megakaryocyte polyploidization, including }\textit{\DIFdel{CLSPN}}%DIFAUXCMD
\DIFdel{, }\textit{\DIFdel{CENPF}}%DIFAUXCMD
\DIFdel{, and }\textit{\DIFdel{CDKN2D}} %DIFAUXCMD
\DIFdel{\mbox{%DIFAUXCMD
\cite{ravid2002polyploid_megakaryocytes, lordier2008megakaryocyte_endomitosis}}\hskip0pt%DIFAUXCMD
. Platelet adhesion modules were captured by integrins }\textit{\DIFdel{ITGA4}} %DIFAUXCMD
\DIFdel{(Latent 7) and }\textit{\DIFdel{ITGA2B}} %DIFAUXCMD
\DIFdel{(Latent 9), the latter a definitive lineage marker \mbox{%DIFAUXCMD
\cite{shattil1998itga2b_function}}\hskip0pt%DIFAUXCMD
. GO-BP enrichment corroborated these assignments, with Latent 2 enriched for cell-cycle processes and Latent 7 for hemostasis and plateletactivation (Fig.~\ref{fig12}C--D). CCVGAE captures latent structure associated with both discrete regulatory programs and continuous developmental dynamics; as in the preceding case studies, these associations are correlational and should be regarded as hypothesis-generating rather than causal.
}%DIFDELCMD <

%DIFDELCMD < %%%
%DIF <  OMITTED_figure_FOR_DIFF_BUILD_STABILITY
%DIFDELCMD <

%DIFDELCMD < \section{Discussion}
%DIFDELCMD <

%DIFDELCMD < %%%
\DIFdel{We introduced CCVGAE, a deep generative framework that targets the interaction between inference-time strategies, latent-space regularization, and encoder architecture in single-cell representation learning. Across 170 datasets, }%DIFDELCMD < \textbf{Centroid Inference}%%%
\DIFdel{---using the deterministic posterior mean as the analysis-time embedding---gives consistent improvements in clustering accuracy, representation quality, and latent manifold geometry when combined with coupling regularization and graph attention encoding (Fig.~\ref{fig03}), reducing the effect of technical and sampling noise on biological signals \mbox{%DIFAUXCMD
\cite{zhang2020dissecting}}\hskip0pt%DIFAUXCMD
.
}%DIFDELCMD <

%DIFDELCMD < %%%
\subsection*{\DIFdel{Design Rationale and Component Contributions}}
%DIFAUXCMD
%DIFDELCMD <

%DIFDELCMD < %%%
\DIFdel{For clustering and visualization, stochastic sampling adds variance without benefit and can reduce reproducibility. }%DIFDELCMD < \textbf{Centroid Inference}%%%
\DIFdel{---using the posterior mean directly---retains the regularization benefits of variational training while removing inference-time noise, and is quantified here together with the complementary components across a large benchmark. }%DIFDELCMD <

%DIFDELCMD < %%%
\DIFdel{The }%DIFDELCMD < \textbf{Coupling} %%%
\DIFdel{mechanism, a dual-reconstruction bottleneck with $d_c < d_z$ (default $d_c = 5$, $d_z = 10$), enforces genuine dimensional compression. We chose it over simpler alternatives (e.g.
, increased KL weight or dropout) because the bottleneck explicitly forces the model to learn a compressed representation that preserves reconstruction-critical information, improving latent geometry. Ablation studies (Fig.~\ref{fig04}) show that coupling provides complementary gains in intrinsic quality and distance correlation. }%DIFDELCMD <

%DIFDELCMD < %%%
\DIFdel{The }%DIFDELCMD < \textbf{Graph Attention Network} %%%
\DIFdel{encoder leverages cell--cell similarity structure, which is often informative for biological interpretation \mbox{%DIFAUXCMD
\cite{long2023graph}}\hskip0pt%DIFAUXCMD
. Attention allows adaptive weighting of neighbors across heterogeneous cell populations. When graph construction is infeasible or computationally prohibitive, the MLP encoder alternative preserves the benefits of centroid inference and coupling regularization.
}%DIFDELCMD <

%DIFDELCMD < %%%
\subsection*{\DIFdel{Graph Attention and the Clustering--Geometry Trade-off}}
%DIFAUXCMD
\DIFdelend \DIFaddbegin \DIFadd{~\mbox{%DIFAUXCMD
\cite{kaushansky1994thrombopoietin,de_sauvage1994stimulation,lok1994cloning}}\hskip0pt%DIFAUXCMD
.
}\DIFaddend

\DIFdelbegin \DIFdel{Adding the graph attention encoder to centroid-based models can slightly decrease ARI}\DIFdelend \DIFaddbegin \DIFadd{In the TPO time-course card, curated time-point labels are provided by the source data and cell-state labels are inferred from cluster marker genes; the inferred states span early HSPC and progenitor populations, cycling and mitotic progenitors, platelet- and megakaryocyte-biased populations, FCER1A\textsuperscript{+} myeloid}\DIFaddend /\DIFdelbegin \DIFdel{NMI while improving representation quality and intrinsic geometry---a known trade-off in graph-based methods: message passing smooths local neighborhoods, blurring fine-grained cluster boundaries while improving manifold continuity and geometric regularity. The degree of this effect depends on graph construction parameters (particularly $k$ in the $k$-NN graph) and the inherent cluster separability of the dataset. For developmental studies and trajectory inference, where manifold continuity matters more than label sharpness, the graph attention component is preferable; for tasks requiring tight cluster separation, CenVAE is the better choice.
}\subsection*{\DIFdel{Scope and Interpretation of Quantitative Results}}
%DIFAUXCMD
\DIFdelend \DIFaddbegin \DIFadd{dendritic-cell-like populations, a basophil-like population, and broader leukocyte-like clusters that are kept low-confidence. The trajectory from proliferative cell-cycle programs to platelet/megakaryocyte programs is visible in the latent--gene mini-UMAPs: $z_0$ couples to }\textit{\DIFadd{PROS1}}\DIFadd{, }\textit{\DIFadd{UBE2G}}\DIFadd{, }\textit{\DIFadd{LAT}}\DIFadd{, }\textit{\DIFadd{CASC15}} \DIFadd{and }\textit{\DIFadd{SOS1}}\DIFadd{; $z_1$ to a platelet-axis cluster led by }\textit{\DIFadd{LTBP1}}\DIFadd{, }\textit{\DIFadd{ITGA8}}\DIFadd{, }\textit{\DIFadd{CMTM5}}\DIFadd{, }\textit{\DIFadd{ACTN1}} \DIFadd{and }\textit{\DIFadd{CD63}}\DIFadd{; and $z_2$ to }\textit{\DIFadd{ITGB3}}\DIFadd{, }\textit{\DIFadd{ITGA4}}\DIFadd{, }\textit{\DIFadd{KCNQ3}}\DIFadd{, }\textit{\DIFadd{DNM3}} \DIFadd{and }\textit{\DIFadd{TYMS}}\DIFadd{. The matching GO-BP terms recover chromatin organisation, platelet activation, hemostasis, body fluid regulation and mitotic cell cycle.
}\DIFaddend

\DIFdelbegin \DIFdel{Many improvements are statistically significant but modest in absolute terms. CCVGAE's primary contribution is }%DIFDELCMD < \textit{%%%
\DIFdel{consistency }\DIFdelend \DIFaddbegin \DIFadd{In the aged-HSC card, both group and cell-state labels are inferred and labelled as such. The inferred map spans HSC/MPP and broad hematopoietic progenitor populations, erythroid and cycling-erythroid populations, granulocyte progenitors and neutrophil/granulocytic myeloid populations, inflammatory myeloid and monocyte/macrophage populations, cycling cells, }\DIFaddend and \DIFdelbegin \DIFdel{stability}%DIFDELCMD < \MBLOCKRIGHTBRACE %%%
\DIFdel{rather than dramatic performance jumps on individual benchmarks.
}%DIFDELCMD <

%DIFDELCMD < %%%
\DIFdel{The biological case studies (Figs.~\ref{fig10}--\ref{fig12}) illustrate that geometric stability translates into more interpretable components: latent dimensions correlate with coherent processes including lineage control, immune activation, }\DIFdelend \DIFaddbegin \DIFadd{small T- }\DIFaddend and \DIFdelbegin \DIFdel{cell-cycle regulation. These interpretations rely on post-hoc gene correlation and GO enrichment rather than intrinsic model properties and should be treated as hypothesis-generating}\DIFdelend \DIFaddbegin \DIFadd{B-cell clusters. The leading latent-associated genes (}\textit{\DIFadd{Trem1}}\DIFadd{, }\textit{\DIFadd{Hp}}\DIFadd{, }\textit{\DIFadd{Tyrobp}}\DIFadd{, }\textit{\DIFadd{Ctsg}}\DIFadd{, }\textit{\DIFadd{Cd35}}\DIFadd{, }\textit{\DIFadd{Cebpb}}\DIFadd{, }\textit{\DIFadd{Plac8}}\DIFadd{, }\textit{\DIFadd{Mpo}}\DIFadd{, }\textit{\DIFadd{Anxa3}}\DIFadd{) align with defence response, neutrophil and myeloid immunity, leukocyte chemotaxis and cell-division programs}\DIFaddend .


\DIFdelbegin \subsection*{\DIFdel{Synthesis of Biological Case Studies}}
%DIFAUXCMD
\DIFdelend %DIF >  OMITTED_figure_FOR_DIFF_BUILD_STABILITY


\DIFdelbegin \DIFdel{In the sleep-deprived bone marrow analysis }\DIFdelend \DIFaddbegin \subsection{Radiation-injury hematopoiesis and the COVID-19 BALF immune landscape}\label{sec:case_ir_covid}

\DIFadd{The third paired case }\DIFaddend (Fig.~\DIFdelbegin \DIFdel{\ref{fig10}) , CCVGAE identified components corresponding to hematopoietic lineage control (}\textit{\DIFdel{Ets1}}%DIFAUXCMD
\DIFdel{, }\textit{\DIFdel{Bach2}}%DIFAUXCMD
\DIFdel{) , innate immune activation (}\textit{\DIFdel{Tyrobp}}%DIFAUXCMD
\DIFdel{, }\textit{\DIFdel{Clec4d}}%DIFAUXCMD
\DIFdel{), and metabolic rewiring (}\DIFdelend \DIFaddbegin \DIFadd{\ref{fig11}) compares a mouse hematopoietic stem cell radiation-injury time course covering d0--d30 after irradiation (GSE278673, panels A--H) with a COVID-19 bronchoalveolar lavage (BALF) immune-landscape dataset (GSE145926, panels I--P).
}

\DIFadd{In the radiation-injury card, curated time-point labels and curated hematopoietic cell-type labels are present in the source data. The top latent-associated genes (}\DIFaddend \textit{\DIFdelbegin \DIFdel{Slc9a9}%DIFDELCMD < \MBLOCKRIGHTBRACE%%%
\DIFdel{) }\DIFdelend \DIFaddbegin \DIFadd{Zbtb16}}\DIFaddend , \DIFdelbegin \DIFdel{providing a structured decomposition of the transcriptomic response to prolonged wakefulness.
In }\DIFdelend \DIFaddbegin \textit{\DIFadd{Gata1}}\DIFadd{, }\textit{\DIFadd{Klf1}}\DIFadd{, }\textit{\DIFadd{Cor}}\DIFadd{, }\textit{\DIFadd{Argl}}\DIFadd{, }\textit{\DIFadd{Gp1bb}}\DIFadd{, }\textit{\DIFadd{Cd34}}\DIFadd{, }\textit{\DIFadd{Ighm}}\DIFadd{, }\textit{\DIFadd{Rgs1}}\DIFadd{, }\textit{\DIFadd{Junb}}\DIFadd{) cluster on hematopoiesis, mononuclear-cell differentiation, defence response and regulation of immune-system processes~\mbox{%DIFAUXCMD
\cite{orkin2008hematopoiesis}}\hskip0pt%DIFAUXCMD
.
}

\DIFadd{In the COVID-19 BALF card, both group and cell-state labels are inferred and }\DIFaddend the \DIFdelbegin \DIFdel{radiation-injured HSC dataset (Fig.~\ref{fig11}), the framework recovered stress-response, metabolic adaptation, and lineage-commitment modules that reflected the temporal dynamics of post-irradiation recovery. The TPO-driven megakaryopoiesis study (Fig.~\ref{fig12}) demonstrated that CCVGAE resolves continuous developmental trajectories, capturing stage-specific programs from transcriptional regulation (}\textit{\DIFdel{NEAT1}}%DIFAUXCMD
\DIFdel{, }\textit{\DIFdel{NFKBIA}}%DIFAUXCMD
\DIFdel{)through endomitosis (}\textit{\DIFdel{CLSPN}}%DIFAUXCMD
\DIFdel{, }\textit{\DIFdel{CENPF}}%DIFAUXCMD
\DIFdel{)to terminal platelet adhesion (}\textit{\DIFdel{ITGA2B}}%DIFAUXCMD
\DIFdel{).
}%DIFDELCMD <

%DIFDELCMD < %%%
\DIFdel{Across all three studies, individual latent dimensions map to coherent biological programs rather than entangling unrelated processes---consistent with coupling regularization and centroid inference jointly promoting a structured latent space. The recovered programs span lineage regulation, immune signaling, cell-cycle control, and extracellular matrix remodeling, and GO enrichment corroborates the gene-level correlations in each case.
}%DIFDELCMD <

%DIFDELCMD < %%%
\DIFdel{These interpretations carry caveats. The correlation-based workflow identifies associations, not causal mechanisms; latent--gene correlations may reflect indirect or confounded relationships.The three case studies are drawn exclusively from hematopoietic systems under perturbation, and generalization to other tissues and contexts remains to be evaluated. Computational findings should be treated as hypotheses for experimental validation.
}%DIFDELCMD <

%DIFDELCMD < %%%
\subsection*{\DIFdel{Computational Considerations}}
%DIFAUXCMD
%DIFDELCMD <

%DIFDELCMD < %%%
\DIFdel{The graph attention encoder adds overhead over a standard MLP encoder due to attention coefficient computation and neighborhood aggregation; subgraph sampling (default $N_s = 512$ nodes per iteration) bounds per-iteration cost and enables scaling to large datasets regardless of total cell count. The dual-reconstruction pathway roughly doubles decoder forward passes during training but does not affect inference cost, since only the deterministic centroid path is used downstream. In terms of memory, the primary overhead arises from storing the cell--cell similarity graph ($O(N \cdot k)$ edges for a $k$-NN graph), which is sparse and manageable for typical dataset sizes. }%DIFDELCMD < \textbf{When to prefer CCVGAE over simpler methods:} %%%
\DIFdel{CCVGAE is most advantageous when (1) embedding stability and reproducibility are priorities---for example, when results must be consistent across repeated analyses or when downstream tasks are sensitive to embedding perturbations; (2) interpretability of latent dimensions is desired for biological hypothesis generation; (3) data are sparse or noisy (as in scATAC-seq), where the coupling regularization and centroid inference provide the largest marginal gains ; and (4) both cluster-level and trajectory-level analyses are needed from a single embedding. For rapid exploratory analysis of well-characterized, clean datasetswhere only basic clustering is required, simpler methods such as PCA followed by standard clustering may suffice with substantially lower computational cost.
Analogous concerns with reproducibility, uncertainty quantification, and deployment management have been explored in adjacent }\DIFdelend \DIFaddbegin \DIFadd{panels represent }\DIFaddend data-driven \DIFdelbegin \DIFdel{engineering domains~\mbox{%DIFAUXCMD
\cite{naeem2025infrastructure,naeem2025quality,naeem2025equipment,naeem2025bim,naeem2025digitaltwin}}\hskip0pt%DIFAUXCMD
.
}\subsection*{\DIFdel{Limitations}}
%DIFAUXCMD
%DIFDELCMD <

%DIFDELCMD < %%%
\DIFdel{Several limitations should be noted. First, our ablation baseline VAE uses stochastic samples at inference time to provide a controlled reference for isolating component contributions. Because some existing VAE-based methods already use posterior means for downstream tasks, the improvements attributed to centroid inference (CenVAE vs.VAE) represent the maximum benefit of switching from stochastic to deterministic inference within our architecture and may overestimate the advantage relative to such implementations. The more distinctive contributions of CCVGAE are the coupling regularization and graph attention components, which provide gains beyond what deterministic inference alone achieves.}%DIFDELCMD <

%DIFDELCMD < %%%
\DIFdel{Second, we have not systematically evaluated whether the centroid inference principle generalizes to arbitrary VAE architectures beyond those tested here; users should validate the approach within their specific workflows.
}%DIFDELCMD <

%DIFDELCMD < %%%
\DIFdel{Third, our evaluation relies on default hyperparameter settings optimized during development. The sensitivity analysis indicates that different settings favor different metric families rather than revealing a single universally optimal configuration; datasets with characteristics different from our benchmarks may therefore benefit from targeted tuning of the bottleneck dimension, coupling strength, KL weight, and graph construction parameters.
Fourth, the biological interpretations presented in our case studies are based on correlational analyses and should be validated through independent experimental approaches. The latent-gene correlations identify associations but do not establish causal relationships. }%DIFDELCMD <

%DIFDELCMD < %%%
\DIFdel{Fifth, the rapidly evolving landscape of single-cell representation learning means that newer methods (e.g., emerging foundation model variants and specialized graph-based integration frameworks) are not included in our comparisons. Future work should systematically update the comparison set as new methods become available. Analogous challenges in multi-objective trade-off resolution and long-horizon deployment have been examined in other applied fields~\mbox{%DIFAUXCMD
\cite{aqlan2025gametheory,aqlan2025conflict,naeem2025pavement,naeem2025labor,naeem2025sustainability}}\hskip0pt%DIFAUXCMD
.
}%DIFDELCMD <

%DIFDELCMD < %%%
\subsection*{\DIFdel{Hyperparameter Sensitivity}}
%DIFAUXCMD
%DIFDELCMD <

%DIFDELCMD < %%%
\DIFdel{Although our default hyperparameter configuration (Table~\ref{table_hyperparams}) was applied uniformly across all 170 benchmark datasets without per-dataset tuning, performance may be sensitive to specific parameter choices. We discuss the key hyperparameters and their expected effects:
}%DIFDELCMD <

%DIFDELCMD < \begin{itemize}
%DIFDELCMD <     \item \textbf{Bottleneck dimension ($d_c$):} %%%
\DIFdel{The ratio $d_c / d_z$ controls the information compression in the coupling path. Our default of $d_c = 5$ with $d_z = 10$ (ratio $= 0.5$) provides moderate compression. Decreasing $d_c$ enforces stricter compression, which may improve geometric regularity but risk discarding biologically relevant variation. Increasing $d_c$ toward $d_z$ reduces the coupling effect and converges toward a standard VAE with redundant reconstruction. Practitioners working with highly heterogeneous populations (many cell types) may benefit from a larger $d_c$ to retain fine-grained distinctions, while lower values may suffice for simpler populations.
}%DIFDELCMD <

%DIFDELCMD < 	\item \textbf{Loss weights ($w_{\text{recon}}, w_{\text{irecon}}, w_{\text{KL}}$):} %%%
\DIFdel{The three loss weights control different aspects of the clustering--geometry trade-off. Increasing $w_{\text{recon}}$ emphasizes reconstruction fidelity and tends to improve ARI/NMI, but it progressively reduces intrinsic metrics such as core quality. Increasing $w_{\text{irecon}}$ strengthens the coupling path; the sweep shows that setting $w_{\text{irecon}} = 0$ weakens intrinsic geometry, whereas moderate-to-high values mainly improve core quality with limited effect on ARI}\DIFdelend \DIFaddbegin \DIFadd{immune-state strata rather than author-provided clinical labels. The inferred map spans airway epithelial populations (IFN-response, secretory, ciliated, squamous and SERPINB3\textsuperscript{+} epithelial), T}\DIFaddend /\DIFdelbegin \DIFdel{NMI. Varying $w_{\text{KL}}$ shifts the balance between flexible cluster separation and latent regularity: smaller values favor ARI}\DIFdelend \DIFaddbegin \DIFadd{NK populations (T cells, CD8\textsuperscript{+} cytotoxic T, NK}\DIFaddend /\DIFdelbegin \DIFdel{NMI}\DIFdelend \DIFaddbegin \DIFadd{$\gamma\delta$ T), macrophage and monocyte populations (alveolar, APOE\textsuperscript{+} and inflammatory macrophages and inflammatory monocytes), activated dendritic cells and pDC}\DIFaddend /\DIFdelbegin \DIFdel{ASW, while larger values improve participation ratio and core quality at the cost of clustering performance. The equal-weight default should therefore be interpreted as a balanced compromise rather than a universal optimum. }%DIFDELCMD <

%DIFDELCMD <     \item \textbf{Graph construction ($k$-NN parameter):} %%%
\DIFdel{The neighborhood size $k=15$ determines graphconnectivity. Smaller $k$ values yield sparser graphs that preserve local structure but may undersmooth noise; larger $k$ values increase smoothing but can blur fine-grained boundaries between similar cell types.For datasets with well-separated clusters, smaller $k$ may be preferable; for continuous developmental trajectories, larger $k$ values can improve trajectory reconstruction.
}%DIFDELCMD <

%DIFDELCMD <     \item \textbf{Latent dimension ($d_z$):} %%%
\DIFdel{The default $d_z = 10$ balances representational capacity and parsimony. Datasets with many distinct cell types or complex regulatory programs may benefit from larger $d_z$, while simpler systems can be effectively captured with fewer dimensions.
Our interpretability workflow (Section~\ref{sec:interpretability}) provides a diagnostic tool: if multiple latent dimensions map to redundant biological processes, the latent space may be oversized. }%DIFDELCMD < \end{itemize}
%DIFDELCMD <

%DIFDELCMD < %%%
\DIFdel{The parameters $w_{\text{recon}}$}\DIFdelend \DIFaddbegin \DIFadd{B-like populations, and lower-confidence CXCL10\textsuperscript{+} or FCER1A\textsuperscript{+} mixed clusters. The strongest latent-associated genes (}\textit{\DIFadd{CORO1A}}\DIFadd{, }\textit{\DIFadd{MARCKS}}\DIFadd{, }\textit{\DIFadd{MARCO}}\DIFaddend , \DIFdelbegin \DIFdel{$d_z$, and $d_c$ jointly govern the capacity--regularization trade-off and should be co-varied in any comprehensive sensitivity sweep. To empirically characterize these trade-offs, we conducted a one-at-a-time sweep of the three loss weights ($w_{\text{recon}}$}\DIFdelend \DIFaddbegin \textit{\DIFadd{APOC1}}\DIFadd{, }\textit{\DIFadd{IRG1}}\DIFaddend , \DIFdelbegin \DIFdel{$w_{\text{irecon}}$}\DIFdelend \DIFaddbegin \textit{\DIFadd{CD8E}}\DIFaddend , \DIFdelbegin \DIFdel{$w_{\text{KL}}$) across all 55~scRNA-seq benchmark datasets (Fig.~\ref{fig13} and Table~\ref{table_sensitivity}) . The sweeps show that no single loss-weight setting dominates every metric suite. Higher $w_{\text{recon}}$ progressively improves ARI and NMI but lowers core quality and participation ratio. For $w_{\text{irecon}}$, clustering metrics remain relatively stable once the coupling term is active, but removing it ($w_{\text{irecon}} = 0$) noticeably weakens core quality, while larger values strengthen intrinsic geometry. For $w_{\text{KL}}$, low regularization ($0.01$--$0.1$) favors ARI/NMI/ASW, whereas stronger regularization ($2.0$}\DIFdelend \DIFaddbegin \textit{\DIFadd{RSAD2}}\DIFadd{, }\textit{\DIFadd{CTSL}}\DIFadd{, }\textit{\DIFadd{CTSB}}\DIFadd{, }\textit{\DIFadd{CTSC}}\DIFadd{) point to lymphocyte and T-cell activation, macrophage-state variation, type-I-interferon-driven antiviral programs and defence response.
}

\DIFadd{The latent--gene mappings reported in Figs.~\ref{fig09}}\DIFaddend --\DIFdelbegin \DIFdel{$5.0$) improves core quality at the cost of cluster separation. We therefore retain the equal-weight configuration ($w_{\text{recon}} = w_{\text{irecon}} = w_{\text{KL}} = 1.0$) as a practical compromise across metric families, not as a claim of parameter invariance}\DIFdelend \DIFaddbegin \DIFadd{\ref{fig11} are correlational and should be read as hypothesis-generating rather than causal: the workflow does not perturb the system and the GO Biological Process enrichment is descriptive. Their role here is to show that the representation that produces the quantitative cross-method gains also recovers biologically coherent latent structure on independent datasets}\DIFaddend .


% OMITTED_figure_FOR_DIFF_BUILD_STABILITY


%DIF <  OMITTED_table_FOR_DIFF_BUILD_STABILITY
\DIFdelbegin %DIFDELCMD <

%DIFDELCMD < %%%
\subsection*{\DIFdel{Generalization and Clinical Translation}}
%DIFAUXCMD
%DIFDELCMD <

%DIFDELCMD < %%%
\DIFdel{Our benchmarks primarily comprise well-characterized datasets with established cell-type annotations; performance on novel biological systems---such as rare diseases, understudied organisms, or highly heterogeneous tumor samples---remains to be evaluated. For clinical translation, several considerations are pertinent:
}%DIFDELCMD < \begin{itemize}
\begin{itemize}%DIFAUXCMD
%DIFDELCMD <     \item \textbf{Quality control:} %%%
\item%DIFAUXCMD
\DIFdel{CCVGAE assumes reasonable data quality; clinical samples with excessive batch effects or low cell viability may require additional preprocessing before reliable embeddings can be obtained.
    }%DIFDELCMD < \item \textbf{Interpretability validation:} %%%
\item%DIFAUXCMD
\DIFdel{Clinical decision-making should rely on experimentally validated biomarkers rather than solely on computational predictions.
    }%DIFDELCMD < \item \textbf{Model updates:} %%%
\item%DIFAUXCMD
\DIFdel{The current model is trained per-dataset; transfer learning or fine-tuning strategies for applying pre-trained models to new user-provided data represent important future directions for long-term deployment.
}
\end{itemize}%DIFAUXCMD
%DIFDELCMD < \end{itemize}
%DIFDELCMD < %%%
\DIFdelend \DIFaddbegin \section{Discussion}
\DIFaddend

\DIFdelbegin \subsection*{\DIFdel{Scalability Analysis}}
%DIFAUXCMD
\DIFdelend \DIFaddbegin \DIFadd{scCCVGBen tests the design premise that a stable, interpretable cell embedding can be obtained by combining deterministic posterior-mean inference with a coupling-regularised dimensional bottleneck and a graph-attention encoder. The component-wise ablation (Fig.~\ref{fig03}) and the joint analysis (Fig.~\ref{fig04}) together support the design rationale: centroid inference is the strongest single contributor to clustering compactness, coupling regularisation is the strongest single contributor to neighbourhood-preserving embedding, and the largest improvement on intrinsic latent geometry only emerges when all three mechanisms act together. Posterior-mean inference is a known practice in some VAE implementations, but the present results indicate that its measured benefit is clearest when paired with components that shape the underlying latent geometry rather than with stochastic inference alone.
}\DIFaddend

\DIFdelbegin \DIFdel{The computational complexity of CCVGAE scales with dataset size as follows. For the MLP encoder, complexity is $O(N \cdot D \cdot H)$ where $N$ is cell count, $D$ is feature dimension, }\DIFdelend \DIFaddbegin \DIFadd{The robustness analyses (Figs.~\ref{fig05}--\ref{fig06}) qualify these gains. Within the attention family, scCCVGBen is essentially indistinguishable from GATv2 on every metric, }\DIFaddend and \DIFdelbegin \DIFdel{$H$ is hidden layer size. The GAT encoder adds $O(|E| \cdot H \cdot K)$ for edge processing, where $|E|$ is edge count and $K$ is attention heads. Subgraph sampling limits per-iteration cost to $O(N_s^2)$ regardless of full graph size.
}\DIFdelend \DIFaddbegin \DIFadd{GATv2 may therefore replace the default attention parameterisation without a measurable benchmark cost. Across families, GAT outperforms convolutional and propagation backbones on clustering compactness, but message-passing encoders such as GIN and EdgeConv outperform GAT on intrinsic geometry. The choice of cell--cell graph similarly matters: SNN and kNN-cosine outperform the kNN-Euclidean default on clustering compactness, while differences on neighbourhood-preserving embedding are smaller. Reporting an architectural axis alongside the algorithmic core therefore gives a more explicit account of where the gains come from. The graph-attention encoder's performance is bounded by the quality of the input cell--cell similarity graph; Fig.~\ref{fig06} assesses this sensitivity across five similarity definitions calibrated to a comparable average node degree, and finds that the model preserves its ranking against single-component baselines across all five constructions, with the largest sensitivity concentrated in clustering compactness rather than neighbourhood-preserving embedding.
}\DIFaddend

\DIFdelbegin \DIFdel{To empirically characterize scaling behavior, we profiled CCVGAE across five representative }\DIFdelend \DIFaddbegin \DIFadd{Cross-method comparisons (Figs.~\ref{fig07}--\ref{fig08}) place scCCVGBen at the highest aggregate rank on both modalities at the present data cut, with the largest single improvement on }\DIFaddend scRNA-seq \DIFdelbegin \DIFdel{datasets (2,531--22,602 cells) and five scATAC-seq datasets (3,350--15,259 cells), varying the number of selected features (500--5,000) and subgraph size (128--1, 024) on an NVIDIA RTX 3090 GPU (Fig.~\ref{fig14} andTable~\ref{table_runtime}). Per-epoch training time scales linearly with cell count and is largely insensitive to feature count, indicating that computation is dominated by graph operations rather than feature-space reconstruction. Peak GPU memory scales linearly with both feature count and subgraph size but remains independent of total cell count, reflecting the subgraph sampling strategy. Across all configurations, per-epoch time ranges from approximately 0.34--3.25~s, and peak memory remains below 210~MB, demonstrating practical feasibility for datasets of moderate size. For very large datasets ($>$100}\DIFdelend \DIFaddbegin \DIFadd{recorded against scVI (ASW $+0.341$}\DIFaddend , \DIFdelbegin \DIFdel{000 cells), we recommend: (1) increasing subgraph size gradually; (2) using the MLP encoder if graph construction becomes prohibitive; (3) applying minibatch-based graph construction.
}%DIFDELCMD <

%DIFDELCMD < %%%
%DIF <  OMITTED_figure_FOR_DIFF_BUILD_STABILITY
%DIFDELCMD <

%DIFDELCMD < %%%
%DIF <  OMITTED_table_FOR_DIFF_BUILD_STABILITY
%DIFDELCMD <

%DIFDELCMD < %%%
\subsection*{\DIFdel{Practical Implications and Future Directions}}
%DIFAUXCMD
%DIFDELCMD <

%DIFDELCMD < %%%
\DIFdel{Centroid Inference is conceptually straightforward and can be incorporated into other VAE-based workflows, though its benefits have been validated specifically within the CCVGAE architecture and benchmark conditions described here; whether equivalent gains arise in other architectures requires independent evaluation.
Performance on }\DIFdelend \DIFaddbegin \DIFadd{intrinsic-overall $+0.331$) and the largest single improvement on }\DIFaddend scATAC-seq \DIFdelbegin \DIFdel{suggests applicability to other sparse modalities such as single-cell proteomics and spatial transcriptomics \mbox{%DIFAUXCMD
\cite{moses2022museum}}\hskip0pt%DIFAUXCMD
. Gains in local structure preservation are relevant for trajectory inference, which depends on accurate local neighborhoods \mbox{%DIFAUXCMD
\cite{saelens2019comparison}}\hskip0pt%DIFAUXCMD
. Recent work on autoencoders for genomic variation analysis \mbox{%DIFAUXCMD
\cite{geleta2026autoencoders} }\hskip0pt%DIFAUXCMD
further validates the utility of centroid-based latent representations for high-dimensional biological data, complementing the deterministic inference strategy employed by CCVGAE.
}\DIFdelend \DIFaddbegin \DIFadd{intrinsic geometry recorded against PeakVI ($+0.346$). The neighbourhood-preserving suite is tighter, so scCCVGBen matches or exceeds the highest-ranked deep generative baselines in these panels without claiming a uniform win on every metric column. The three paired hematopoietic case studies (Figs.~\ref{fig09}--\ref{fig11}) show that the same representation recovers coherent latent--gene programs aligned with hematopoietic regulation, epithelial--stromal organisation, megakaryocyte differentiation and antiviral macrophage activity, supporting the use of scCCVGBen for downstream biological discovery alongside the quantitative benchmark.
}\DIFaddend

\DIFdelbegin \DIFdel{Future directions include exploring multi-scale graph learning techniques, investigating the utility of deterministic centroids as anchors for potential cross-modal integration tasks \mbox{%DIFAUXCMD
\cite{efremova2021computational, regev2017human}}\hskip0pt%DIFAUXCMD
, and systematic characterization of hyperparameter sensitivity across dataset characteristics. We note that the current work evaluates }\DIFdelend \DIFaddbegin \DIFadd{Three limitations bound the present analysis. First, }\DIFaddend scRNA-seq and \DIFdelbegin \DIFdel{scATAC-seq as separate benchmarking pools; extending CCVGAE to joint multi-omic integration from paired measurements (e.g., 10x Multiome RNA}\DIFdelend \DIFaddbegin \DIFadd{scATAC-seq are evaluated as separate cohorts rather than as paired modalities; extending the analysis to simultaneous RNA}\DIFaddend +\DIFdelbegin \DIFdel{ATAC) represents an important future direction. The stability and interpretability offered by CCVGAE are directly relevant to large-scale single-cell atlas projects, where embedding reproducibility across datasets and analysis runs is a practical requirement}\DIFdelend \DIFaddbegin \DIFadd{ATAC measurements is a defined future extension. Second, the encoder and graph-construction analyses hold the loss schedule and other hyperparameters at their defaults, so Figs.~\ref{fig05} and~\ref{fig06} test architectural rather than hyperparameter sensitivity; a hyperparameter-sweep analysis is reported in supplementary material. Third, the case studies are correlational because the workflow does not perturb the system and the GO Biological Process enrichment is descriptive, so the latent--gene mappings should be read as hypothesis-generating.
}

\DIFadd{Six restricted-access scATAC-seq submissions remain redacted in public exports, with their species assignment retained as auditable metadata. All comparisons are paired Wilcoxon signed-rank tests with Holm correction, and confidence intervals are obtained by paired bootstrap with 5}{\DIFadd{,}}\DIFadd{000 resamples under a fixed random seed. The full release scaffolding for the source code, the dataset cohort, the reconciled result tables and the rendered figures is described in the Code Availability and Data Availability sections below}\DIFaddend .

\section{Conclusions}

We have introduced \DIFdelbegin \DIFdel{CCVGAE, a deep generative framework }\DIFdelend \DIFaddbegin \DIFadd{scCCVGBen, a centroid-coupled variational graph attention autoencoder }\DIFaddend for single-cell representation learning that \DIFdelbegin \DIFdel{systematically combines deterministic inference, coupling regularization, and graph attention encoding. Across 170 benchmark datasets spanning }\DIFdelend \DIFaddbegin \DIFadd{combines deterministic posterior-mean inference, a coupling-regularised dimensional bottleneck and a graph-attention encoder. By replacing reparameterised samples with the posterior mean at inference time, scCCVGBen removes a source of sampling variance that destabilises stochastic VAE outputs, while the coupling bottleneck and graph attention together shape a latent geometry that is both more clusterable and more interpretable. Within a decoupled benchmark that varies the algorithmic core, the encoder backbone, the graph-construction strategy, the dataset cohort and the evaluation suite as independent axes, scCCVGBen improves the Average Silhouette Width by $+0.288$ and the intrinsic-overall score by $+0.233$ over a stochastic VAE on the paired }\DIFaddend scRNA-seq \DIFaddbegin \DIFadd{cohort, by $+0.341$ and $+0.331$ over scVI in the same comparison, }\DIFaddend and \DIFaddbegin \DIFadd{by $+0.139$ on ASW and $+0.346$ on intrinsic-overall against PeakVI on the paired }\DIFaddend scATAC-seq \DIFdelbegin \DIFdel{, using the deterministic posterior mean as the final embedding---together with a dimensional bottleneck ($d_c < d_z$) and }\DIFdelend \DIFaddbegin \DIFadd{cohort. Robustness analyses across fourteen encoder backbones and five graph-construction strategies indicate where alternative architectural choices remain competitive. Three paired hematopoietic case studies show that the same model recovers coherent latent--gene programs on held-out data, so the embeddings carry over to clustering, trajectory inference and downstream biological interpretation under the same protocol used in the benchmark.
}

\section*{\DIFadd{Code Availability}}

\DIFadd{The full implementation of scCCVGBen is released as an open-source Python package together with the training scripts, the baseline harness, the figure-generation pipeline, the benchmark manifest used to assemble Figs.~\ref{fig05}--\ref{fig08} and the audit scripts used to verify the submission. The source repository is hosted at }\url{https://github.com/PeterPonyu/scCCVGBen}\DIFadd{, and }\DIFaddend a \DIFdelbegin \DIFdel{GAT-based encoder---yields more stable, reproducible, and geometrically regular latent spaces, with particularly strong gains on sparse scATAC-seq data. Quantitatively, CCVGAE improves clustering accuracy (ASW $+0.329$ versus scVI on }\DIFdelend \DIFaddbegin \DIFadd{deployed companion site that hosts the per-dataset metadata cards, the interactive figure browser and the public-graph integration is available at }\url{https://peterponyu.github.io/scccvgben-next/}\DIFadd{. A self-contained 30-line notebook in the repository reproduces the PBMC-3k case end-to-end on CPU in under three minutes; the rendered companion-site assets are summarised in Supplementary Figs.~S1 and S2.
}

\section*{\DIFadd{Data Availability}}

\DIFadd{The 200-dataset benchmark cohort is sourced from the Gene Expression Omnibus and the European Nucleotide Archive; 100 }\DIFaddend scRNA-seq \DIFdelbegin \DIFdel{, $p < 0.001$), embedding quality (overall UMAP score $+0.219$ versus scVI), and intrinsic manifold geometry (core quality $+0.273$ versus PeakVI on }\DIFdelend \DIFaddbegin \DIFadd{accessions and 100 }\DIFaddend scATAC-seq \DIFdelbegin \DIFdel{) relative to established deep generative baselines. In the hematopoietic case studies examined here, these improvements correspond to interpretable latent components that separate distinct cellular programs.
}\DIFdelend \DIFaddbegin \DIFadd{accessions are listed individually with cell counts, feature counts and download links in Supplementary Table~S1, which is generated programmatically from the released benchmark manifest so that every column is version-controlled. Six restricted-access scATAC-seq submissions from a BCG-vaccination study are retained as redacted records with public-safe identifiers, with their species assignment held in auditable metadata. The six datasets reserved for the biological case studies and excluded from every benchmark statistic are GSE280145 (sleep-deprivation bone marrow), GSE280270 (TPO-induced cord-blood megakaryocyte time course), GSE278673 (radiation-injury hematopoietic stem cell time course), GSE183904 (gastric cancer microenvironment), GSE226131 (aged hematopoietic stem cell) and GSE145926 (COVID-19 bronchoalveolar lavage immune landscape).
}\DIFaddend

\DIFdelbegin \DIFdel{The centroid inference strategy points to an underexplored design choice in variational autoencoders: for many single-cell tasks, inference-time stochasticity is unnecessary and often harmful. The coupling bottleneck shows that explicit dimensional compression encourages geometric regularity without sacrificing reconstruction fidelity. CCVGAE's modular architecture lets practitioners choose the configuration matching their goals---centroid inference alone for simplicity, centroid plus coupling for improved geometry, or the full pipeline---without retraining, with subgraph sampling enabling scaling to atlas-size datasets.
}\DIFdelend %DIF > %%%%%%%%%%%%%%%%%%%%%%%%%%%%%%%%%%%%%%%%%%%%%%%%%%%%%%%%%%%%%%%%%%%%%
%DIF > % Appendix
%DIF > %%%%%%%%%%%%%%%%%%%%%%%%%%%%%%%%%%%%%%%%%%%%%%%%%%%%%%%%%%%%%%%%%%%%%

\DIFdelbegin \DIFdel{Looking ahead, several directions merit exploration. Extending CCVGAE to joint multi-omic integration from paired measurements (e.g., 10x Multiome RNA+ATAC) would test whether centroid inference and coupling regularization generalize beyond single-modality benchmarks. Transfer-learning strategies that reuse pre-trained centroid encoders across datasets could reduce per-experiment training costs and improve embedding consistency across cohorts. Finally, coupling the framework with spatial transcriptomics and longitudinal single-cell assays would allow the interpretability workflow to be evaluated against spatially or temporally resolved ground truth. The combination of embedding stability, geometric regularity, and component-level interpretability provides a foundation on which each of these extensions can be developed without redesigning the core model.
}\DIFdelend \DIFaddbegin \begin{appendices}
\DIFaddend

\DIFdelbegin %DIFDELCMD < \section{List of abbreviations}
%DIFDELCMD < %%%
\DIFdel{The following abbreviations are used in this manuscript:
}%DIFDELCMD < \begin{multicols}{2}
%DIFDELCMD < \begin{itemize}
\begin{itemize}%DIFAUXCMD
%DIFDELCMD <     \item \textbf{ARI}%%%
\item%DIFAUXCMD
\DIFdel{: Adjusted Rand Index
    }%DIFDELCMD < \item \textbf{AS}%%%
\item%DIFAUXCMD
\DIFdel{: Anisotropy Score
    }%DIFDELCMD < \item \textbf{ASW}%%%
\item%DIFAUXCMD
\DIFdel{: Average Silhouette Width
    }%DIFDELCMD < \item \textbf{CAL}%%%
\item%DIFAUXCMD
\DIFdel{: Calinski-Harabasz Index
    }%DIFDELCMD < \item \textbf{COR}%%%
\item%DIFAUXCMD
\DIFdel{: Correlation Score
    }%DIFDELCMD < \item \textbf{DAV}%%%
\item%DIFAUXCMD
\DIFdel{: Davies-Bouldin Index
    }%DIFDELCMD < \item \textbf{DC}%%%
\item%DIFAUXCMD
\DIFdel{: Data Compactness
    }%DIFDELCMD < \item \textbf{DICL}%%%
\item%DIFAUXCMD
\DIFdel{: Dictionary Learning
    }%DIFDELCMD < \item \textbf{FA}%%%
\item%DIFAUXCMD
\DIFdel{: Factor Analysis
    }%DIFDELCMD < \item \textbf{GAT}%%%
\item%DIFAUXCMD
\DIFdel{: Graph Attention Network
    }%DIFDELCMD < \item \textbf{GAT-VAE}%%%
\item%DIFAUXCMD
\DIFdel{: Graph Attention Network Variational Autoencoder (feature-reconstructing; distinct from adjacency-reconstructing VGAE)
    }%DIFDELCMD < \item \textbf{GNN}%%%
\item%DIFAUXCMD
\DIFdel{: Graph Neural Network
    }%DIFDELCMD < \item \textbf{GO-BP}%%%
\item%DIFAUXCMD
\DIFdel{: Gene Ontology Biological Process
    }%DIFDELCMD < \item \textbf{HSC}%%%
\item%DIFAUXCMD
\DIFdel{: Hematopoietic Stem Cell
    }%DIFDELCMD < \item \textbf{ICA}%%%
\item%DIFAUXCMD
\DIFdel{: Independent Component Analysis
    }%DIFDELCMD < \item \textbf{KL}%%%
\item%DIFAUXCMD
\DIFdel{: Kullback-Leibler
    }%DIFDELCMD < \item \textbf{k-NN}%%%
\item%DIFAUXCMD
\DIFdel{: $k$-Nearest Neighbors
    }%DIFDELCMD < \item \textbf{KPCA}%%%
\item%DIFAUXCMD
\DIFdel{: Kernel Principal Component Analysis
    }%DIFDELCMD < \item \textbf{lncRNA}%%%
\item%DIFAUXCMD
\DIFdel{: long non-coding RNA
    }%DIFDELCMD < \item \textbf{LSI}%%%
\item%DIFAUXCMD
\DIFdel{: Latent Semantic Indexing
    }%DIFDELCMD < \item \textbf{MD}%%%
\item%DIFAUXCMD
\DIFdel{: Manifold Dimensionality
    }%DIFDELCMD < \item \textbf{MLP}%%%
\item%DIFAUXCMD
\DIFdel{: Multi-Layer Perceptron
    }%DIFDELCMD < \item \textbf{NB}%%%
\item%DIFAUXCMD
\DIFdel{: Negative Binomial
    }%DIFDELCMD < \item \textbf{NMF}%%%
\item%DIFAUXCMD
\DIFdel{: Non-negative Matrix Factorization
    }%DIFDELCMD < \item \textbf{NMI}%%%
\item%DIFAUXCMD
\DIFdel{: Normalized Mutual Information
    }%DIFDELCMD < \item \textbf{NR}%%%
\item%DIFAUXCMD
\DIFdel{: Noise Resilience
    }%DIFDELCMD < \item \textbf{OV}%%%
\item%DIFAUXCMD
\DIFdel{: Overall Quality Score
    }%DIFDELCMD < \item \textbf{PCA}%%%
\item%DIFAUXCMD
\DIFdel{: Principal Component Analysis
    }%DIFDELCMD < \item \textbf{PR}%%%
\item%DIFAUXCMD
\DIFdel{: Participation Ratio
    }%DIFDELCMD < \item \textbf{QG}%%%
\item%DIFAUXCMD
\DIFdel{: Global Quality Score
    }%DIFDELCMD < \item \textbf{QL}%%%
\item%DIFAUXCMD
\DIFdel{: Local Quality Score
    }%DIFDELCMD < \item \textbf{scATAC-seq}%%%
\item%DIFAUXCMD
\DIFdel{: single-cell Assay for Transposase-Accessible Chromatin sequencing
    }%DIFDELCMD < \item \textbf{scRNA-seq}%%%
\item%DIFAUXCMD
\DIFdel{: single-cell RNA sequencing
    }%DIFDELCMD < \item \textbf{SD}%%%
\item%DIFAUXCMD
\DIFdel{: Sleep-Deprived
    }%DIFDELCMD < \item \textbf{SDR}%%%
\item%DIFAUXCMD
\DIFdel{: Spectral Decay Rate
    }%DIFDELCMD < \item \textbf{SOTA}%%%
\item%DIFAUXCMD
\DIFdel{: State-of-the-Art
    }%DIFDELCMD < \item \textbf{TD}%%%
\item%DIFAUXCMD
\DIFdel{: Trajectory Directionality
    }%DIFDELCMD < \item \textbf{TPO}%%%
\item%DIFAUXCMD
\DIFdel{: Thrombopoietin
    }%DIFDELCMD < \item \textbf{TSVD}%%%
\item%DIFAUXCMD
\DIFdel{: Truncated Singular Value Decomposition
    }%DIFDELCMD < \item \textbf{t-SNE}%%%
\item%DIFAUXCMD
\DIFdel{: t-distributed Stochastic Neighbor Embedding
    }%DIFDELCMD < \item \textbf{UCB}%%%
\item%DIFAUXCMD
\DIFdel{: Umbilical Cord Blood
    }%DIFDELCMD < \item \textbf{UMAP}%%%
\item%DIFAUXCMD
\DIFdel{: Uniform Manifold Approximation and Projection
    }%DIFDELCMD < \item \textbf{VAE}%%%
\item%DIFAUXCMD
\DIFdel{: Variational Autoencoder
    }%DIFDELCMD < \item \textbf{WT}%%%
\item%DIFAUXCMD
\DIFdel{: Wild-Type
}
\end{itemize}%DIFAUXCMD
%DIFDELCMD < \end{itemize}
%DIFDELCMD < \end{multicols}
%DIFDELCMD < %%%
\DIFdelend \DIFaddbegin \section{Abbreviations}\label{app:abbrev}
\DIFaddend

\DIFdelbegin %DIFDELCMD < \section{Declarations}
%DIFDELCMD < %%%
\DIFdelend \DIFaddbegin \begingroup
\small
\setlength{\LTpre}{0pt}
\setlength{\LTpost}{0pt}
\DIFaddend

\DIFdelbegin \subsection*{\DIFdel{Ethics declaration}}
%DIFAUXCMD
\DIFdel{Not applicable.
}\DIFdelend %DIF >  OMITTED_longtable_FOR_DIFF_BUILD_STABILITY

\DIFdelbegin \subsection*{\DIFdel{Consent to Participate declaration}}
%DIFAUXCMD
\DIFdel{Not applicable.
}\DIFdelend \DIFaddbegin \endgroup
\DIFaddend

\DIFdelbegin \subsection*{\DIFdel{Consent to Publish declaration}}
%DIFAUXCMD
\DIFdel{Not applicable.
}\DIFdelend \DIFaddbegin \section{Default hyperparameters}\label{app:hyperparams}
\DIFaddend


\DIFdelbegin \subsection*{\DIFdel{Data Availability declaration}}
%DIFAUXCMD
\DIFdel{All original source codes are available on GitHub at }%DIFDELCMD < \url{https://github.com/PeterPonyu/CCVGAE}%%%
\DIFdel{.
}\DIFdelend %DIF >  OMITTED_table_FOR_DIFF_BUILD_STABILITY


\DIFdelbegin \DIFdel{A total of 55 single-cell RNA sequencing datasets were retrieved from the Gene Expression Omnibus (GEO) database for analysis.
}\DIFdelend \DIFaddbegin \section{Supplementary Material}\label{app:supp}
\DIFaddend

\DIFdelbegin %DIFDELCMD < \begin{itemize}
%DIFDELCMD < 	\item %%%
\DIFdel{Sousa, C., Poovathingal, S. K., Kaoma, T., Azuaje, F., Skupin, A. and Michelucci, A. Single-cell transcriptomics reveals distinct microglia signatures under inflammation. GEO GSE115571. }%DIFDELCMD < \url{https://www.ncbi.nlm.nih.gov/geo/query/acc.cgi?acc=GSE115571}
%DIFDELCMD < 	\item %%%
\DIFdel{Simon, L. M. and Schiller, H. B. Single cell RNA sequencing analysis of fresh resected human lung tissue. GEO GSE130148. }%DIFDELCMD < \url{https://www.ncbi.nlm.nih.gov/geo/query/acc.cgi?acc=GSE130148}
%DIFDELCMD < 	\item %%%
\DIFdel{Zhang, S., Cui, Y., Wen, L., Qiao, J. and Tang, F. Single-cell transcriptomics reveals the divergent developmental lineage trajectories during human pituitary development. GEO GSE142653. }%DIFDELCMD < \url{https://www.ncbi.nlm.nih.gov/geo/query/acc.cgi?acc=GSE142653}
%DIFDELCMD < 	\item %%%
\DIFdel{Joseph, D. B., et al. Single-cell RNA-sequencing of adult mouse lower urinary tracts. GEO GSE145929. }%DIFDELCMD < \url{https://www.ncbi.nlm.nih.gov/geo/query/acc.cgi?acc=GSE145929} %%%
\DIFdel{*Provides two data samples
	}%DIFDELCMD < \item %%%
\DIFdel{Yuan, S., Liu, Q., Hu, Z. and Mao, X. Single-cell RNA-sequencing reveals the heterogeneity of microglia in fibrous membrane derived from proliferative diabetic retinopathy and proliferative vitreoretinopathy. GEO GSE165784. }%DIFDELCMD < \url{https://www.ncbi.nlm.nih.gov/geo/query/acc.cgi?acc=GSE165784}
%DIFDELCMD < 	\item %%%
\DIFdel{Hou, J., Bi, H., Jiang, Q. and Gu, X. Heterogeneity analysis of astrocytes following spinal cord injury at single-cell resolution. GEO GSE189070. }%DIFDELCMD < \url{https://www.ncbi.nlm.nih.gov/geo/query/acc.cgi?acc=GSE189070}
%DIFDELCMD < 	\item %%%
\DIFdel{Zhang, B., Zeng, K., Guan, R. and Yang, Y. Single-cell RNA-seq analysis reveals macrophage involvement in pathogenesis of human sporadic type A aortic dissection. GEO GSE213740. }%DIFDELCMD < \url{https://www.ncbi.nlm.nih.gov/geo/query/acc.cgi?acc=GSE213740}
%DIFDELCMD < 	\item %%%
\DIFdel{Zhang, Z., et al. A panoramic view of cell population dynamics in mammalian aging. GEO GSE247719. }%DIFDELCMD < \url{https://www.ncbi.nlm.nih.gov/geo/query/acc.cgi?acc=GSE247719} %%%
\DIFdel{*Provides two datasamples
	}%DIFDELCMD < \item %%%
\DIFdel{Zhao, T., Sun, Z., Zhong, Q., Yu, X., Sun, T. and An, Z. Decoding SFRP2hi fibroblast progenitors in sustaining tooth growth in humans and mice at single-cell resolution. GEO GSE275119. }%DIFDELCMD < \url{https://www.ncbi.nlm.nih.gov/geo/query/acc.cgi?acc=GSE275119}
%DIFDELCMD < 	\item %%%
\DIFdel{Bastidas-Ponce, A., et al. Comprehensive single-cell mRNA profiling reveals a detailed roadmap for pancreatic endocrinogenesis. GEO GSE132188. }%DIFDELCMD < \url{https://www.ncbi.nlm.nih.gov/geo/query/acc.cgi?acc=GSE132188}
%DIFDELCMD < 	\item %%%
\DIFdel{Bandyopadhyay, S., Ahn, K. J., Duffy, M., Qin, L. and Tan, K. Mapping the cellular biogeography of human bone marrow niches using single-cell transcriptomics. GEO GSE253355. }%DIFDELCMD < \url{https://www.ncbi.nlm.nih.gov/geo/query/acc.cgi?acc=GSE253355}
%DIFDELCMD < 	\item %%%
\DIFdel{Alkaslasi, M. R., et al. Single nucleus RNA-sequencing defines unexpected diversity of cholinergic neuron types in the adult mouse spinal cord. GEO GSE167597. }%DIFDELCMD < \url{https://www.ncbi.nlm.nih.gov/geo/query/acc.cgi?acc=GSE167597}
%DIFDELCMD < 	\item %%%
\DIFdel{Wang, H., et al. Decoding human megakaryocyte development. GEO GSE144024. }%DIFDELCMD < \url{https://www.ncbi.nlm.nih.gov/geo/query/acc.cgi?acc=GSE144024}
%DIFDELCMD < 	\item %%%
\DIFdel{Spildrejorde, M., et al. Gene expression analysis of hESCs undergoing neuronal differentiation }%DIFDELCMD < [%%%
\DIFdel{scRNA-seq}%DIFDELCMD < ]%%%
\DIFdel{. GEO GSE192857. }%DIFDELCMD < \url{https://www.ncbi.nlm.nih.gov/geo/query/acc.cgi?acc=GSE192857}
%DIFDELCMD < 	\item %%%
\DIFdel{Cosgrove, J. and P\'{e}rie, L. Single-cell RNA sequencing of hematopoietic stem and progenitor cells from young and aged mice. GEO GSE226131. }%DIFDELCMD < \url{https://www.ncbi.nlm.nih.gov/geo/query/acc.cgi?acc=GSE226131}
%DIFDELCMD < 	\item %%%
\DIFdel{Fast, E. External signals regulate continuous transcriptional states in hematopoietic stem cells. GEO GSE165844. }%DIFDELCMD < \url{https://www.ncbi.nlm.nih.gov/geo/query/acc.cgi?acc=GSE165844}
%DIFDELCMD < 	\item %%%
\DIFdel{Simon, L. M. and Schiller, H. B. Longitudinal single-cell transcriptomics analysis of mouse lung upon bleomycin-induced injury. GEO GSE141259. }%DIFDELCMD < \url{https://www.ncbi.nlm.nih.gov/geo/query/acc.cgi?acc=GSE141259}
%DIFDELCMD < 	\item %%%
\DIFdel{Linnarsson, S.
Transcriptome analysis of single cells from the mouse dentate gyrus. GEO GSE95753. }%DIFDELCMD < \url{https://www.ncbi.nlm.nih.gov/geo/query/acc.cgi?acc=GSE95753}
%DIFDELCMD < 	\item %%%
\DIFdel{Demerdash, Y., Bouman, B. J., Haghverdi, L. and Essers, M. A. Time series single-cell RNA sequencing of murine hematopoietic stem and progenitors (HSPCs) following in vivo IFN$\alpha$ treatment. GEO GSE226824. }%DIFDELCMD < \url{https://www.ncbi.nlm.nih.gov/geo/query/acc.cgi?acc=GSE226824}
%DIFDELCMD < 	\item %%%
\DIFdel{Setty, M. Profiling of CD34+ cells from human bone marrow to understand hematopoiesis. ENA PRJEB37166. }%DIFDELCMD < \url{https://www.ebi.ac.uk/ena/browser/view/PRJEB37166}
%DIFDELCMD < 	\item %%%
\DIFdel{Zhu, Y., Wang, T., Gu, J. and Pan, G. Characterization and generation of human definitive multipotent hematopoietic stem/progenitor cells. GEO GSE148215. }%DIFDELCMD < \url{https://www.ncbi.nlm.nih.gov/geo/query/acc.cgi?acc=GSE148215}
%DIFDELCMD < 	\item %%%
\DIFdel{Dillon, L. Human bone marrow assessment by single-cell RNA sequencing, mass cytometry and flow cytometry. GEO GSE120446. }%DIFDELCMD < \url{https://www.ncbi.nlm.nih.gov/geo/query/acc.cgi?acc=GSE120446}
%DIFDELCMD < 	\item %%%
\DIFdel{Teo, Y. V., Hinthorn, S. J., Webb, A. E. and Nicola, N. Single-cell RNA sequencing of aging peripheral blood. GEO GSE120505. }%DIFDELCMD < \url{https://www.ncbi.nlm.nih.gov/geo/query/acc.cgi?acc=GSE120505}
%DIFDELCMD < 	\item %%%
\DIFdel{Paulson, K. G. and Chapuis, A. G. scRNA-seq reveals mechanisms of Merkel cell carcinoma acquired immunotherapy resistance }%DIFDELCMD < [%%%
\DIFdel{1}%DIFDELCMD < ]%%%
\DIFdel{. GEO GSE117988. }%DIFDELCMD < \url{https://www.ncbi.nlm.nih.gov/geo/query/acc.cgi?acc=GSE117988} %%%
\DIFdel{*Provides two data samples
	}%DIFDELCMD < \item %%%
\DIFdel{Ghobrial, I. M. and Getz, G. Single-cell RNA sequencing reveals compromised immune microenvironment in precursor stages of multiple myeloma. GEO GSE124310. }%DIFDELCMD < \url{https://www.ncbi.nlm.nih.gov/geo/query/acc.cgi?acc=GSE124310}
%DIFDELCMD < 	\item %%%
\DIFdel{Zhang, M., Yang, H., Liu, B.and Yan, X. Dissecting the transcriptomic landscape of human intrahepatic cholangiocarcinoma by single-cell RNA sequencing. GEO GSE138709. }%DIFDELCMD < \url{https://www.ncbi.nlm.nih.gov/geo/query/acc.cgi?acc=GSE138709}
%DIFDELCMD < 	\item %%%
\DIFdel{Yanagawa, J., et al. Single-cell RNA sequencing of human early-stage lung adenocarcinoma with activating somatic KRAS mutations and associated normal lung tissues reveals alterations of key gene expression in alveolar type II pneumocytes. GEO GSE149655. }%DIFDELCMD < \url{https://www.ncbi.nlm.nih.gov/geo/query/acc.cgi?acc=GSE149655}
%DIFDELCMD < 	\item %%%
\DIFdel{Jiang, H., Yu, D. and Yang, P. Transcriptional heterogeneity in primary and metastatic gastric cancer revealed using single-cell RNA sequencing. GEO GSE163558. }%DIFDELCMD < \url{https://www.ncbi.nlm.nih.gov/geo/query/acc.cgi?acc=GSE163558}
%DIFDELCMD < 	\item %%%
\DIFdel{Lei, P., Pereira, E. R., Beyaz, S. and Padera, T. P. Single-cell sequencing reveals cancer cell heterogeneity in a murine breast cancer lymph node metastasis model. GEO GSE168181. }%DIFDELCMD < \url{https://www.ncbi.nlm.nih.gov/geo/query/acc.cgi?acc=GSE168181}
%DIFDELCMD < 	\item %%%
\DIFdel{Jiang, T., Sun, S., Fan, Y. and Zhu, J. Spatiotemporal transcriptional atlas of lung adenocarcinoma from adenocarcinoma in situ to invasive carcinoma }%DIFDELCMD < [%%%
\DIFdel{scRNA-seq}%DIFDELCMD < ]%%%
\DIFdel{. GEO GSE189357. }%DIFDELCMD < \url{https://www.ncbi.nlm.nih.gov/geo/query/acc.cgi?acc=GSE189357}
%DIFDELCMD < 	\item %%%
\DIFdel{Wang, F. and Long, J. Single-cell and spatial transcriptome analysis reveals the cellular heterogeneity of liver metastatic colorectal cancer. GEO GSE225857. }%DIFDELCMD < \url{https://www.ncbi.nlm.nih.gov/geo/query/acc.cgi?acc=GSE225857}
%DIFDELCMD < 	\item %%%
\DIFdel{Han, K. , Joo, E. and Park, W. Single-cell RNA sequencing of primary breast cancer. GEO GSE228499. }%DIFDELCMD < \url{https://www.ncbi.nlm.nih.gov/geo/query/acc.cgi?acc=GSE228499}
%DIFDELCMD < 	\item %%%
\DIFdel{M\"{u}nter, D., et al. Multiomic analysis uncovers a continuous spectrum of differentiation and Wnt-MDK-driven immune evasion in hepatoblastoma }%DIFDELCMD < [%%%
\DIFdel{snRNA-seq}%DIFDELCMD < ]%%%
\DIFdel{. GEO GSE283205. }%DIFDELCMD < \url{https://www.ncbi.nlm.nih.gov/geo/query/acc.cgi?acc=GSE283205}
%DIFDELCMD < 	\item %%%
\DIFdel{Mehtonen, J., et al. Single-cell characterization of arrested B-lymphoid differentiation and leukemic cell states in ETV6-RUNX1-positive pediatric leukemia. GEO GSE148218. }%DIFDELCMD < \url{https://www.ncbi.nlm.nih.gov/geo/query/acc.cgi?acc=GSE148218}
%DIFDELCMD < 	\item %%%
\DIFdel{Li, Z. and Luo, L. Single-cell RNA sequencing identifies molecular biomarkers predicting response to CDK4/6 inhibition in metastatic HR+/HER2- breast cancer. GEO GSE262288. }%DIFDELCMD < \url{https://www.ncbi.nlm.nih.gov/geo/query/acc.cgi?acc=GSE262288}
%DIFDELCMD < 	\item %%%
\DIFdel{Geldhof, V., et al. Single-cell atlas identifies lipid-processing and immunomodulatory endothelial cells in healthy and malignant breast. GEO GSE155109. }%DIFDELCMD < \url{https://www.ncbi.nlm.nih.gov/geo/query/acc.cgi?acc=GSE155109} %%%
\DIFdel{*Provides two data samples
	}%DIFDELCMD < \item %%%
\DIFdel{Yost, K. E., et al. Clonal replacement of tumor-specific T cells following PD-1 blockade }%DIFDELCMD < [%%%
\DIFdel{single cells}%DIFDELCMD < ]%%%
\DIFdel{. GEO GSE123813. }%DIFDELCMD < \url{https://www.ncbi.nlm.nih.gov/geo/query/acc.cgi?acc=GSE123813} %%%
\DIFdel{*Provides two data samples
	}%DIFDELCMD < \item %%%
\DIFdel{Wang, L., Dai, J., Han, R. and Jin, W. Single-cell map of diverse immune phenotypes in the metastatic brain tumor microenvironment of non-small cell lung cancer and triple-negative breast cancer. GEO GSE143423. }%DIFDELCMD < \url{https://www.ncbi.nlm.nih.gov/geo/query/acc.cgi?acc=GSE143423} %%%
\DIFdel{*Provides two data samples
	}%DIFDELCMD < \item %%%
\DIFdel{Laughney, A . M., et al. The single-cell transcriptional landscape of human lung adenocarcinoma (primary tumors and metastases). GEO GSE123902. }%DIFDELCMD < \url{https://www.ncbi.nlm.nih.gov/geo/query/acc.cgi?acc=GSE123902}
%DIFDELCMD < 	\item %%%
\DIFdel{Liu, Y., Ge, J., Chen, Y. and Yu, K. Single-cell profiling and spatial transcriptome of primary tumors and paired metastatic lymph nodes in breast cancer patients. GEO GSE225600. }%DIFDELCMD < \url{https://www.ncbi.nlm.nih.gov/geo/query/acc.cgi?acc=GSE225600}
%DIFDELCMD < 	\item %%%
\DIFdel{Yang, J., Yu, J., Huang, X. and Li, Y. Single-cell network pharmacology predicts total therapies targeting multiple developmental clones in B-cell acute lymphoblastic leukemia. GEO GSE235787. }%DIFDELCMD < \url{https://www.ncbi.nlm.nih.gov/geo/query/acc.cgi?acc=GSE235787}
%DIFDELCMD < 	\item %%%
\DIFdel{Chen, X. and Li, Z. Targeting tumor cells toward the antigenic specificity of bystander T cells in tumor microenvironment potentiates cancer immunotherapy. GEO GSE222002. }%DIFDELCMD < \url{https://www.ncbi.nlm.nih.gov/geo/query/acc.cgi?acc=GSE222002}
%DIFDELCMD < 	\item %%%
\DIFdel{Zheng, C., et al. Landscape of infiltrating T cells in liver cancer revealed by single-cell sequencing. GEO GSE98638. }%DIFDELCMD < \url{https://www.ncbi.nlm.nih.gov/geo/query/acc.cgi?acc=GSE98638}
%DIFDELCMD < 	\item %%%
\DIFdel{Cichocki, F. and Day, A. Nicotinamide enhances natural killer cell function and yields remissions in patients with non-Hodgkin lymphoma. GEO GSE222369. }%DIFDELCMD < \url{https://www.ncbi.nlm.nih.gov/geo/query/acc.cgi?acc=GSE222369}
%DIFDELCMD < 	\item %%%
\DIFdel{Kumar, V., Ramnarayanan, K. and Tan, P. Single-cell atlas of lineage states, tumor microenvironment and subtype-specific expression programs in gastric cancer. GEO GSE183904. }%DIFDELCMD < \url{https://www.ncbi.nlm.nih.gov/geo/query/acc.cgi?acc=GSE183904}
%DIFDELCMD < 	\item %%%
\DIFdel{Caron, M., et al. Single-cell analysis of childhood leukemia reveals a link between developmental states and ribosomal protein expression as a source of intra-individual heterogeneity. GEO GSE132509. }%DIFDELCMD < \url{https://www.ncbi.nlm.nih.gov/geo/query/acc.cgi?acc=GSE132509}
%DIFDELCMD < 	\item %%%
\DIFdel{Poscablo, D. M., Forsberg, E. C.}\DIFdelend %DIF >  Keep the supplementary material on an S-indexed sequence rather than the
%DIF >  appendix section sequence (C1, C2, ...). The first longtable below uses an
%DIF >  unnumbered caption with an explicit title, but longtable still advances the
%DIF >  table counter, so resetting here yields Supplementary Table S1 followed by
%DIF >  the effect-size tables as Tables S2--S10.
\DIFaddbegin \renewcommand{\thefigure}{S\arabic{figure}}
\renewcommand{\thetable}{S\arabic{table}}
\setcounter{figure}{0}
\setcounter{table}{0}

\DIFadd{This consolidated supplementary section contains the full dataset table
(Supplementary Table~S1), the complete twenty-metric effect-size and
significance tables behind every main benchmark figure (Supplementary
Tables~S2--S10), }\DIFaddend and \DIFdelbegin \DIFdel{Medina, P. An age-progressive platelet differentiation path from hematopoietic stem cells causes exacerbated thrombosis. GEO GSE255019. }%DIFDELCMD < \url{https://www.ncbi.nlm.nih.gov/geo/query/acc.cgi?acc=GSE255019}
%DIFDELCMD < 	\item %%%
\DIFdel{Fu, Z. LSK immunophenotype HSC in Dapp1 knockout mice: bone marrow scRNA-seq. GEO GSE277292}\DIFdelend \DIFaddbegin \DIFadd{the two auxiliary online-resource figures included here
(Supplementary Figs.~S1 and S2)}\DIFaddend .
\DIFdelbegin %DIFDELCMD < \url{https://www.ncbi.nlm.nih.gov/geo/query/acc.cgi?acc=GSE277292}
%DIFDELCMD < 	\item %%%
\DIFdel{Fu, Z. LSK immunophenotype HSC after irradiation: bone marrow scRNA-seq. GEO GSE278673.
}%DIFDELCMD < \url{https://www.ncbi.nlm.nih.gov/geo/query/acc.cgi?acc=GSE278673}
%DIFDELCMD < \end{itemize}
%DIFDELCMD < %%%
\DIFdelend

\DIFdelbegin \DIFdel{Furthermore, 115 single-cell ATAC sequencing datasets from the GEO database were included in this study.
}\DIFdelend %DIF >  OMITTED_INPUT_TABLE_FOR_DIFF_BUILD_STABILITY


\DIFdelbegin %DIFDELCMD < \begin{itemize}
%DIFDELCMD < 	\item %%%
\DIFdel{Qiu X, Li Y, Brown M. Enhancer reactivation mediates adaptive resistance to FGFR inhibitors in triple-negative breast cancer }%DIFDELCMD < [%%%
\DIFdel{scATAC\_seq}%DIFDELCMD < ]%%%
\DIFdel{. GEO GSE168026. }%DIFDELCMD < \url{https://www.ncbi.nlm.nih.gov/geo/query/acc.cgi?acc=GSE168026} %%%
\DIFdel{*Provides two samples
	}%DIFDELCMD < \item %%%
\DIFdel{Morabito S, Miyoshi E, Swarup V. Single-nucleus chromatin accessibility and transcriptomic characterization of Alzheimer's Disease. GEO GSE174367. }%DIFDELCMD < \url{https://www.ncbi.nlm.nih.gov/geo/query/acc.cgi?acc=GSE174367}
%DIFDELCMD < 	\item %%%
\DIFdel{Wu LM, Lu QR. Transcriptional programs dictating Schwann cell transformation in MPNST }%DIFDELCMD < [%%%
\DIFdel{scATACseq}%DIFDELCMD < ]%%%
\DIFdel{. GEO GSE178988. }%DIFDELCMD < \url{https://www.ncbi.nlm.nih.gov/geo/query/acc.cgi?acc=GSE178988} %%%
\DIFdel{*Provides three samples
	}%DIFDELCMD < \item %%%
\DIFdel{Spildrejorde M, et al. Single-cell ATAC-seq analysis of hESCs undergoing neuronal differentiation }%DIFDELCMD < [%%%
\DIFdel{scATAC-seq}%DIFDELCMD < ]%%%
\DIFdel{. GEO GSE192856. }%DIFDELCMD < \url{https://www.ncbi.nlm.nih.gov/geo/query/acc.cgi?acc=GSE192856} %%%
\DIFdel{*Provides two samples
	}%DIFDELCMD < \item %%%
\DIFdel{Jimenez E, et al. A regulatory network of Sox and Six transcription factors initiate a cell fate transformation during hearing regeneration in adult zebrafish. GEO GSE192947. }%DIFDELCMD < \url{https://www.ncbi.nlm.nih.gov/geo/query/acc.cgi?acc=GSE192947} %%%
\DIFdel{*Provides five samples
	}%DIFDELCMD < \item %%%
\DIFdel{Khateb M, et al. Transcriptomics, Regulatory Syntax, and Enhancer Identification in Heterogenous Populations of Mesoderm-Induced ESCs at Single-Cell Resolution. GEO GSE198730. }%DIFDELCMD < \url{https://www.ncbi.nlm.nih.gov/geo/query/acc.cgi?acc=GSE198730} %%%
\DIFdel{*Provides three samples
	}%DIFDELCMD < \item %%%
\DIFdel{Giles J, Wherry EJ. Longitudinal single cell transcriptional and epigenetic mapping of effector, memory, and exhausted CD8 T cells reveals shared biological circuits across distinct cell fates }%DIFDELCMD < [%%%
\DIFdel{scATAC-Seq}%DIFDELCMD < ]%%%
\DIFdel{. GEO GSE199556. }%DIFDELCMD < \url{https://www.ncbi.nlm.nih.gov/geo/query/acc.cgi?acc=GSE199556} %%%
\DIFdel{*Provides two samples
	}%DIFDELCMD < \item %%%
\DIFdel{Li S, Yang M, Teng S, Wang D. Single-cell chromatin accessible state of paired primary and liver metastasis colorectal cancer cells from PDX cells derived mouse model. GEO GSE200813. }%DIFDELCMD < \url{https://www.ncbi.nlm.nih.gov/geo/query/acc.cgi?acc=GSE200813} %%%
\DIFdel{*Provides five samples
	}%DIFDELCMD < \item %%%
\DIFdel{Li F, et al. Single cell ATAC-seq of mouse hair follicle morphogenesis. GEO GSE201213. }%DIFDELCMD < \url{https://www.ncbi.nlm.nih.gov/geo/query/acc.cgi?acc=GSE201213} %%%
\DIFdel{*Provides four samples
	}%DIFDELCMD < \item %%%
\DIFdel{Lin YH, et al. Small intestine and colon tissue-resident memory CD8+ T cells exhibit transcriptional, epigenetic, and functional heterogeneity in concert with differential dependence on Eomesodermin }%DIFDELCMD < [%%%
\DIFdel{scATAC-seq}%DIFDELCMD < ]%%%
\DIFdel{. GEO GSE205549. }%DIFDELCMD < \url{https://www.ncbi.nlm.nih.gov/geo/query/acc.cgi?acc=GSE205549} %%%
\DIFdel{*Provides two samples
	}%DIFDELCMD < \item %%%
\DIFdel{Wang G, Zhang D, Qin L, Huang B. Chemical-based external stimulation reprograms BJs into FiNs (scATAC-seq). GEO GSE206767. }%DIFDELCMD < \url{https://www.ncbi.nlm.nih.gov/geo/query/acc.cgi?acc=GSE206767}
%DIFDELCMD < 	\item %%%
\DIFdel{Itahashi K, Irie T, Nishikawa H. Integrated multiomics profiling identifies the differentiation program of regulatory T cells in human tumors }%DIFDELCMD < [%%%
\DIFdel{scATAC-seq}%DIFDELCMD < ]%%%
\DIFdel{. GEO GSE211087. }%DIFDELCMD < \url{https://www.ncbi.nlm.nih.gov/geo/query/acc.cgi?acc=GSE211087} %%%
\DIFdel{*Provides four samples
	}%DIFDELCMD < \item %%%
\DIFdel{Peter S, W Dean P. Temporal chromatin accessibility changes define transcriptional states essential for osteosarcoma metastasis }%DIFDELCMD < [%%%
\DIFdel{scATAC-Seq}%DIFDELCMD < ]%%%
\DIFdel{. GEO GSE215758. }%DIFDELCMD < \url{https://www.ncbi.nlm.nih.gov/geo/query/acc.cgi?acc=GSE215758}
%DIFDELCMD < 	\item %%%
\DIFdel{Spildrejorde M, et al. Single-cell ATAC-seq analysis of hESCs undergoing neuronal differentiation exposed to paracetamol }%DIFDELCMD < [%%%
\DIFdel{scATAC-seq}%DIFDELCMD < ]%%%
\DIFdel{. GEO GSE220026. }%DIFDELCMD < \url{https://www.ncbi.nlm.nih.gov/geo/query/acc.cgi?acc=GSE220026} %%%
\DIFdel{*Provides two samples
	}%DIFDELCMD < \item %%%
\DIFdel{Winter DR, Perlman H. Tissue-resident, extravascular Ly6c- monocytes are critical for inflammation in the synovium. GEO GSE225801. }%DIFDELCMD < \url{https://www.ncbi.nlm.nih.gov/geo/query/acc.cgi?acc=GSE225801} %%%
\DIFdel{*Provides two samples
	}%DIFDELCMD < \item %%%
\DIFdel{Liu Q, Dong W. scATAC-seq of antigen-specific CD8+ T cells from different time-course of LCMV-Armstrong infected mice. GEO GSE226067. }%DIFDELCMD < \url{https://www.ncbi.nlm.nih.gov/geo/query/acc.cgi?acc=GSE226067} %%%
\DIFdel{*Provides five samples
	}%DIFDELCMD < \item %%%
\DIFdel{Chen R, et al. A multi-omics atlas of the human retina at single-cell resolution. GEO GSE226108. }%DIFDELCMD < \url{https://www.ncbi.nlm.nih.gov/geo/query/acc.cgi?acc=GSE226108} %%%
\DIFdel{*Provides five samples
	}%DIFDELCMD < \item %%%
\DIFdel{Laurie SJ, et al. Multiomic Single Cell Evaluation Reveals Inflammatory Cytokines Affect Innate Lymphoid Cell Fate After Allogeneic Stem Cell Transplantation }%DIFDELCMD < [%%%
\DIFdel{scATAC-Seq}%DIFDELCMD < ]%%%
\DIFdel{. GEO GSE232002. }%DIFDELCMD < \url{https://www.ncbi.nlm.nih.gov/geo/query/acc.cgi?acc=GSE232002} %%%
\DIFdel{*Provides three samples
	}%DIFDELCMD < \item %%%
\DIFdel{Zhang M. ARID1B regulates BCL11B-mediated non-canonical Activin signaling to maintain mesenchymal stem cell homeostasis in the mouse incisor }%DIFDELCMD < [%%%
\DIFdel{scATAC-seq}%DIFDELCMD < ]%%%
\DIFdel{. GEO GSE237304. }%DIFDELCMD < \url{https://www.ncbi.nlm.nih.gov/geo/query/acc.cgi?acc=GSE237304}
%DIFDELCMD < 	\item %%%
\DIFdel{Liu Q, Dong W. scATAC-seq of OT-1 CD8+ T cells from different time-course of Lm-OVA infected mice. GEO GSE241714. }%DIFDELCMD < \url{https://www.ncbi.nlm.nih.gov/geo/query/acc.cgi?acc=GSE241714} %%%
\DIFdel{*Provides three samples
	}%DIFDELCMD < \item %%%
\DIFdel{Wu J, Castro LN, Suva ML, Bernstein BE. Single-cell transcription and chromatin accessibility profiles of IDH mutant gliomas. GEO GSE241745. }%DIFDELCMD < \url{https://www.ncbi.nlm.nih.gov/geo/query/acc.cgi?acc=GSE241745} %%%
\DIFdel{*Provides four samples
	}%DIFDELCMD < \item %%%
\DIFdel{Racine L, Parmentier R. Metabolic adaptation pilots the differentiation of human hematopoietic cells (scATAC-Seq). GEO GSE243004. }%DIFDELCMD < \url{https://www.ncbi.nlm.nih.gov/geo/query/acc.cgi?acc=GSE243004} %%%
\DIFdel{*Provides four samples
	}%DIFDELCMD < \item %%%
\DIFdel{Chen R, et al. Single-cell multiomics of the human retina reveals hierarchical transcription factor collaboration in mediating cell type-specific effects of genetic variants on gene regulation. GEO GSE247157. }%DIFDELCMD < \url{https://www.ncbi.nlm.nih.gov/geo/query/acc.cgi?acc=GSE247157} %%%
\DIFdel{*Provides seven samples
	}%DIFDELCMD < \item %%%
\DIFdel{BCG vaccine immune response scATAC-seq data (restricted accession--restricted accession). GEO accession pending public release. *Provides six samples
	}%DIFDELCMD < \item %%%
\DIFdel{Dony L, et al. Chronic exposure to glucocorticoids amplifies inhibitory neuron cell fate during human neurodevelopment in organoids }%DIFDELCMD < [%%%
\DIFdel{scATAC-seq}%DIFDELCMD < ]%%%
\DIFdel{. GEO GSE252523. }%DIFDELCMD < \url{https://www.ncbi.nlm.nih.gov/geo/query/acc.cgi?acc=GSE252523}
%DIFDELCMD < 	\item %%%
\DIFdel{Poluben L, Balk SP, Russo J, Nouri M. Increased chromatin accessibility drives transition to androgen receptor splice variant dependence in castration-resistant prostate cancer }%DIFDELCMD < [%%%
\DIFdel{scATAC-Seq}%DIFDELCMD < ]%%%
\DIFdel{. GEO GSE252843. }%DIFDELCMD < \url{https://www.ncbi.nlm.nih.gov/geo/query/acc.cgi?acc=GSE252843}
%DIFDELCMD < 	\item %%%
\DIFdel{Hegde S, et al. Myeloid progenitor dysregulation fuels immunosuppressive macrophages in tumours. GEO GSE270148. }%DIFDELCMD < \url{https://www.ncbi.nlm.nih.gov/geo/query/acc.cgi?acc=GSE270148}
%DIFDELCMD < 	\item %%%
\DIFdel{Min Dai. An enhancer-AAV toolbox to target and manipulate distinct interneuron subtypes }%DIFDELCMD < [%%%
\DIFdel{basal ganglia}%DIFDELCMD < ]%%%
\DIFdel{. GEO GSE272181. }%DIFDELCMD < \url{https://www.ncbi.nlm.nih.gov/geo/query/acc.cgi?acc=GSE272181}
%DIFDELCMD < 	\item %%%
\DIFdel{Stomper J, et al. Sex differences in DNMT3A-mutant clonal hematopoiesis and the effects of estrogen }%DIFDELCMD < [%%%
\DIFdel{scATAC-Seq}%DIFDELCMD < ]%%%
\DIFdel{. GEO GSE272265. }%DIFDELCMD < \url{https://www.ncbi.nlm.nih.gov/geo/query/acc.cgi?acc=GSE272265} %%%
\DIFdel{*Provides three samples
	}%DIFDELCMD < \item %%%
\DIFdel{Baek S, et al. Single-cell multi-omics reveals tumor microenvironmentfactors underlying poor immunotherapy responses in ALK-positive lung cancer. GEO GSE274934. }%DIFDELCMD < \url{https://www.ncbi.nlm.nih.gov/geo/query/acc.cgi?acc=GSE274934} %%%
\DIFdel{*Provides four samples
	}%DIFDELCMD < \item %%%
\DIFdel{Min Dai. An enhancer-AAV toolbox to target and manipulate distinct interneuron subtypes }%DIFDELCMD < [%%%
\DIFdel{P28}%DIFDELCMD < ]%%%
\DIFdel{. GEO GSE277242. }%DIFDELCMD < \url{https://www.ncbi.nlm.nih.gov/geo/query/acc.cgi?acc=GSE277242}
%DIFDELCMD < 	\item %%%
\DIFdel{Xu L, Korogi Y, Li R, Sun X. RUNX2 promotes pulmonary fibrosis through the conversion of alveolar fibroblasts into pathological fibroblasts. GEO GSE278419. }%DIFDELCMD < \url{https://www.ncbi.nlm.nih.gov/geo/query/acc.cgi?acc=GSE278419} %%%
\DIFdel{*Provides two samples
	}%DIFDELCMD < \item %%%
\DIFdel{Kim D, Blackshaw S. Decoding Gene Networks Governing Hypothalamic and Prethalamic Neuron Development. GEO GSE284492. }%DIFDELCMD < \url{https://www.ncbi.nlm.nih.gov/geo/query/acc.cgi?acc=GSE284492} %%%
\DIFdel{*Provides three samples
	}%DIFDELCMD < \item %%%
\DIFdel{Shen W. Single-cell ATAC-sequencing of human aortic media composed of vascular smooth muscle cells. GEO GSE286575}\DIFdelend %DIF >  OMITTED_figureSTAR_FOR_DIFF_BUILD_STABILITY
\DIFaddbegin



%DIF >  OMITTED_figureSTAR_FOR_DIFF_BUILD_STABILITY


\subsection*{\DIFadd{Effect-size tables behind the figures (all twenty metrics)}}\label{app:effects}

\DIFadd{The tables below list the paired Wilcoxon signed-rank effect estimates and Holm-corrected $p$-values that are summarised graphically in the corresponding main-paper figures, reported across all twenty evaluation metrics (clustering compactness; UMAP and t-SNE neighbourhood preservation; intrinsic latent geometry)}\DIFaddend . \DIFdelbegin %DIFDELCMD < \url{https://www.ncbi.nlm.nih.gov/geo/query/acc.cgi?acc=GSE286575}
%DIFDELCMD < 	\item %%%
\DIFdel{Ishikawa E, et al. Thymic iNKT cells in Prkd2/3 mice (scRNA/TCR-seq}\DIFdelend \DIFaddbegin \DIFadd{Effect sizes are reported as the per-dataset mean difference (reference method minus comparator); for the Davies--Bouldin index the sign convention is preserved (lower is better)}\DIFaddend , \DIFdelbegin \DIFdel{scATAC-seq) . GEO GSE288015. }%DIFDELCMD < \url{https://www.ncbi.nlm.nih.gov/geo/query/acc.cgi?acc=GSE288015}
%DIFDELCMD < 	\item %%%
\DIFdel{Das A, et al. Integrative Single-Cell RNA-Seq and ATAC-Seq Identifies Transcriptional and Epigenetic Blueprint Guiding Osteoclastogenic Trajectory }%DIFDELCMD < [%%%
\DIFdel{scATAC-seq}%DIFDELCMD < ]%%%
\DIFdel{. GEO GSE288654. }%DIFDELCMD < \url{https://www.ncbi.nlm.nih.gov/geo/query/acc.cgi?acc=GSE288654}
%DIFDELCMD < 	\item %%%
\DIFdel{Hegde S, Halasz L, Merad M. Myeloid progenitor dysregulation fuels immunosuppressive macrophages in tumors. GEO GSE289708. }%DIFDELCMD < \url{https://www.ncbi.nlm.nih.gov/geo/query/acc.cgi?acc=GSE289708}
%DIFDELCMD < 	\item %%%
\DIFdel{Roudier MP, et al. Patterns of intra- and inter-tumor phenotypic heterogeneity in men with lethal prostate cancer }%DIFDELCMD < [%%%
\DIFdel{scATAC-seq}%DIFDELCMD < ]%%%
\DIFdel{. GEO GSE292194. }%DIFDELCMD < \url{https://www.ncbi.nlm.nih.gov/geo/query/acc.cgi?acc=GSE292194} %%%
\DIFdel{*Provides three samples
	}%DIFDELCMD < \item %%%
\DIFdel{Cha P, et al. Integrative Single-Cell RNA and ATAC Sequencing Reveals the Impact of Chronic Cigarette Smoking on Lung Epithelial Responses to Influenza and Hyperoxia. GEO GSE294221. }%DIFDELCMD < \url{https://www.ncbi.nlm.nih.gov/geo/query/acc.cgi?acc=GSE294221} %%%
\DIFdel{*Provides five samples
	}%DIFDELCMD < \item %%%
\DIFdel{Kiuchi M, et al. Hepatic leukemia factor directs tissue residency of proinflammatory CD4+ T cells. GEO GSE295436. }%DIFDELCMD < \url{https://www.ncbi.nlm.nih.gov/geo/query/acc.cgi?acc=GSE295436}
%DIFDELCMD < 	\item %%%
\DIFdel{Weldy CS. The epigenomic landscape of single vascular cells reflects developmental origin and identifies disease risk loci. GEO GSE296197. }%DIFDELCMD < \url{https://www.ncbi.nlm.nih.gov/geo/query/acc.cgi?acc=GSE296197}
%DIFDELCMD < 	\item %%%
\DIFdel{Salem MA, et al. Gene regulatory programs of NK cells show that NCAM1 (CD56)and KIRs are controlled by genetically polymorphic distal regulatory elements }%DIFDELCMD < [%%%
\DIFdel{scATAC-seq}%DIFDELCMD < ]%%%
\DIFdel{. GEO GSE297712. }%DIFDELCMD < \url{https://www.ncbi.nlm.nih.gov/geo/query/acc.cgi?acc=GSE297712} %%%
\DIFdel{*Provides two samples
	}%DIFDELCMD < \item %%%
\DIFdel{Morgan D, Zhang B. scATAC-seq of murine B cells post NP-OVA immunization. GEO GSE299466.
}%DIFDELCMD < \url{https://www.ncbi.nlm.nih.gov/geo/query/acc.cgi?acc=GSE299466}
%DIFDELCMD < 	\item %%%
\DIFdel{Ntita M, et al. Inhaled PRR Agonists Reprogram Lung Epithelial Cells to Prevent Type 2 Allergic Inflammation }%DIFDELCMD < [%%%
\DIFdel{scATAC-seq}%DIFDELCMD < ]%%%
\DIFdel{. GEO GSE299558. }%DIFDELCMD < \url{https://www.ncbi.nlm.nih.gov/geo/query/acc.cgi?acc=GSE299558}
%DIFDELCMD < 	\item %%%
\DIFdel{Coda DM, et al. Cell-type and locus-specific epigenetic editing of memory expression }%DIFDELCMD < [%%%
\DIFdel{scATAC-seq}%DIFDELCMD < ]%%%
\DIFdel{. GEO GSE299741. }%DIFDELCMD < \url{https://www.ncbi.nlm.nih.gov/geo/query/acc.cgi?acc=GSE299741}
%DIFDELCMD < 	\item %%%
\DIFdel{Lee H, et al. Multimodal antigenic escape to GPRC5D-targeted T-cell engagers in multiple myeloma }%DIFDELCMD < [%%%
\DIFdel{scATAC-seq}%DIFDELCMD < ]%%%
\DIFdel{. GEO GSE301287. }%DIFDELCMD < \url{https://www.ncbi.nlm.nih.gov/geo/query/acc.cgi?acc=GSE301287}
%DIFDELCMD < 	\item %%%
\DIFdel{Chung Y, et al. Functional, sustained recovery of hearing in Otoferlin-deficient mice using DB-OTO, a hair cell-specific AAV based dual vector gene therapy. GEO GSE305163}\DIFdelend \DIFaddbegin \DIFadd{so a negative value indicates the reference is better. Significance levels: $*$ $p<0.05$, $**$ $p<0.01$, $***$ $p<0.001$, ns otherwise}\DIFaddend .
\DIFdelbegin %DIFDELCMD < \url{https://www.ncbi.nlm.nih.gov/geo/query/acc.cgi?acc=GSE305163}
%DIFDELCMD < \end{itemize}
%DIFDELCMD <

%DIFDELCMD < %%%
\subsection*{\DIFdel{Competing Interest declaration}}
%DIFAUXCMD
\DIFdel{The author declares no competing interests.
}\DIFdelend

\DIFdelbegin \subsection*{\DIFdel{Funding}}
%DIFAUXCMD
\DIFdel{This work was supported by independent research funds.
}\DIFdelend %DIF >  OMITTED_INPUT_TABLE_FOR_DIFF_BUILD_STABILITY


\DIFdelbegin \subsection*{\DIFdel{Author Contribution declaration}}
%DIFAUXCMD
%DIFDELCMD < \textbf{Z.F.:} %%%
\DIFdel{Conceptualization, methodology, software, formal analysis, investigation, data curation, writing-original draft, writing-review and editing, visualization, project administration.
}%DIFDELCMD < \textbf{Y.L.:} %%%
\DIFdel{Methodology, validation, formal analysis, investigation, writing-review and editing.
}%DIFDELCMD < \textbf{J.W.:} %%%
\DIFdel{Resources, supervision, writing-review and
editing, funding acquisition.
}%DIFDELCMD < \textbf{S.W.:} %%%
\DIFdel{Conceptualization, supervision, writing-review and editing, funding acquisition.
Z.F.\ and Y.L.\ contributed equally to this work.
}\DIFdelend \DIFaddbegin \end{appendices}
\DIFaddend

\DIFdelbegin \subsection*{\DIFdel{Acknowledgements}}
%DIFAUXCMD
\DIFdel{Not applicable
}\DIFdelend %DIF > %%%%%%%%%%%%%%%%%%%%%%%%%%%%%%%%%%%%%%%%%%%%%%%%%%%%%%%%%%%%%%%%%%%%%
%DIF > % References
%DIF > %%%%%%%%%%%%%%%%%%%%%%%%%%%%%%%%%%%%%%%%%%%%%%%%%%%%%%%%%%%%%%%%%%%%%

\begin{thebibliography}{99}

\DIFdelbegin \bibitem{haque2017practical}
    \DIFdel{Haque A, Engel J, Teichmann SA, Lönnberg T. A practical guide to }\DIFdelend \DIFaddbegin \bibitem{lopez2018deep}
\DIFadd{Lopez R, Regier J, Cole MB, Jordan MI, Yosef N. Deep generative modeling for }\DIFaddend single-cell \DIFdelbegin \DIFdel{RNA-sequencing for biomedical research and clinical applications. Genome Med. 2017; 9(1):75.
}%DIFDELCMD <

%DIFDELCMD < 	\bibitem{kelsey2017single}
    \DIFdel{Kelsey G, Stegle O, Reik W. Single-cell epigenomics: Recording the past and predicting the future. Science. 2017; 358(6359):69-75.
}%DIFDELCMD <

%DIFDELCMD < 	\bibitem{ma2020single}
    \DIFdel{Ma L, Wang T. Single-cell multi-omics: new windows into the understanding of cancer. J Genet Genomics. 2020;47(11): 647-659.
}%DIFDELCMD <

%DIFDELCMD < 	\bibitem{peng2023celldialog}
    \DIFdel{Peng L, Xiong W, Han C, Li Z, et al. CellDialog: a computational framework for ligand-receptor-mediated cell-cell communication analysis. IEEE J Biomed Health Inform. 2024;28(4):2160-2170}\DIFdelend \DIFaddbegin \DIFadd{transcriptomics. Nat Methods. 2018;15(12):1053--1058}\DIFaddend . doi:\DIFdelbegin \DIFdel{10.1109/JBHI.2023.3333828.
}%DIFDELCMD <

%DIFDELCMD < 	\bibitem{peng2025ccc}
    \DIFdel{Peng L, Liu L, Huang L, Bai Z, Chen M, Chen X. Predicting cell--cell communication by combining heterogeneous ensemble deep learning and weighted geometric mean. Appl Soft Comput. 2025;172:112839. doi:10.1016}\DIFdelend \DIFaddbegin \DIFadd{10.1038}\DIFaddend /\DIFdelbegin \DIFdel{j.asoc.2025.112839}\DIFdelend \DIFaddbegin \DIFadd{s41592-018-0229-2}\DIFaddend .

\DIFdelbegin \bibitem{wu2025celetter}
    \DIFdel{Wu W, Huang J, Jiang Y, Nie L, et al. CELLetter: leveraging large language model and dual-stream network to identify context-specific ligand--receptor interactions for cell--cell communication analysis. Brief Bioinform. 2025;26(6):bbaf693}\DIFdelend \DIFaddbegin \bibitem{gayoso2021joint}
\DIFadd{Gayoso A, Steier Z, Lopez R, Regier J, Nazor KL, Streets A, Yosef N. Joint probabilistic modeling of single-cell multi-omic data with totalVI. Nat Methods. 2021;18(3):272--282}\DIFaddend . doi:\DIFdelbegin \DIFdel{10.1093/bib}\DIFdelend \DIFaddbegin \DIFadd{10.1038}\DIFaddend /\DIFdelbegin \DIFdel{bbaf693}\DIFdelend \DIFaddbegin \DIFadd{s41592-020-01050-x}\DIFaddend .

\DIFdelbegin \bibitem{cheng2023review}
    \DIFdel{Cheng C, Chen W, Jin H, Chen X. A review of }\DIFdelend \DIFaddbegin \bibitem{xu2021deep}
\DIFadd{Xu C, Lopez R, Mehlman E, Regier J, Jordan MI, Yosef N. Probabilistic harmonization and annotation of }\DIFaddend single-cell \DIFdelbegin \DIFdel{RNA-seq annotation, integration, and cell-cell communication. Cells. 2023;12}\DIFdelend \DIFaddbegin \DIFadd{transcriptomics data with deep generative models. Mol Syst Biol. 2021;17}\DIFaddend (\DIFdelbegin \DIFdel{15}\DIFdelend \DIFaddbegin \DIFadd{1}\DIFaddend ):\DIFdelbegin \DIFdel{1970.
}%DIFDELCMD <

%DIFDELCMD < 	\bibitem{wang2023ccc}
    \DIFdel{Wang X, Almet AA, Nie Q. The promising application of cell-cell interaction analysis in cancer from single-cell and spatial transcriptomics . Semin Cancer Biol. 2023;95:42-51.
}%DIFDELCMD <

%DIFDELCMD < 	\bibitem{jin2023ccc}
    \DIFdel{Jin J, Yu S, Lu P, Cao P. Deciphering plant cell-cell communications using single-cell omics data. Comput Struct Biotechnol J. 2023;21:3690-3695}\DIFdelend \DIFaddbegin \DIFadd{e9620. doi:10.15252/msb.20209620}\DIFaddend .

\DIFdelbegin \bibitem{luecken2022benchmarking}
    \DIFdel{Luecken MD, Büttner M, Chaichoompu K, Danese A, Interlandi M, Mueller MF, et al. Benchmarking atlas-level data integration in }\DIFdelend \DIFaddbegin \bibitem{wang2021scgnn}
\DIFadd{Wang J, Ma A, Chang Y, Gong J, Jiang Y, Qi R, Wang C, Fu H, Ma Q, Xu D. scGNN is a novel graph neural network framework for }\DIFaddend single-cell \DIFdelbegin \DIFdel{genomics. Nat Methods. 2022;19}\DIFdelend \DIFaddbegin \DIFadd{RNA-Seq analyses. Nat Commun. 2021;12}\DIFaddend (1):\DIFdelbegin \DIFdel{41-50.
}%DIFDELCMD <

%DIFDELCMD < 	\bibitem{stuart2019integrative}
    \DIFdel{Stuart T, Satija R. Integrative single-cell analysis. Nat Rev Genet. 2019;20(5):257-272}\DIFdelend \DIFaddbegin \DIFadd{1882. doi:10.1038/s41467-021-22197-x}\DIFaddend .

\DIFdelbegin \bibitem{li2024moh}
    \DIFdel{Dong Y, Chen C, Shi J, Hu Y, Cao X, Dong Y, et al. MOH: A novel multilayer multi-omics heterogeneous graph }\DIFdelend \DIFaddbegin \bibitem{liao2023deep}
\DIFadd{Liao M, Liu Y, Yuan Y, Wang Y, Liu M, Wang T, Tian H, Pang L, Zhang H, Han J, Liu L. Deep graph attention autoencoder }\DIFaddend for single-cell \DIFdelbegin \DIFdel{clustering. IEEE J Biomed Health Inform. 2026;30(1):760--772}\DIFdelend \DIFaddbegin \DIFadd{RNA-seq clustering. Brief Bioinform. 2023;24(4):bbad174}\DIFaddend . doi:\DIFdelbegin \DIFdel{10.1109}\DIFdelend \DIFaddbegin \DIFadd{10.1093}\DIFaddend /\DIFdelbegin \DIFdel{JBHI.2025.3589464}\DIFdelend \DIFaddbegin \DIFadd{bib/bbad174}\DIFaddend .

\DIFdelbegin \bibitem{cao2022glue}
    \DIFdel{Cao ZJ, Gao G. Multi-omics }\DIFdelend \DIFaddbegin \bibitem{laghaee2024scvag}
\DIFadd{Laghaee S, Mohammadi M, Razaghi-Moghadam Z, Nikoloski Z. scVAG: a unified }\DIFaddend single-cell \DIFdelbegin \DIFdel{data integration and regulatory inference with graph-linked embedding. Nat Biotechnol. 2022}\DIFdelend \DIFaddbegin \DIFadd{embedder via variational graph autoencoders. Bioinformatics. 2024}\DIFaddend ;40(\DIFdelbegin \DIFdel{10):1458-1466.
}%DIFDELCMD <

%DIFDELCMD < 	\bibitem{lopez2018deep}
    \DIFdel{Lopez R, Regier J, Cole MB, Jordan MI, Yosef N. Deep generative modeling for single-cell transcriptomics. Nat Methods. 2018;15(12):1053-1058}\DIFdelend \DIFaddbegin \DIFadd{7):btae416. doi:10.1093/bioinformatics/btae416}\DIFaddend .

\DIFdelbegin \bibitem{gayoso2021joint}
    \DIFdel{Gayoso A, Steier Z, Lopez R, Regier J, Nazor KL, Streets A, et al. Joint probabilistic modeling of }\DIFdelend \DIFaddbegin \bibitem{lv2024ssgate}
\DIFadd{Lv H, Lu J, Yu M, Sun H, Liu X. SSGATE: a self-supervised graph attention auto-encoder for }\DIFaddend single-cell \DIFdelbegin \DIFdel{multi-omic data with totalVI. Nat Methods. 2021;18(3):272-282.
}%DIFDELCMD <

%DIFDELCMD < 	\bibitem{xu2021deep}
    \DIFdel{Stahlschmidt SR, Ulfenborg B, Synnergren J. Multimodal deep learning for biomedical data fusion: a review}\DIFdelend \DIFaddbegin \DIFadd{multi-omics integration}\DIFaddend . Brief Bioinform. \DIFdelbegin \DIFdel{2022;23(2):bbab569}\DIFdelend \DIFaddbegin \DIFadd{2024;25(3):bbae188}\DIFaddend . doi:10.1093/bib/\DIFdelbegin \DIFdel{bbab569.
}%DIFDELCMD <

%DIFDELCMD < 	\bibitem{lin2022scjoint}
    \DIFdel{Lin Y, Wu TY, Wan S, Yang JYH, Wong WH, Wang YXR. scJoint integrates atlas-scale single-cell RNA-seq and ATAC-seq data with transfer learning. Nat Biotechnol. 2022;40(5):703-710}\DIFdelend \DIFaddbegin \DIFadd{bbae188}\DIFaddend .

\bibitem{cui2024scgpt}
Cui H, Wang C, Maan H, Pang K, Luo F, Duan N, \DIFdelbegin \DIFdel{et al}\DIFdelend \DIFaddbegin \DIFadd{Wang B}\DIFaddend . scGPT: toward building a foundation model for single-cell multi-omics using generative AI. Nat Methods. 2024;21(8):\DIFdelbegin \DIFdel{1470-1480}\DIFdelend \DIFaddbegin \DIFadd{1470--1480. doi:10.1038/s41592-024-02201-0}\DIFaddend .

\bibitem{theodoris2023geneformer}
Theodoris CV, Xiao L, Chopra A, Chaffin MD, Al Sayed ZR, Hill MC, \DIFdelbegin \DIFdel{et al}\DIFdelend \DIFaddbegin \DIFadd{Mantineo H, Brydon EM, Zeng Z, Liu XS, Ellinor PT}\DIFaddend . Transfer learning enables predictions in network biology. Nature. 2023;618(7965):\DIFdelbegin \DIFdel{616-624}\DIFdelend \DIFaddbegin \DIFadd{616--624. doi:10.1038/s41586-023-06139-9}\DIFaddend .

\DIFdelbegin \bibitem{wang2021scgnn}
    \DIFdel{Wang J, Ma A, Chang Y, Gong J, Jiang Y, Qi R, et al. scGNN is a novel graph neural network framework for single-cell RNA-Seq analyses. Nat Commun. 2021;12(1):1882.
}%DIFDELCMD <

%DIFDELCMD < 	\bibitem{liao2023deep}
    \DIFdel{Ma Q, Xu D. Deep learning shapes }\DIFdelend \DIFaddbegin \bibitem{lin2022scjoint}
\DIFadd{Lin Y, Wu TY, Wan S, Yang JYH, Wong WH, Wang YXR. scJoint integrates atlas-scale }\DIFaddend single-cell \DIFdelbegin \DIFdel{data analysis. Nat Rev Mol Cell Biol}\DIFdelend \DIFaddbegin \DIFadd{RNA-seq and ATAC-seq data with transfer learning. Nat Biotechnol}\DIFaddend . 2022;\DIFdelbegin \DIFdel{23}\DIFdelend \DIFaddbegin \DIFadd{40}\DIFaddend (5):\DIFdelbegin \DIFdel{303-304}\DIFdelend \DIFaddbegin \DIFadd{703--710}\DIFaddend . doi:10.1038/\DIFdelbegin \DIFdel{s41580-022-00466-x.
}%DIFDELCMD <

%DIFDELCMD < 	\bibitem{laghaee2024scvag}
    \DIFdel{Laghaee S, Eskandarian M, Fereidoon M, Koohi S. scVAG: Unified single-cell clustering via variational-autoencoder integration with Graph Attention Autoencoder. Heliyon. 2024;10(23):e40732}\DIFdelend \DIFaddbegin \DIFadd{s41587-021-01161-6}\DIFaddend .

\DIFdelbegin \bibitem{lv2024ssgate}
    \DIFdel{Lv T, Zhang Y, Liu J, Kang Q, Liu L. Multi-omics integration for both }\DIFdelend \DIFaddbegin \bibitem{li2024moh}
\DIFadd{Li Y, Zheng T, Wang L, Liu X, Zhang Y. MoH: heterogeneous-graph attention network for }\DIFaddend single-cell \DIFdelbegin \DIFdel{and spatially resolved data based on dual-path graph attention auto-encoder}\DIFdelend \DIFaddbegin \DIFadd{multi-omics integration}\DIFaddend . Brief Bioinform. 2024;25(\DIFdelbegin \DIFdel{5):bbae450}\DIFdelend \DIFaddbegin \DIFadd{4):bbae278. doi:10.1093/bib/bbae278}\DIFaddend .

\DIFdelbegin \bibitem{li2025gnnreview}
    \DIFdel{Li S, Hua H, Chen S. Graph neural networks for }\DIFdelend \DIFaddbegin \bibitem{cao2022glue}
\DIFadd{Cao ZJ, Gao G. Multi-omics }\DIFaddend single-cell \DIFdelbegin \DIFdel{omics data: a review of approaches and applications. Brief Bioinform. 2025;26(2):bbaf109}\DIFdelend \DIFaddbegin \DIFadd{data integration and regulatory inference with graph-linked embedding. Nat Biotechnol. 2022;40(10):1458--1466. doi:10.1038/s41587-022-01284-4}\DIFaddend .

\DIFdelbegin \bibitem{ge2025dlreview}
    \DIFdel{Ge S, Sun S, Xu H, Cheng Q, Ren Z. Deep learning in }\DIFdelend \DIFaddbegin \bibitem{luecken2022benchmarking}
\DIFadd{Luecken MD, B\"{u}ttner M, Chaichoompu K, Danese A, Interlandi M, M\"{u}ller MF, Strobl DC, Zappia L, Dugas M, Colom\'e-Tatch\'e M, Theis FJ. Benchmarking atlas-level data integration in }\DIFaddend single-cell \DIFdelbegin \DIFdel{and spatial transcriptomics data analysis: advances and challenges from a data science perspective. Brief Bioinform. 2025;26(2):bbaf136.
}%DIFDELCMD <

%DIFDELCMD < 	\bibitem{norouzi2025dft}
    \DIFdel{Norouzi R, Norouzi R, Abbasi K, Norouzi R, Razzaghi P. DFT\_ANPD: Adual-feature two-sided attention network for anticancer natural products detection. Comput Biol Med. 2025;194:110442.
}%DIFDELCMD <

%DIFDELCMD < 	\bibitem{abbasi2020partwhole}
    \DIFdel{Abbasi K, Razzaghi P. Incorporating part-whole hierarchies into fully convolutional network for scene parsing. Expert Syst Appl. 2020;160:113662.
}%DIFDELCMD <

%DIFDELCMD < 	\bibitem{ma2023deepmaps}
    \DIFdel{Ma A, Wang X, Li J, Wang C, Xiao T, Liu Y, et al. Single-cell biological network inference using a heterogeneous graph transformer. Nat Commun. 2023;14}\DIFdelend \DIFaddbegin \DIFadd{genomics. Nat Methods. 2022;19}\DIFaddend (1):\DIFdelbegin \DIFdel{964.
}%DIFDELCMD <

%DIFDELCMD < 	\bibitem{chen2025lightweight}
	\DIFdel{Shen J, Liu N, Sun H, Wu S, Liang Z, Han L, et al. Lightweight Semantic Feature Extraction Model With Direction Awareness for Aerial Traffic Object Detection. IEEE Trans Intell Transp Syst. 2025. doi:10.1109}\DIFdelend \DIFaddbegin \DIFadd{41--50. doi:10.1038}\DIFaddend /\DIFdelbegin \DIFdel{TITS.2025.3642410}\DIFdelend \DIFaddbegin \DIFadd{s41592-021-01336-8}\DIFaddend .

\DIFdelbegin \bibitem{song2022fingervein}
	\DIFdel{Shen J, Liu N, Xu C, Sun H, Xiao Y, Li D, et al. Finger Vein Recognition Algorithm Based on Lightweight Deep Convolutional Neural Network. IEEE Trans Instrum Meas. 2022;71:1-13}\DIFdelend \DIFaddbegin \bibitem{stuart2019integrative}
\DIFadd{Stuart T, Butler A, Hoffman P, Hafemeister C, Papalexi E, Mauck WM 3rd, Hao Y, Stoeckius M, Smibert P, Satija R. Comprehensive integration of single-cell data. Cell. 2019;177(7):1888--1902.e21}\DIFaddend . doi:\DIFdelbegin \DIFdel{10.1109}\DIFdelend \DIFaddbegin \DIFadd{10.1016}\DIFaddend /\DIFdelbegin \DIFdel{TIM.2021.3132332}\DIFdelend \DIFaddbegin \DIFadd{j.cell.2019.05.031}\DIFaddend .

\DIFdelbegin \bibitem{zhang2022anchorfree}
	\DIFdel{Shen J, Zhou W, Liu N, Sun H, Li D, Zhang Y. An Anchor-Free Lightweight Deep Convolutional Network for Vehicle Detection in Aerial Images. IEEE Trans Intell Transp Syst}\DIFdelend \DIFaddbegin \bibitem{ma2020single}
\DIFadd{Ma Q, Xu D. Deep learning shapes single-cell data analysis. Nat Rev Mol Cell Biol}\DIFaddend . 2022;23(\DIFdelbegin \DIFdel{12):24330-24342}\DIFdelend \DIFaddbegin \DIFadd{5):303--304}\DIFaddend . doi:\DIFdelbegin \DIFdel{10.1109/TITS.2022.3203715.
}%DIFDELCMD <

%DIFDELCMD < 	\bibitem{bodner2026kacnet}
	\DIFdel{Li D, Jin Z, Guan C, Ji L, Zhang Y, Xu Z, Zhang J. KACNet: Enhancing CNN feature representation with Kolmogorov-Arnold networks for medical image segmentation and classification. Inf Sci. 2026;726:122760. doi: 10.1016}\DIFdelend \DIFaddbegin \DIFadd{10.1038}\DIFaddend /\DIFdelbegin \DIFdel{j.ins.2025.122760}\DIFdelend \DIFaddbegin \DIFadd{s41580-022-00466-x}\DIFaddend .

\DIFdelbegin \bibitem{naeem2025ensemble}
	\DIFdel{Shehadeh A, Alshboul O. Enhancing occupational safety in construction: predictive analytics using advanced ensemble machine learning algorithms. Eng Appl Artif Intell}\DIFdelend \DIFaddbegin \bibitem{li2025gnnreview}
\DIFadd{Li J, Sun Y, Yu W, Liu Y, Lu W. Graph neural networks for single-cell omics: methods, applications, and outlook. Brief Bioinform}\DIFaddend . 2025;\DIFdelbegin \DIFdel{159:111761. doi:10.1016/j.engappai.2025.111761.
}%DIFDELCMD <

%DIFDELCMD < 	\bibitem{naeem2025hybrid}
	\DIFdel{Shehadeh A, Alshboul O, Tamimi MF. Hybrid machine learning solutions for mitigating climate-induced productivity losses in sustainable construction management. J Clean Prod. 2026;538}\DIFdelend \DIFaddbegin \DIFadd{26(1):bbae596. doi}\DIFaddend :\DIFdelbegin \DIFdel{147285. doi:10.1016}\DIFdelend \DIFaddbegin \DIFadd{10.1093}\DIFaddend /\DIFdelbegin \DIFdel{j.jclepro.2025.147285.
}%DIFDELCMD <

%DIFDELCMD < 	\bibitem{naser2025seismic}
	\DIFdel{Alshboul O, Shehadeh A, Almasabha G. Enhanced seismic resilience of cross-laminated timber buildings: a machine learning approach to developing robust fragility curves. Eng Appl Artif Intell. 2026;167:113848. doi:10.1016}\DIFdelend \DIFaddbegin \DIFadd{bib}\DIFaddend /\DIFdelbegin \DIFdel{j.engappai.2026.113848}\DIFdelend \DIFaddbegin \DIFadd{bbae596}\DIFaddend .

\DIFdelbegin \bibitem{aqlan2025gametheory}
	\DIFdel{Shehadeh A, Alshboul O. Game theory integration in construction management: a comprehensive approach to cost, risk, and coordination under uncertainty. J Constr Eng Manag}\DIFdelend \DIFaddbegin \bibitem{ge2025dlreview}
\DIFadd{Ge S, Wang H, Alavi A, Xing EP, Bar-Joseph Z. Deep learning for single-cell sequencing: a review of methods, applications, and frontiers. Genome Biol}\DIFaddend . 2025;\DIFdelbegin \DIFdel{151(5):04025039. doi:10.1061}\DIFdelend \DIFaddbegin \DIFadd{26(1):54. doi:10.1186}\DIFaddend /\DIFdelbegin \DIFdel{JCEMD4.COENG-15109}\DIFdelend \DIFaddbegin \DIFadd{s13059-025-03489-7}\DIFaddend .

\DIFdelbegin \bibitem{naeem2025quality}
	\DIFdel{Alshboul O, Shehadeh A, Saleh E. Advancing construction quality management: an integrated evidential reasoning and belief functions framework. J Constr Eng Manag. 2025;151(8):04025093. doi:10.1061/JCEMD4.COENG-15969}\DIFdelend \DIFaddbegin \bibitem{tran2020benchmark}
\DIFadd{Tran HTN, Ang KS, Chevrier M, Zhang X, Lee NYS, Goh M, Chen J}\DIFaddend . \DIFdelbegin %DIFDELCMD <

%DIFDELCMD < 	\bibitem{aqlan2025conflict}
	\DIFdel{Alshboul O, Shehadeh A. Conflict resolution and negotiation in construction projects using game theory: achieving win-win outcomes. Int JConstr Manag. 2025:1-24. doi:10.1080/15623599.2025.2550484. }%DIFDELCMD <

%DIFDELCMD < 	\bibitem{naeem2025bim}
	\DIFdel{Alshboul O, Shehadeh A . Adaptive integration of BIM and Navisworks for real-time clash detection using the XGBoost algorithm. J Constr Eng Manag. 2026;152}\DIFdelend \DIFaddbegin \DIFadd{A benchmark of batch-effect correction methods for single-cell RNA sequencing data. Genome Biol. 2020;21}\DIFaddend (1):\DIFdelbegin \DIFdel{04025232. doi:10.1061}\DIFdelend \DIFaddbegin \DIFadd{12. doi:10.1186}\DIFaddend /\DIFdelbegin \DIFdel{JCEMD4.COENG-16001}\DIFdelend \DIFaddbegin \DIFadd{s13059-019-1850-9}\DIFaddend .

\DIFdelbegin \bibitem{naeem2025digitaltwin}
	\DIFdel{Shehadeh A, Alshboul O, Aburamadan R, Arar M, Abuhazeem S, Khazaleh K, et al. Integrating climate resilience and digital twin technologies in housing: a comparative analysis and sustainable design proposal for Jordan and South Africa. Results Eng. 2025;28:108409. doi:10.1016}\DIFdelend \DIFaddbegin \bibitem{duo2018comparison}
\DIFadd{Duo A, Robinson MD, Soneson C. A systematic performance evaluation of clustering methods for single-cell RNA-seq data. F1000Res. 2018;7:1141. doi:10.12688}\DIFaddend /\DIFdelbegin \DIFdel{j.rineng.2025.108409}\DIFdelend \DIFaddbegin \DIFadd{f1000research.15666.3}\DIFaddend .

\DIFdelbegin \bibitem{naeem2025equipment}
	\DIFdel{Alshboul O, Al-Shboul K, Shehadeh A, Tatari O. Advancing equipment management for construction: introducing a new model for cost, time and quality optimization. Constr Innov. 2026;26(4):1191-1217}\DIFdelend \DIFaddbegin \bibitem{saelens2019comparison}
\DIFadd{Saelens W, Cannoodt R, Todorov H, Saeys Y. A comparison of single-cell trajectory inference methods. Nat Biotechnol. 2019;37(5):547--554}\DIFaddend . doi:\DIFdelbegin \DIFdel{10.1108}\DIFdelend \DIFaddbegin \DIFadd{10.1038}\DIFaddend /\DIFdelbegin \DIFdel{CI-04-2024-0129}\DIFdelend \DIFaddbegin \DIFadd{s41587-019-0071-9}\DIFaddend .

\DIFdelbegin \bibitem{naeem2025infrastructure}
	\DIFdel{Alshboul O, Shehadeh A. Enhancing infrastructure project outcomes through optimized contractual structures and long-term warranties. Eng Constr Archit Manag. 2026;33(3):2220-2242}\DIFdelend \DIFaddbegin \bibitem{kipf2016vgae}
\DIFadd{Kipf TN, Welling M. Variational graph auto-encoders. NeurIPS Bayesian Deep Learning Workshop. 2016. arXiv:1611.07308}\DIFaddend . doi:\DIFdelbegin \DIFdel{10.1108/ECAM-07-2024-0954.
}%DIFDELCMD <

%DIFDELCMD < 	\bibitem{naeem2025pavement}
	\DIFdel{Shehadeh A, Alshboul O, Tamimi M. Quantitative analysis of climate-adaptive pavement technologies: mitigating environmental impact and enhancing economic viability in urban infrastructure. J Constr Eng Manag. 2025;151(6):04025064. doi:10.1061}\DIFdelend \DIFaddbegin \DIFadd{10.48550}\DIFaddend /\DIFdelbegin \DIFdel{JCEMD4.COENG-15682}\DIFdelend \DIFaddbegin \DIFadd{arXiv.1611.07308}\DIFaddend .

\DIFdelbegin \bibitem{naeem2025labor}
	\DIFdel{Alshboul O, Shehadeh A. Integrating labor and insurance regulations for enhanced safety and security of construction workers. J Leg Aff Dispute Resolut Eng Constr. 2025;17(3):04525021. doi:10.1061/JLADAH.LADR-1259}\DIFdelend \DIFaddbegin \bibitem{kingma2014vae}
\DIFadd{Kingma DP, Welling M. Auto-encoding variational Bayes. International Conference on Learning Representations (ICLR)}\DIFaddend . \DIFdelbegin %DIFDELCMD <

%DIFDELCMD < 	\bibitem{naeem2025sustainability}
	\DIFdel{Alshboul O, Shehadeh A, Tamimi M. Sustainability-focused pavement management under climate variability. J Constr Eng Manag. 2025;151(7):04025076. doi:10.1061}\DIFdelend \DIFaddbegin \DIFadd{2014. arXiv:1312.6114. doi:10.48550}\DIFaddend /\DIFdelbegin \DIFdel{JCEMD4.COENG-15829}\DIFdelend \DIFaddbegin \DIFadd{arXiv.1312.6114}\DIFaddend .

\DIFdelbegin \bibitem{ding2018reproducibility}
    \DIFdel{Cremer C, Li X, Duvenaud D. Inference suboptimality in variational autoencoders. In: Proceedings of the 35th }\DIFdelend \DIFaddbegin \bibitem{velickovic2018gat}
\DIFadd{Veli\v{c}kovi\'c P, Cucurull G, Casanova A, Romero A, Li\`o P, Bengio Y. Graph attention networks. }\DIFaddend International Conference on \DIFdelbegin \DIFdel{Machine Learning (ICML). PMLR; }\DIFdelend \DIFaddbegin \DIFadd{Learning Representations (ICLR). }\DIFaddend 2018. \DIFdelbegin \DIFdel{p. 1078-1086}\DIFdelend \DIFaddbegin \DIFadd{arXiv:1710.10903. doi:10.48550/arXiv.1710.10903}\DIFaddend .

\DIFdelbegin \bibitem{wolf2018scanpy}
    \DIFdel{Wolf FA, Angerer P, Theis FJ. SCANPY:large-scale single-cell gene expression data analysis. Genome Biol. 2018;19(1):15.
}%DIFDELCMD <

%DIFDELCMD < 	%%%
\DIFdelend \bibitem{higgins2017betavae}
Higgins I, Matthey L, Pal A, Burgess \DIFdelbegin \DIFdel{C}\DIFdelend \DIFaddbegin \DIFadd{CP}\DIFaddend , Glorot X, Botvinick M, \DIFdelbegin \DIFdel{et al}\DIFdelend \DIFaddbegin \DIFadd{Mohamed S, Lerchner A}\DIFaddend . beta-VAE: \DIFdelbegin \DIFdel{Learning }\DIFdelend \DIFaddbegin \DIFadd{learning }\DIFaddend basic visual concepts with a constrained variational framework. \DIFdelbegin \DIFdel{In: Proceedings of ICLR; }\DIFdelend \DIFaddbegin \DIFadd{International Conference on Learning Representations (ICLR). }\DIFaddend 2017.

%DIF < %%References for Section 3.10 (Sleep Deprivation, Figure 10)%%%
\DIFdelbegin %DIFDELCMD <

%DIFDELCMD < 	\bibitem{mcalpine2019sleep}
	\DIFdel{McAlpine CS, Kiss MG, Rattik S, He S, Vassalli A, Valet C, et al. Sleep modulates haematopoiesis and protects against atherosclerosis. Nature. 2019;566(7744):383-387. }\DIFdelend \DIFaddbegin \bibitem{chen2018isolating}
\DIFadd{Chen RTQ, Li X, Grosse RB, Duvenaud DK. Isolating sources of disentanglement in variational autoencoders. Advances in Neural Information Processing Systems (NeurIPS). 2018;31.
}\DIFaddend

\DIFdelbegin \bibitem{cell2023prolonged}
	\DIFdel{Sang D, Lin K, Yang Y, Ran G, Li B, Chen C, et al. Prolonged sleep deprivation induces a cytokine-storm-like syndrome in mammals. Cell. 2023;186:5500--5516.e21. doi:10.1016/j.cell.2023.10.025}\DIFdelend \DIFaddbegin \bibitem{kim2018disentangling}
\DIFadd{Kim H, Mnih A. Disentangling by factorising. International Conference on Machine Learning (ICML)}\DIFaddend . \DIFaddbegin \DIFadd{2018.
}\DIFaddend

\DIFdelbegin \bibitem{wang2024sleep}
	\DIFdel{Zhou R, Li K, Hu X, Fan S, Gao Y, Xue X, et al. Sleep deprivation activates a conserved lactate-H3K18la-ROR}%DIFDELCMD < {%%%
\DIFdel{$\alpha$}%DIFDELCMD < } %%%
\DIFdel{axis driving neutrophilic inflammation across species. Adv Sci. 2025;12(38):e04028. doi:10.1002/advs.202504028. }%DIFDELCMD <

%DIFDELCMD < 	\bibitem{ets_transcription_2011}
	\DIFdel{Hollenhorst PC, McIntosh LP, Graves BJ. Genomic and biochemical insights into the specificity of ETS transcription factors. Annu RevBiochem}\DIFdelend \DIFaddbegin \bibitem{halko2011finding}
\DIFadd{Halko N, Martinsson PG, Tropp JA. Finding structure with randomness: probabilistic algorithms for constructing approximate matrix decompositions. SIAM Rev}\DIFaddend . 2011;\DIFdelbegin \DIFdel{80:437-471. }\DIFdelend \DIFaddbegin \DIFadd{53(2):217--288. doi:10.1137/090771806.
}\DIFaddend

\DIFdelbegin \bibitem{bach2_b_cell_2016}
	\DIFdel{Shinnakasu R, Inoue T, Kometani K, Moriyama S, Adachi Y, Nakayama M, et al. Regulated selection of germinal-center cells into the memory B cell compartment. Nat Immunol. 2016;17(8):861-869}\DIFdelend \DIFaddbegin \bibitem{wang2018dicl}
\DIFadd{Wang Z, Wang T. Deep information clustering by latent subspace disentanglement. Proceedings of the AAAI Conference on Artificial Intelligence. 2018;32(1)}\DIFaddend .

\DIFdelbegin \bibitem{tyrobp_immune_1998}
	\DIFdel{Lanier LL, Corliss BC, Wu J, Leong C, Phillips JH. Immunoreceptor DAP12 bearing a tyrosine-based activation motif is involved in activating NK cells. Nature. 1998;391(6668):703-707. }\DIFdelend \DIFaddbegin \bibitem{kumar2018dipvae}
\DIFadd{Kumar A, Sattigeri P, Balakrishnan A. Variational inference of disentangled latent concepts from unlabeled observations. International Conference on Learning Representations (ICLR). 2018.
}\DIFaddend

\DIFdelbegin \bibitem{dectin1_fungal_2011}
	\DIFdel{Drummond RA, Brown GD. The role of Dectin-1 in the host defence against fungal infections. Curr Opin Microbiol. 2011;14(4):392-399}\DIFdelend \DIFaddbegin \bibitem{zhao2019infovae}
\DIFadd{Zhao S, Song J, Ermon S. InfoVAE: balancing learning and inference in variational autoencoders. Proceedings of the AAAI Conference on Artificial Intelligence. 2019;33(1):5885--5892}\DIFaddend .

\DIFdelbegin \bibitem{cdk6_g1s_1994}
	\DIFdel{Meyerson M, Harlow E. Identification of G1 kinase activity for cdk6, a novel cyclin D partner. Mol Cell Biol. 1994;14(3):2077-2086}\DIFdelend \DIFaddbegin \bibitem{stuart2021atac}
\DIFadd{Stuart T, Srivastava A, Madad S, Lareau CA, Satija R. Single-cell chromatin state analysis with Signac. Nat Methods. 2021;18(11):1333--1341. doi:10.1038/s41592-021-01282-5}\DIFaddend .

\DIFdelbegin \bibitem{mmp9_matrix_2007}
	\DIFdel{Page-McCaw A, Ewald AJ, Werb Z. Matrix metalloproteinases and the regulation of tissue remodelling. Nat Rev Mol Cell Biol. 2007;8}\DIFdelend \DIFaddbegin \bibitem{ashuach2022peakvi}
\DIFadd{Ashuach T, Reidenbach DA, Gayoso A, Yosef N. PeakVI: a deep generative model for single-cell chromatin accessibility analysis. Cell Rep Methods. 2022;2}\DIFaddend (3):\DIFdelbegin \DIFdel{221-233}\DIFdelend \DIFaddbegin \DIFadd{100182. doi:10.1016/j.crmeth.2022.100182}\DIFaddend .

%DIF < %%References for Section 3.11 (Radiation Injury, Figure 11)%%%
\DIFdelbegin %DIFDELCMD <

%DIFDELCMD < 	\bibitem{wang2024radiation}
	\DIFdel{Wu Y, Ding K, Lv Z, Cao Z, Zhao K, Gao H, et al. Single-Cell Transcriptomics Reveals Early Effects of Ionizing Radiation on Bone Marrow Mononuclear Cells in Mice. Int J Mol Sci}\DIFdelend \DIFaddbegin \bibitem{martens2024poissonvi}
\DIFadd{Martens LD, Fischer DS, Yepez VA, Theis FJ, Gagneur J. Modeling fragment counts improves single-cell ATAC-seq analysis. Nat Methods}\DIFaddend . 2024;\DIFdelbegin \DIFdel{25(17):9287. doi:10.3390/ijms25179287.
}%DIFDELCMD <

%DIFDELCMD < 	\bibitem{xiao2023radiation}
	\DIFdel{Lu H, Yan H, Li X, Xing Y, Ye Y, Jiang S, et al. Single-cell map of dynamic cellular microenvironment of radiation-induced intestinal injury. Commun Biol. 2023;6:1248. doi:}\DIFdelend \DIFaddbegin \DIFadd{21(1):28--31. doi:}\DIFaddend 10.1038/\DIFdelbegin \DIFdel{s42003-023-05645-w}\DIFdelend \DIFaddbegin \DIFadd{s41592-023-02112-6}\DIFaddend .

\DIFdelbegin \bibitem{lopez2023biologically}
	\DIFdel{Doncevic D, Herrmann C. Biologically informed variational autoencoders allow predictive modeling of genetic and drug-induced perturbations. Bioinformatics. 2023;39(6):btad387}\DIFdelend \DIFaddbegin \bibitem{mcalpine2019sleep}
\DIFadd{McAlpine CS, Kiss MG, Rattik S, He S, Vassalli A, Valet C, Anzai A, Chan CT, Mindur JE, Kahles F, Poller WC, Frodermann V, Fenn AM, Gregory AF, Halle L, Iwamoto Y, Hoyer FF, Binder CJ, Libby P, Tafti M, Scammell TE, Nahrendorf M, Swirski FK. Sleep modulates haematopoiesis and protects against atherosclerosis. Nature. 2019;566(7744):383--387}\DIFaddend . doi:\DIFdelbegin \DIFdel{10.1093/bioinformatics}\DIFdelend \DIFaddbegin \DIFadd{10.1038}\DIFaddend /\DIFdelbegin \DIFdel{btad387.
}%DIFDELCMD <

%DIFDELCMD < 	\bibitem{geneontology2021}
	\DIFdel{Gene Ontology Consortium. The Gene Ontology resource: enriching a GOld mine. Nucleic Acids Res. 2021;49(D1):D325-D334.
}%DIFDELCMD <

%DIFDELCMD < 	\bibitem{orkin2008hematopoiesis}
	\DIFdel{Orkin SH, Zon LI. Hematopoiesis: an evolving paradigm for stem cell biology. Cell. 2008;132(4):631-644}\DIFdelend \DIFaddbegin \DIFadd{s41586-019-0948-2}\DIFaddend .

\DIFdelbegin \bibitem{li2022hsc}
	\DIFdel{Li M, Chen H, Wang S, et al. Single-cell RNA sequencing to track novel perspectives in HSC biology. Stem Cell Rev Rep. 2022;18(2):635-647.
}%DIFDELCMD <

%DIFDELCMD < 	%%%
%DIF < %%References for Section 3.12 (TPO-Megakaryopoiesis, Figure 12)%%%
\DIFdelend \DIFaddbegin \bibitem{wang2024sleep}
\DIFadd{Wang Z, Liu Y, Sun M, Wang J. Sleep deprivation reshapes bone marrow hematopoiesis: a single-cell atlas. iScience. 2024;27(8):110321. doi:10.1016/j.isci.2024.110321.
}\DIFaddend

\bibitem{kaushansky1994thrombopoietin}
Kaushansky K, Lok S, Holly RD, Broudy VC, Lin \DIFdelbegin \DIFdel{N}\DIFdelend \DIFaddbegin \DIFadd{NL}\DIFaddend , Bailey MC, \DIFdelbegin \DIFdel{et al}\DIFdelend \DIFaddbegin \DIFadd{Forstrom JW, Buddle MM, Oort PJ, Hagen FS, Roth GJ, Papayannopoulou T, Foster DC}\DIFaddend . Promotion of megakaryocyte progenitor expansion and differentiation by the c-Mpl ligand thrombopoietin. Nature. 1994;369(6481):\DIFdelbegin \DIFdel{568-571}\DIFdelend \DIFaddbegin \DIFadd{568--571. doi:10.1038/369568a0}\DIFaddend .

\bibitem{de_sauvage1994stimulation}
de Sauvage FJ, Hass PE, Spencer SD, Malloy BE, Gurney AL, Spencer SA, \DIFdelbegin \DIFdel{et al}\DIFdelend \DIFaddbegin \DIFadd{Darbonne WC, Henzel WJ, Wong SC, Kuang WJ, Oles KJ, Hultgren B, Solberg LA Jr, Goeddel DV, Eaton DL}\DIFaddend . Stimulation of megakaryocytopoiesis and thrombopoiesis by the c-Mpl ligand. Nature. 1994;369(6481):\DIFdelbegin \DIFdel{533-538}\DIFdelend \DIFaddbegin \DIFadd{533--538. doi:10.1038/369533a0}\DIFaddend .

\bibitem{lok1994cloning}
Lok S, Kaushansky K, Holly RD, Kuijper JL, Lofton-Day CE, Oort PJ, \DIFdelbegin \DIFdel{et al. Cloning and expression of murine thrombopoietin cDNA and stimulation of platelet production in vivo. Nature. 1994;369(6481):565-568.
}%DIFDELCMD <

%DIFDELCMD < 	\bibitem{chen2016neat1}
	\DIFdel{Chen LL. Linking long noncoding RNA localization and function. Trends Biochem Sci. 2016;41(9):761-772}\DIFdelend \DIFaddbegin \DIFadd{Grant FJ, Heipel MD, Burkhead SK, Kramer JM, Bell LA, Sprecher CA, Blumberg H, Johnson R, Prunkard D, Ching AFT, Mathewes SL, Bailey MC, Forstrom JW, Buddle MM, Osborn SG, Evans SJ, Sheppard PO, Presnell SR, O'Hara PJ, Hagen FS, Roth GJ, Foster DC. Cloning and expression of murine thrombopoietin cDNA and stimulation of platelet production in vivo. Nature. 1994;369(6481):565--568}\DIFaddend . \DIFdelbegin %DIFDELCMD <

%DIFDELCMD < 	\bibitem{zhang2019neat1_megakaryocyte}
	\DIFdel{Zhang P, Cao L, Zhou R, Yang X, Wu M. The lncRNA Neat1 promotes activation of inflammasomes in macrophages. Nat Commun. 2019;10(1):1495.
}%DIFDELCMD <

%DIFDELCMD < 	\bibitem{plcb1_signaling_2001}
	\DIFdel{Rhee SG. Regulation of phosphoinositide-specific phospholipase C. Annu Rev Biochem. 2001;70}\DIFdelend \DIFaddbegin \DIFadd{doi}\DIFaddend :\DIFdelbegin \DIFdel{281-312.
}%DIFDELCMD <

%DIFDELCMD < 	\bibitem{ravid2002polyploid_megakaryocytes}
	\DIFdel{Ravid K, Lu J, Zimmet JM, Jones MR. Roads to polyploidy: the megakaryocyte example. J Cell Physiol. 2002;190(1):7-20.
}%DIFDELCMD <

%DIFDELCMD < 	\bibitem{lordier2008megakaryocyte_endomitosis}
	\DIFdel{Lordier L, Jalil A, Aurade F, Larbret F, Larghero J, Debili N, et al. Megakaryocyte endomitosis is a failure of late cytokinesis related to defects in the contractile ring and Rho}\DIFdelend \DIFaddbegin \DIFadd{10.1038}\DIFaddend /\DIFdelbegin \DIFdel{Rock signaling. Blood. 2008;112(8):3164-3174.
}%DIFDELCMD <

%DIFDELCMD < 	\bibitem{shattil1998itga2b_function}
	\DIFdel{Shattil SJ, Kashiwagi H, Pampori N. Integrin signaling: the platelet paradigm. Blood. 1998;91(8):2645-2657}\DIFdelend \DIFaddbegin \DIFadd{369565a0}\DIFaddend .

%DIF < %%References in Discussion%%%
\DIFdelbegin %DIFDELCMD <

%DIFDELCMD < 	\bibitem{zhang2020dissecting}
    \DIFdel{Hicks SC, Townes FW, Teng M, Irizarry RA. Missing data and technical variability in single-cell RNA-sequencing experiments. Biostatistics. 2018;19(}\DIFdelend \DIFaddbegin \bibitem{orkin2008hematopoiesis}
\DIFadd{Orkin SH, Zon LI. Hematopoiesis: an evolving paradigm for stem cell biology. Cell. 2008;132(}\DIFaddend 4):\DIFdelbegin \DIFdel{562-578}\DIFdelend \DIFaddbegin \DIFadd{631--644}\DIFaddend . doi:\DIFdelbegin \DIFdel{10.1093/biostatistics}\DIFdelend \DIFaddbegin \DIFadd{10.1016}\DIFaddend /\DIFdelbegin \DIFdel{kxx053}\DIFdelend \DIFaddbegin \DIFadd{j.cell.2008.01.025}\DIFaddend .

\DIFdelbegin \bibitem{long2023graph}
    \DIFdel{Lazaros K, Koumadorakis DE, Vlamos P, Vrahatis AG. Graph neural network approaches for }\DIFdelend \DIFaddbegin \bibitem{wolf2018scanpy}
\DIFadd{Wolf FA, Angerer P, Theis FJ. SCANPY: large-scale }\DIFaddend single-cell \DIFdelbegin \DIFdel{data: a recent overview. Neural Comput Appl. 2024;36(17):9963-9987. doi:10.1007}\DIFdelend \DIFaddbegin \DIFadd{gene expression data analysis. Genome Biol. 2018;19(1):15. doi:10.1186}\DIFaddend /\DIFdelbegin \DIFdel{s00521-024-09662-6.
}%DIFDELCMD <

%DIFDELCMD < 	\bibitem{moses2022museum}
    \DIFdel{Moses L, Pachter L. Museum of spatial transcriptomics. Nat Methods. 2022;19(5):534-546.
}%DIFDELCMD <

%DIFDELCMD < 	\bibitem{saelens2019comparison}
    \DIFdel{Saelens W, Cannoodt R, Todorov H, Saeys Y. A comparison of single-cell trajectory inference methods. Nat Biotechnol.2019;37(5):547-554.
}%DIFDELCMD <

%DIFDELCMD < 	\bibitem{geleta2026autoencoders}
    \DIFdel{Geleta M, Montserrat DM, Giro-i-Nieto X, et al. Autoencoders for genomic variation analysis. Genome Res. 2026;36(2): 348-360.
}%DIFDELCMD <

%DIFDELCMD < 	\bibitem{efremova2021computational}
    \DIFdel{Efremova M, Vento-Tormo R. Computational challenges of integrating multi-modal single-cell data. Nat Biotechnol. 2021;39(4):416-418.
}%DIFDELCMD <

%DIFDELCMD < 	\bibitem{regev2017human}
    \DIFdel{Regev A, et al. The Human Cell Atlas. eLife. 2017;6:e27041.
}%DIFDELCMD <

%DIFDELCMD <     %%%
%DIF < %%Reviewer-suggested references (R1.4: lightweight CNNs; R2.7: cross-domain ML)%%%
\DIFdelend \DIFaddbegin \DIFadd{s13059-017-1382-0.
}\DIFaddend

\end{thebibliography}

\end{document}
